%%%%%%%%%%%%%%%%%%%%%%%%%%%%%%%%%%%%%%%%%%%%%%%%%%%%%%%%%%%%%%%%%%%%%%
% Lớp 11 | Đường thẳng, mặt phẳng. Quan hệ song song trong | Đường thẳng và mặt phẳng song song
% Dạng 1H4?3-2: Đường thẳng song song với mặt phẳng
% Số câu: 117  (sinh tự động bởi tach_cau.py; từ nay thầy sửa trực tiếp file này)
%%%%%%%%%%%%%%%%%%%%%%%%%%%%%%%%%%%%%%%%%%%%%%%%%%%%%%%%%%%%%%%%%%%%%%

%%%%%%%%%%%%%%%% TRẮC NGHIỆM 4 PHƯƠNG ÁN %%%%%%%%%%%%%%%%

% nguồn: Ok_Toan11_TracNghiem_Lan2
\begin{ex}%[1H4N3-2]%[ma:cfa1b60a]
	Cho hình chóp $S.ABCD$ có đáy $ABCD$ là hình bình hành, $M$ và $N$ là hai điểm trên $SA,SB$ sao cho $\dfrac{SM}{SA}=\dfrac{SN}{SB}=\dfrac{1}{3} $. Vị trí tương đối giữa $MN$  và $( ABCD)$ là
	 \choice
		{$MN$  nằm trên $( ABCD) $}
		 {$MN$ cắt $( ABCD) $}
		{\True $MN$ song song $( ABCD)  $}
		{$MN$ và $( ABCD)$ chéo  nhau}
	\loigiai{ 
		Theo định lí Talet, ta có $\dfrac{SM}{SA}=\dfrac{SN}{SB}$ suy ra $MN$ song song với $AB$.\\ 
		Mà $AB$ nằm trong mặt phẳng $( ABCD)$ suy ra $MN  \parallel ( ABCD) $.
	} 
\end{ex}

% nguồn: Ok_Toan11_TracNghiem_Lan2
\begin{ex}%[1H4N3-2]%[ma:e6b212e6]
	Cho  hình chóp $S.ABCD$ có đáy $ABCD$ là hình bình hành. $M$, $N$, $K$ lần lượt là trung điểm của $DC$, $BC$, $SA$. Gọi $H$ là giao điểm của $AC$ và $MN$. Trong các khẳng định sau, khẳng định nào \textbf{sai}?
	 \choice
		{$MN$ chéo $SC$} 
		{$MN\parallel \left(SBD\right)$} 
		{\True  $MN\parallel \left(ABCD\right)$}
		{$MN\cap \left(SAC\right)=H$} 
	\loigiai{ 
		\begin{center}
			\begin{tikzpicture}[declare function={a=2; b=4; h=4;},line join=round]
					\path (0,0) coordinate  (A)
						(-145:a)  coordinate (B)
						(b,0) coordinate  (D)
						 ($(B)!0.5!(D)$) coordinate (O)
						 ($(O)+(0,h)$) coordinate (S);
					 \path ($(D)-(A)+(B)$) coordinate (C);
					\path  ($(D)!0.5!(C)$) coordinate (M)
						($(B)!0.5!(C)$) coordinate (N) 
						($(S)!0.5!(A)$)  coordinate (K);
					\path (intersection of A--C and M--N) coordinate (H); 
					\draw[dashed] (S)--(A)--(B) 
						(A)--(D)--(B) 
						(A)--(C) 
						 (M)--(N);
					 \draw (B)--(C)--(D)
						 (B)--(S)
						 (D)--(S)--(C);
					 \foreach \t/\g in {A/150,B/-90,C/-90,D/0,S/90,M/-30,N/-90,K/135,H/90}{
						\draw[fill=black] (\t) circle (1pt) node[shift={(\g:7pt)},font=\scriptsize]{$  \t $};
					 }
				\end{tikzpicture}
		\end{center}
		 Vì $MN\subset \left(ABCD\right)$ nên $MN$ không song song với mặt phẳng $\left(ABCD\right)\Rightarrow MN\parallel \left(ABCD\right)$ sai.
	} 
\end{ex}

% nguồn: Ok_Toan11_TracNghiem_Lan2
\begin{ex}%[1H4N3-2]%[ma:44e5b07b]
	Cho hình chóp $S.ABCD$ có đáy $ABCD$ là hình thang, đáy lớn $AB$. Gọi $P$, $Q$ lần lượt là hai điểm nằm trên cạnh $SA$ và $SB$ sao cho $\dfrac{SP}{SA}= \dfrac{SQ}{SB} =\dfrac{1}{3}$. Khẳng định nào sau đây là  đúng?
	\choice 
		 {$PQ$ cắt $\left(ABCD\right)$}
		{$PQ\subset \left(ABCD\right)$} 
		 {\True $PQ\parallel \left(ABCD\right)$}
		 {$PQ$ và $CD$ chéo nhau}
	\loigiai{ 
		\begin{center}
			\begin{tikzpicture}
					 \tikzset{declare function={a=5; b=3; h=3.5;}}
					\path  (0,0) coordinate (A)
						(a+1,0)  coordinate (B)
						 (-45:b) coordinate (D)
						($(D)+({0.45*a},0)$) coordinate  (C)
						(0,h) coordinate  (S);
					\path  ($(S)!1/3!(A)$) coordinate (P);
					\path ($(S)!1/3!(B)$) coordinate  (Q);
					 \draw[dashed] (A)--(B)
						 (P)--(Q);
					\draw (S)--(D)--(C)--(B)--cycle (D)--(A)--(S)--(C); 
					\foreach  \t/\g in {A/180,D/-90,C/-90,B/0,S/120,P/180,Q/45}{
						\draw[fill=white] (\t) circle (1pt) node[shift={(\g:9pt)},font=\scriptsize]{$\t$}; 
					} 
				\end{tikzpicture}
		\end{center}
		Ta có  $\heva{&PQ\parallel AB \\ &AB\subset \left(ABCD\right) \\& PQ\subset \left(ABCD\right)}\Rightarrow PQ\parallel \left(ABCD\right)$.
	} 
\end{ex}

% nguồn: Ok_Toan11_TracNghiem_Lan2
\begin{ex}%[1H4N3-2]%[ma:518ca77d]
	Cho hình chóp $S.ABCD$ có đáy là hình bình hành. Gọi $A'$, $B'$ lần lượt  là trung điểm của $SA$, $SB$. Đường thẳng $A'B'$ song song với mặt phẳng nào sau đây?
	\choice 
		{$(  SAB)$}
		{$( SBC)$} 
		{\True $(  SCD)$}
		{$(  SAD)$}
	 \loigiai{
		\begin{center}
			\begin{tikzpicture}[scale=1,font=\footnotesize,line join=round,line cap=round]
					\foreach \x\y\t in  {0/0/A,-1.7/-1.6/B,2.5/-1.6/C,0.4/2.8/S}
					\coordinate (\t) at (\x,\y); 
					\coordinate (D) at ($(A)+(C)-(B)$); 
					\coordinate (A')  at ($(S)!1/2!(A)$);
					\coordinate  (B') at ($(S)!1/2!(B)$);
					\draw (S)--(B)--(C)--(S)--(D)--(C); 
					 \draw[dashed](B)--(A)--(D)
						(S)--(A) 
						(A')--(B'); 
					 \foreach \t/\g in {S/90,A/160,B/-140,C/-45,D/0,A'/-20,B'/120}
					\draw[fill=black] (\t)  circle (1.0pt)
						 node[shift={(\g:7pt)}]{$\t$};
				\end{tikzpicture}
		\end{center}
		Vì $A'B'$ song song với $AB$ và $AB$ song song với $CD$ nên $A'B'$ song song với $CD$.\\ 
		Hơn nữa, $A'B'$ không chứa trong $( SCD)$ nên $A'B'$ song  song với $( SCD) $.
	 }
\end{ex}

% nguồn: Ok_Toan11_TracNghiem_Lan2
\begin{ex}%[1H4N3-2]%[ma:f39111f5]
	Cho hình chóp $S.ABCD$ có đáy $ABCD$ là hình bình hành. Đường thẳng $AD$ song  song với mặt phẳng nào trong các mặt phẳng dưới đây?
	 \choice
		 {\True $( SBC)$}
		{$( ABCD)$} 
		 {$SAC)$}
		{$(  SAB)$}
	 \loigiai{
		Do $AD \parallel BC$, $AD\not \subset (  SBC)$, $BC \subset ( SBC)$ nên $AD \parallel ( SBC)$.
	} 
\end{ex}

% nguồn: Ok_Toan11_TracNghiem_Lan2
\begin{ex}%[1H4N3-2]%[ma:e6ff34b8]
	Cho hình chóp tứ  giác $S.ABCD$. Gọi $M$ và $N$ lần lượt là trung điểm của $SA$ và $SC.$ Khẳng định nào sau đây đúng?
	 \choice
		{\True $MN \parallel ( ABCD) $} 
		 {$MN \parallel ( SAB) $}
		 {$MN \parallel ( SCD) $}
		{$MN \parallel ( SBC)  $}
	 \loigiai{
		Xét tam giác $SAC$ có $M$, $N$ lần lượt là trung  điểm của $SA$, $SC$.\\
		Suy ra  $MN \parallel AC$ mà $AC \subset ( ABCD) \Rightarrow MN \parallel ( ABCD) $.
	 }
\end{ex}

% nguồn: Ok_Toan11_TracNghiem_Lan2
\begin{ex}%[1H4N3-2]%[ma:76b23c2f]
	Cho tứ diện $ABCD$. Gọi $M$,  $N$, $P$ lần lượt là trung điểm của $AB$, $AC$, $AD$. Đường thẳng $MN$ song song với mặt phẳng nào trong các mặt phẳng sau đây?
	 \choice
		{$(PCD)$} 
		 {$(ABC)$}
		{\True  $(BCD)$}
		{$(BCD)$} 
	 \loigiai{
		\begin{center}
			\begin{tikzpicture}[scale=1,font=\footnotesize,line join=round,line cap=round,>=stealth]
					\coordinate (B) at (0,0); \coordinate (C) at (1.5,-2); \coordinate (D) at (6,0); 
					 \coordinate (A) at ($(B)+(2,4)$);
					 \coordinate (M) at ($(A)!1/2!(B)$); \coordinate (N) at ($(C)!1/2!(A)$); \coordinate (P) at ($(A)!1/2!(D)$);
					 \draw (A)--(B)--(C)--(D)--(A)--(C)
						 (M)--(N)--(P);
					\draw[dashed] (B)--(D) 
						(M)--(P); 
					\foreach  \x/\g in {A/90,B/180,C/-90,D/0,M/130,N/-30,P/30} \fill (\x) circle (1pt)
						 ($(\x)+(\g:3mm)$) node{$\x$};
				\end{tikzpicture}
		\end{center}
		Theo đề bài, ta có $MN$ là đường trung bình của tam giác $ABC$ nên $MN$ song song với $BC$, $MN$ không nằm trong mặt phẳng $(BCD)$ nên đường thẳng $MN$  song song với mặt phẳng $(BCD)$.
	} 
\end{ex}

% nguồn: Ok_TracNghiemLop11_Lan1
\begin{ex}%[1H4N3-2]%[ma:946804ee]
	Cho hình chóp $S.ABCD$ có đáy là hình  bình hành. Gọi $E$, $F$ lần lượt thuộc $SA$, $SB$ sao cho $SA = 2SE$, $SB = 2SF$. Khẳng định nào sau đây {\bf sai}?
	 \choice
		{$EF  \parallel (ABCD)$}
		{$EF \parallel  (SCD)$}
		{\True  $EF\parallel (SBD)$}
		 {$EF\subset (SAB)$}
	 \loigiai{
		\begin{center}
			\begin{tikzpicture}[scale=0.8,line join = round,line cap = round]
					 \def\a{4}
					\path  (0:0) coordinate (A)
						 ++(0:\a) coordinate (D)
						 ++(-130:\a/2) coordinate (C)
						 ($(A)+(C)-(D)$) coordinate (B)
						($(A)+(80:\a)$) coordinate  (S)
						($(S)!.5!(A)$) coordinate (E) 
						($(S)!.5!(B)$) coordinate  (F);
					 \draw[dashed] (B)--(A)--(D)
						(A)--(S) 
						(E)--(F); 
					\draw (B)-- (C)--(D) 
						 (B)--(S)
						 (C)--(S)
						(D)--(S); 
					 \foreach \x/\g in {A/135,B/-135,C/-45,D/45,S/90,E/0,F/180}
					\fill[black] (\x)  circle (1pt)
						($(\g:4mm)+(\x)$)  node {$\x$};
				\end{tikzpicture}
		\end{center}
		Ta có $E\in SA$,  $F\in SB$, suy ra $EF\subset (SAB)$.\\
		 Lại có $\dfrac{SE}{SA}=\dfrac{SF}{SB}=\dfrac{1}{2}$ nên $EF \parallel AB\parallel CD$.\\
		Mà $AB \subset (ABCD)$, $CD  \subset (SCD)$ và $EF\not\subset(ABCD)$, $EF\not\subset(SCD)$.\\
		Từ đó  suy ra $EF \parallel (ABCD)$ và $EF \parallel (SCD)$.\\
		 Vậy $EF \parallel (SBD)$ là khẳng định sai.
	} 
\end{ex}

% nguồn: Ok_TracNghiemLop11_Lan1
\begin{ex}%[1H4N3-2]%[ma:a1ee72dd]
	Cho hình chóp $S.ABCD$ với đáy $ABCD$ là hình bình hành. Gọi $I$, $J$, $K$, $L$ lần lượt là trung điểm của $SA$, $SB$, $SC$, $SD$. Mệnh đề nào sau đây là đúng? 
	\choice 
		{$AC\parallel (SIK)$} 
		 {$AB\parallel (SIK)$}
		 {\True $CD\parallel (IJK)$}
		{$CD\parallel (SIK)$} 
	 \loigiai{
		\begin{center}
			\begin{tikzpicture}[scale=1,font=\footnotesize,line join=round,line cap=round,>=stealth]
					\coordinate (A)  at (0,0);
					\coordinate  (B) at (-1.3,-1.7);
					\coordinate (D) at  (4,0);
					\coordinate (C) at ($(B)+(D)-(A)$); 
					\coordinate  (S) at (1.3,4);
					\coordinate (I) at  ($(S)!1/2!(A)$);
					\coordinate (J) at ($(S)!1/2!(B)$); 
					\coordinate (K) at ($(S)!1/2!(C)$); 
					 \coordinate (L) at ($(S)!1/2!(D)$);
					\draw (S)--(B)--(C)--(D)--(S)--(C); 
					 \draw[dashed] (S)--(A)--(B)
						(D)--(A); 
					 \draw (J)--(K)--(L);
					\draw[dashed]  (J)--(I)--(L);
					\foreach \x/\g in{S/90,A/135,B/-135,C/-45,D/0,I/-45,J/180,K/0,L/45} \fill (\x)  circle (1pt)+(\g:.35) node{$\x$};
				\end{tikzpicture}
		\end{center}
		Vì $I$, $J$, $K$, $L$ lần lượt là trung điểm của $SA$, $SB$, $SC$, $SD$ nên ta có $KL$ là đường trung bình của tam  giác $SCD$.\\
		Từ đó ta có $CD\parallel  KL$ mà $KL\subset (IJKL)\equiv(IJK)$ suy ra $CD\parallel (IJK)$.
	} 
\end{ex}

% nguồn: Ok_TracNghiemLop11_Lan1
\begin{ex}%[1H4N3-2]%[ma:65503526]
	 Cho hình chóp $S.ABCD$ có đáy $ABCD$ là hình bình hành. Đường thẳng $AD$ song song với mặt phẳng nào trong các mặt phẳng dưới đây?
	\choice 
		{\True  $(SBC)$}
		{$(ABCD)$} 
		 {$(SAC)$}
		{$(SAB)$} 
	 \loigiai{\immini{
		Do $AD \parallel BC$, $AD$ không nằm trong $(SBC)$  và $BC$ lại nằm trong $(SBC)$ nên $AD \parallel (SBC)$.}
		 {
		\begin{tikzpicture}[line join=round,line cap=round,>=stealth,scale=0.6,font=\footnotesize]
				 \tikzset{every node/.style={outer sep=-0.5mm}}
				\foreach  \x/\y/\n in {0/0/A,5/0/D,-2/-2/B,3/-2/C,0/5/S} \coordinate (\n) at (\x,\y);
				 \draw (S)--(B)--(C)--(D)--(S)
					(C)--(S)--(B) 
					(D)--(S); 
				\draw[dashed] (D)--(A)--(S) 
					 (A)--(B);
				\foreach \t/\g in  {A/170,B/180,C/0,D/0,S/90}{
					 \draw[fill] (\t) circle (1pt) node[shift={(\g:10pt)}]{$ \t $};}
				\draw pic[thin,draw,angle  radius=2mm] {right angle= S--A--D};
			\end{tikzpicture}
		} 
	} 
\end{ex}

% nguồn: Ok_TracNghiemLop11_Lan1
\begin{ex}%[1H4N3-2]%[ma:63a5d3f2]
	Cho  hình chóp $S.ABCD$ có đáy $ABCD$ là hình bình hành. Chọn khẳng định đúng.
	\begin{center}
		\begin{tikzpicture}[scale=1,font=\footnotesize,line join=round,line cap=round,>=stealth]
				 \path
					(0,0) coordinate (A) 
					(-1,-1) coordinate  (B)
					 ($(A)+(4,0)$) coordinate (D)
					 ($(B)+(4,0)$) coordinate (C)
					($(A)+(0.5,3)$) coordinate  (S)
					(intersection of  A--C and B--D) coordinate (O);
				\draw[dashed]  (S)--(A)--(B)
					 (C)--(A)--(D)
					 (B)--(D);
				\draw (S)--(B)--(C)--(D)--(S)--(C); 
				\foreach \p/\g in  {S/135,A/135,B/-135,C/-45,D/45,O/-45}
				\fill[black](\p) circle  (1pt)
					 ($(\p)+(\g:3mm)$) node{$\p$};
			\end{tikzpicture}
	\end{center}
	\choice 
		 {$AB \parallel (SBD)$}
		{$BC\parallel (SCD)$} 
		{\True $AD  \parallel (SBC)$}
		{$BD \parallel (SAC)$} 
	\loigiai{Vì $ABCD$ là  hình bình hành nên $AD \parallel BC$, $AD\not\subset (SBC)$ nên $AD \parallel (SBC)$.}
\end{ex}

% nguồn: Ok_Toan11_TracNghiem_Lan2
\begin{ex}%[1H4H3-2]%[ma:e3664275]
	Cho tứ diện $ABC$. Gọi $G$ là trọng tâm tam giác $ABD$,  $M$ là điểm trên cạnh $BC$ sao cho $MB=2MC$. Khi đó đường thẳng $MG$ song song với mặt phẳng nào dưới đây?
	 \choice
		{\True  $\left(ACD\right)$}
		 {$\left(BCD\right)$}
		 {$\left(ABD\right)$}
		 {$\left(ABC\right)$}
	\loigiai{ 
		\begin{center}
			\begin{tikzpicture}[declare function={a=2; b=4; h=3;},line join=round]
					 \path (0,0) coordinate (B)
						 (0:b) coordinate (D)
						 (-45:a) coordinate (C)
						(50:h)  coordinate (A);
					\path ($(A)!0.5!(D)$) coordinate (E); 
					\path ($(B)!2/3!(C)$)  coordinate (M);
					\path ($(B)!2/3!(E)$)  coordinate (G);
					\draw (B)--(C)--(D)--(A)--cycle (A)--(C)--(E); 
					\draw[dashed] (E)--(B)--(D) 
						(M)--(G); 
					\foreach \t/\g in {B/-180,C/-120,D/0,A/90,E/60,M/-110,G/120}{\draw[fill=black] (\t) circle (1pt) node[shift={(\g:7pt)},font=\scriptsize]{$  \t $};}
				\end{tikzpicture}
		\end{center}
		Gọi $E$ là  trung điểm $AD$. \\
		Suy ra $\dfrac{BG}{BE}=\dfrac{BM}{BC}=\dfrac{1}{3}  \Rightarrow MG \parallel CE$.\\
		Mặt khác $CE \subset (ACD)$. Vậy  $MG \parallel (ACD)$.
	 }
\end{ex}

% nguồn: Ok_Toan11_TracNghiem_Lan2
\begin{ex}%[1H4H3-2]%[ma:4b72ee0b]
	Cho tứ diện  $ABCD$, gọi $G_1$, $G_2$ lần lượt là trọng tâm tam giác $BCD$ và $ACD$. Mệnh đề nào sau đây \textbf{sai}?
	 \choice
		 {$G_1 G_2 \parallel \left(ABD\right)$}
		{Ba đường  thẳng $BG_1$, $AG_2$ và $CD$ đồng quy}
		{$G_1  G_2 \parallel \left(ABC\right)$}
		{\True $G_1 G_2=\dfrac{2}{3}  AB$}
	 \loigiai{
		\begin{center}
			\begin{tikzpicture}[declare function={a=2; b=4; h=4;},line join=round]
					 \path (0,0) coordinate (B)
						 (0:b) coordinate (D)
						 (-45:a) coordinate (C)
						(60:h) coordinate (A); 
					\path  ($(D)!0.5!(C)$) coordinate (M);
					\path ($(M)!1/3!(B)$)  coordinate (G_1);
					 \path ($(M)!1/3!(A)$) coordinate (G_2);
					 \draw (B)--(C)--(D)--(A)--cycle (M)--(A)--(C);
					\draw[dashed] (M)--(B)--(D) 
						(G_1)--(G_2); 
					\foreach \t/\g in {B/-180,C/-120,D/0,A/90,M/-45,G_1/-90,G_2/30}{\draw[fill=white] (\t) circle (1pt) node[shift={(\g:7pt)},font=\scriptsize]{$ \t  $};}
				\end{tikzpicture}
		\end{center}
		Gọi $M$ là trung điểm  của $CD$.\\
		 Xét $\triangle ABM$ ta có: $\dfrac{MG_1}{MB}=\dfrac{MG_2}{MA}=\dfrac{1}{3} \Rightarrow \heva{& G_1 G_2 \parallel AB \\ &G_1 G_2=\dfrac{1}{3} AB} \Rightarrow G_1 G_2=\dfrac{2}{3} AB$ sai.\\
		Vì $G_1  G_2 \parallel AB\Rightarrow G_1 G_2 \parallel \left(ABD\right)$ $\Rightarrow G_1 G_2 \parallel \left(ABD\right)$ đúng.\\
		Vì $G_1 G_2 \parallel AB\Rightarrow G_1 G_2  \parallel \left(ABC\right)$ $\Rightarrow G_1 G_2 \parallel \left(ABC\right)$ đúng.\\
		Ba đường $BG_1$, $AG_2$, $CD$, đồng quy tại $M$. 
	 }
\end{ex}

% nguồn: Ok_Toan11_TracNghiem_Lan2
\begin{ex}%[1H4H3-2]%[ma:a70615ac]
	Cho hình chóp $S.ABCD$ có đáy là hình thang, $AB\parallel CD$ và $AB=2CD$. Lấy $E$ thuộc cạnh $SA$, $F$ thuộc cạnh $SC$ sao cho $\dfrac{SE}{SA}=\dfrac{SF}{SC}  =\dfrac{2}{3}$. Khẳng định nào dưới đây đúng?
	 \choice
		{Đường thẳng $EF$ song song với mặt  phẳng $\left(SAC\right)$}
		{Đường thẳng $EF$ cắt  đường thẳng $AC$}
		{\True Đường thẳng $AC$ song song với  mặt phẳng $\left(BEF\right)$}
		 {Đường thẳng $CD$ song song với mặt phẳng $\left(BEF\right)$}
	\loigiai{ 
		\begin{center}
			\begin{tikzpicture}
					\tikzset{declare function={a=5; b=2; h=4;}} 
					\path  (0,0) coordinate (A)
						 (a+1,0) coordinate (B)
						 (-45:b) coordinate (D)
						 ($(D)+({0.45*a},0)$) coordinate (C)
						(70:h) coordinate  (S);
					\path ($(S)!2/3!(A)$)  coordinate (E);
					\path ($(S)!2/3!(C)$) coordinate  (F);
					\draw[dashed]  (C)--(A)--(B)
						 (F)--(E)--(B);
					\draw (S)--(D)--(C)--(B)--cycle (D)--(A)--(S)--(C) 
						(B)--(F); 
					\foreach \t/\g in {A/180,D/-90,C/-90,B/0,S/120,E/180,F/205}{ 
						\draw[fill=white] (\t) circle  (1pt) node[shift={(\g:9pt)},font=\scriptsize]{$\t$};
					 }
				\end{tikzpicture}
		\end{center}
		Vì $\dfrac{SE}{SA}=\dfrac{SF}{SC}=\dfrac{2}{3}$ nên đường thẳng  $EF \parallel AC$. \\
		Mà $EF\subset \left(BEF\right)$, $AC\not\subset \left(BEF\right)$ nên $AC$ song song với mặt phẳng  $\left(BEF\right)$.
	} 
\end{ex}

% nguồn: Ok_Toan11_TracNghiem_Lan2
\begin{ex}%[1H4H3-2]%[ma:e2282a5e]
	Cho hình chóp $S.ABCD$ có đáy $ABCD$ là hình bình hành tâm  $O$, $M$ là trung điểm $SA$. Khẳng định nào sau đây là đúng?
	 \choice
		{\True  $OM\parallel \left(SCD\right)$}
		{$OM\parallel  \left(SBD\right)$}
		{$OM\parallel \left(SAB\right)$} 
		 {$OM\parallel \left(SAD\right)$}
	 \loigiai{
	\immini{Ta có $M$ là trung  điểm $SA$; $O$ là trung điểm $AC$ $\Rightarrow OM$ là đường trung bình $\triangle SAC$.
		\[\Rightarrow OM\parallel SC \quad\left(SC\subset \left(SCD\right); OM\not\subset \left(SCD\right)\right)\Rightarrow  OM\parallel \left(SCD\right).\]
		 }
		{\begin{tikzpicture}[scale=1,font=\footnotesize,line join=round,line cap=round,>=stealth,declare function={gocx=80; goc=-150; a=4; b=a/1.5; h=3;}]
		\path (0,0) coordinate (A)--+(gocx:h) coordinate (S) 
			(a,0)  coordinate (B)
			(goc:b)  coordinate (D)
			+(a,0)  coordinate (C);
		 \path ($(C)!0.5!(A)$) coordinate (O);
		\path ($(S)!0.5!(A)$)  coordinate (M);
		\draw[dashed]  (S)--(A)
			 (D)--(A)--(B)
			(O)--(M) 
			 (A)--(C)
			 (B)--(D);
		 \draw (S)--(D)--(C)--(S)--(B)--(C)--cycle;
		\foreach \x/\goc in {S/90,A/150,B/0,C/-60,D/180,O/-90,M/30} 
		 \draw[fill=white] (\x) node[shift={(\goc:7pt)}]{$\x$} circle (1pt);
	\end{tikzpicture}
		 }
	 }
\end{ex}

% nguồn: Ok_Toan11_TracNghiem_Lan2
\begin{ex}%[1H4H3-2]%[ma:8d982861]
	Cho hình chóp $S.ABCD$, gọi $M$ và $N$ lần lượt là trung điểm của $AB$ và $AD$.  Khẳng định nào sau đây đúng?
	\choice 
		{$MN\parallel  (SBC)$}
		 {\True $MN\parallel (SBD)$}
		{$MN\parallel (ABCD)$} 
		{$MN\parallel  (SCD)$}
	\loigiai{ 
	\immini{Vì $MN$ là đường trung bình của tam giác $ABD$  nên $MN \parallel BD$.\\
		Khi đó ta có $\heva{&MN\not\subset(SBD)\\& BD\subset(SBD)\\&MN \parallel BD}\Rightarrow MN \parallel (SBD)$.}{\begin{tikzpicture}[scale=0.8,font=\footnotesize,line join=round,line cap=round,>=stealth]
		\coordinate (A) at (0,0); 
		\coordinate (B) at  (2,-2);
		\coordinate (C) at (5,0); 
		\coordinate (D) at  (3,1);
		\coordinate  (S) at ($(A)+(60:4)$);
		\coordinate (M)  at ($(A)!.5!(B)$);
		\coordinate  (N) at ($(A)!.5!(D)$);
		\draw  (S)--(A)--(B)--(C)--(S)--(B);
		 \draw [dashed] (A)--(D)--(C)
			(S)--(D)--(B) 
			(M)--(N); 
		\foreach  \x/\g in {A/180,B/-90,C/0,D/60,M/-120,N/90,S/90} \fill (\x) circle (1pt) +(\g:3mm) node{$\x$};
	\end{tikzpicture}
		 }
	} 
\end{ex}

% nguồn: Ok_Toan11_TracNghiem_Lan2
\begin{ex}%[1H4H3-2]%[ma:6ef86e1c]
	Cho hình chóp $S.ABCD$ có đáy là hình bình hành. Gọi $G$ là trọng tâm của tam giác $SAB$ và $I$ là trung điểm của $AB$. Lấy điểm $M$ trên đoạn $AD$ sao cho $AD=3AM$. Đường thẳng qua $M$ và song song với  $AB$ cắt $CI$ tại $J$, đường thẳng $JG$ không song song với mặt phẳng nào sau đây?
	 \choice
		{\True $(SCD)$} 
		{$(SAD)$} 
		 {$(SBC)$}
		{$(SAC)$} 
	\loigiai{ 
		\begin{center}
			\begin{tikzpicture}[scale=1,font=\footnotesize,line join=round,line cap=round,>=stealth]
					\coordinate (A) at  (0,0);
					\coordinate (B) at (-2,-2); 
					\coordinate (D)  at (5,0);
					\coordinate (C)  at ($(B)+(D)-(A)$);
					\coordinate (S) at ($(A)+(0.5,4)$); 
					\coordinate (I) at ($(A)!0.5!(B)$);  \coordinate (G) at ($(S)!2/3!(I)$);
					 \coordinate (M) at ($(A)!1/3!(D)$); \coordinate (J) at ($(I)!1/3!(C)$);
					 \draw (S)--(B)
						(S)--(C)--(D)--(S) 
						(B)--(C); 
					\draw[dashed,thin](A)--(B) 
						(C)--(D) 
						 (S)--(A)--(D)
						(S)--(I)--(C) 
						 (M)--(J)--(G);
					\foreach \i/\g in {S/90,A/60,B/-90,C/-90,D/0,I/-90,J/-90,M/90,G/0}{\draw[fill](\i)  circle (1pt)
						($(\i)+(\g:3mm)$) node[scale=1]{$\i$};} 
				\end{tikzpicture}
		\end{center}
		Ta có \[AI\parallel CD\Rightarrow  \dfrac{AM}{AD}=\dfrac{JI}{IC}=\dfrac{1}{3}.\]
		Mà $\dfrac{GI}{SI}=\dfrac{1}{3}$ nên $\dfrac{GI}{SI}=\dfrac{JI}{IC}=\dfrac{1}{3} \Rightarrow GJ\parallel  SC$.\\
		Hơn  nữa, $SC$ nằm trong các mặt phẳng $(SCD)$, $(SBC)$, $(SAC)$ nên $JG\parallel (SCD)$, $IG\parallel (SBC)$, $JG\parallel (SAC)$.
	 }
\end{ex}

% nguồn: Ok_Toan11_TracNghiem_Lan2
\begin{ex}%[1H4H3-2]%[ma:aa1dd00b]
	Cho hai  hình bình hành $ABCD$ và $ABEF$ không cùng nằm trong một mặt phẳng. Gọi $O_1$, $O_2$ lần lượt là tâm của $ABCD$, $ABEF$. $M$ là trung điểm của $CD$. Chọn khẳng định {\bf sai} trong các khẳng định sau
	 \choice
		{\True $MO_2$ cắt  mặt phẳng $(BEC)$}
		 {$O_1O_2$ song song với mặt phẳng $(BEC)$}
		{$O_1O_2$ song song  với mặt phẳng $(EFM)$}
		{$O_1O_2$ song song  với mặt phẳng $(AFD)$}
	 \loigiai{
	 \immini{Gọi $J$ là giao điểm $AM$ và $BC$.\\
		Ta có $MO_1\parallel AD\parallel BC\Rightarrow MO_1\parallel  CJ$.\\
		Mà $O_1$ là trung điểm của  $AC$ nên $M$ là trung điểm của $AJ$.\\
		Do đó $MO_2\parallel  EJ$.\\
		Từ đó  suy ra $MO_2\parallel (BEC)$ (vì dễ nhận thấy $MO_2$ không nằm trên mặt phẳng $(BEC)$).\\
		Vậy  $MO_2$ không cắt mặt phẳng $(BEC)$.}
		{\begin{tikzpicture}[scale=1,font=\footnotesize,line join=round,line cap=round,>=stealth]
		 \coordinate (A) at (0,0);
		\coordinate (B) at (5,0);  \coordinate (E) at (6.5,1.2);
		\coordinate (D) at (0.3,3); 
		 \coordinate (C) at ($(B)+(D)-(A)$);
		\coordinate (F) at  ($(A)+(E)-(B)$);
		\coordinate (O1) at  ($(A)!0.5!(C)$);
		 \coordinate (O2) at ($(A)!0.5!(E)$);
		\coordinate (M) at  ($(C)!0.5!(D)$);
		\coordinate (J)  at (intersection of A--M and B--C);
		\draw (A)--(D)--(C) 
			 (B)--(A)
			 (B)--(E)
			 (B)--(D)
			 (A)--(C)--(E)--(J)--(A)
			(B)--(J); 
		 \draw[dashed](A)--(F)--(E)--(A)
			 (B)--(F)
			 (O2)--(M);
		\draw[fill] (O1) node[above]{$O_1$} circle (1pt); 
		\draw[fill] (O2) node[below]{$O_2$} circle  (1pt);
		 \foreach \i/\g in {A/-90,B/-90,C/0,D/180,M/90,J/90,E/0,F/30}{\draw[fill](\i) circle (1pt)
			($(\i)+(\g:3mm)$) node[scale=1]{$\i$};} 
	\end{tikzpicture}
		 }
	 }
\end{ex}

% nguồn: Ok_Toan11_TracNghiem_Lan2
\begin{ex}%[1H4H3-2]%[ma:5281a98a]
	Cho tứ diện $ABCD$, $G$ là trọng tâm  $\triangle ABD$ và $M$ là điểm trên cạnh $BC$ sao cho $BM=2MC$. Đường thẳng $MG$ song song với mặt phẳng nào sau đây?
	 \choice
		{\True  $(ACD)$}
		 {$(ABC)$}
		 {$(ABD)$}
		 {$(BCD)$}
	\loigiai{ 
	 \immini{Gọi $P$ là trung điểm của $AD$.\\
		 Ta có $\dfrac{BM}{BC}=\dfrac{BG}{BP}=\dfrac{2}{3} \Rightarrow MG \parallel CP$.\\
		Khi đó ta có $\heva{&MG\not\subset(ACD)\\& CP\subset(ACD)\\&MG \parallel CP}\Rightarrow MG \parallel (ACD)$.}{\begin{tikzpicture}[scale=1,font=\footnotesize,line join=round,line cap=round,>=stealth]
		\coordinate (B) at  (0,0);
		\coordinate (D)  at (2,-2);
		\coordinate  (C) at (5,0);
		\coordinate (A) at ($(B)+(60:4)$); 
		\coordinate  (M) at ($(B)!2/3!(C)$);
		\coordinate  (N) at ($(B)!.5!(D)$);
		\coordinate (P)  at ($(A)!.5!(D)$);
		\coordinate (G) at  ($(A)!2/3!(N)$);
		\draw  (A)--(B)--(D)--(C)--(A)--(D)
			(A)--(N) 
			 (B)--(P)--(C);
		\draw [dashed]  (B)--(C)
			(M)--(G); 
		\foreach \x/\g in {A/90,B/180,C/0,D/-90,M/-70,P/50,G/130} \fill (\x) circle  (1pt) +(\g:3mm) node{$\x$};
	\end{tikzpicture}
		 }
	} 
\end{ex}

% nguồn: Ok_Toan11_TracNghiem_Lan2
\begin{ex}%[1H4H3-2]%[ma:1b2706a5]
	Cho hình chóp $S.ABC$. Gọi  $M$, $N$ lần lượt là trung điểm của các cạnh $SB$, $SC$. Trong các khẳng định sau, khẳng định nào đúng?
	\choice 
		 {\True $MN\parallel (ABC)$}
		{$MN\parallel (SAB)$} 
		{$MN\parallel  (SAC)$}
		{$MN\parallel (SBC)$} 
	 \loigiai{
	\immini{Theo giả thiết thì $M$, $N$ lần lượt là trung điểm của $SB$,  $SC$ nên $MN$ là đường trung bình của $\triangle SBC$, do đó $MN\parallel BC$.\\
		Khi đó ta có $\heva{&MN\not\subset(ABC)\\& BC\subset(ABC)\\&MN \parallel BC}\Rightarrow MN \parallel (ABC)$.}{\begin{tikzpicture}[scale=0.8,font=\footnotesize,line join=round,line cap=round,>=stealth]
		\coordinate  (A) at (0,0);
		\coordinate  (B) at (3,-2);
		\coordinate (C) at (5,0); 
		\coordinate (S)  at ($(A)+(60:4)$);
		\coordinate (M) at  ($(S)!.5!(B)$);
		\coordinate  (N) at ($(S)!.5!(C)$);
		 \draw (S)--(A)--(B)--(C)--(S)--(B)
			 (M)--(N);
		\draw [dashed]  (A)--(C);
		\foreach \x/\g in {A/180,B/-90,C/0,S/90,M/200,N/30} \fill (\x) circle (1pt) +(\g:3mm)  node{$\x$};
	\end{tikzpicture}
		} 
	 }
\end{ex}

% nguồn: Ok_Toan11_TracNghiem_Lan2
\begin{ex}%[1H4H3-2]%[ma:3573a2d5]
	Cho  tứ diện $ABCD$. Gọi $G$ là trọng tâm của tam giác $ABD$, $Q$ thuộc cạnh $AB$ sao cho $AQ=2QB$, $P$ là trung điểm của $AB$, $M$ là trung điểm của $BD$. Khẳng định nào sau đây đúng?
	 \choice
		{$MP \parallel  ( BCD) $}
		{\True $GQ \parallel ( BCD)  $}
		{$MP$ cắt $(  ACD) $}
		 {$Q$ thuộc mặt phẳng $( CDP) $}
	 \loigiai{
		\begin{center}
			\begin{tikzpicture}[scale=1,font=\footnotesize,line join=round,line cap=round]
					\foreach \x\y\t in  {0/0/B,1.3/-1.6/C,4.5/0/D,1/3.5/A}
					\coordinate (\t)  at (\x,\y);
					\coordinate (M) at ($(B)!1/2!(D)$); 
					 \coordinate (P) at ($(A)!1/2!(B)$);
					\coordinate (Q) at ($(A)!2/3!(B)$); 
					\coordinate (G)  at ($(A)!2/3!(M)$);
					 \draw (A)--(B)--(C)--(A)--(D)--(C);
					\draw[dashed](B)--(D)--(P) 
						 (A)--(M)
						 (Q)--(G);
					\foreach  \t/\g in {A/90,B/180,C/-90,D/0,M/-90,Q/150,P/150,G/60}
					\draw[fill=black] (\t) circle  (1.0pt)
						 node[shift={(\g:9pt)}]{$\t$};
				\end{tikzpicture}
		\end{center}
		Ta có  $G$ là trọng tâm tam giác $ABD$ nên $\dfrac{AG}{AM}=\dfrac{2}{3} $.\\
		Điểm $Q \in AB$ sao cho $AQ=2QB \Leftrightarrow  \dfrac{AQ}{AB}=\dfrac{2}{3} \Rightarrow \dfrac{AG}{AM}=\dfrac{AQ}{AB} \Rightarrow GQ \parallel BD$.\\
		Mặt khác $BD$ nằm  trong mặt phẳng $( BCD)$ suy ra $GQ \parallel ( BCD) $.
	} 
\end{ex}

% nguồn: Ok_Toan11_TracNghiem_Lan2
\begin{ex}%[1H4H3-2]%[ma:44d99b8c]
	Cho hình chóp $S.ABCD$ có đáy  là hình bình hành. Gọi $I$ và $J$ lần lượt là trung điểm của $SC$ và $BC$. Chọn khẳng định đúng trong các khẳng định sau:
	 \choice
		 {$JI \parallel (SAC$}
		{\True $JI \parallel  (SAB)$}
		{$JI  \parallel (SBC$}
		{$JI \parallel (SCD$} 
	 \loigiai{
		\begin{center}
			\begin{tikzpicture}[scale=1,font=\footnotesize,line join=round,line cap=round,>=stealth,declare function={gocx=80; goc=-130; a=4; b=a/2; h=3;}]
					\path  (0,0) coordinate (A)--+(gocx:h) coordinate (S)
						(a,0) coordinate  (D)
						 (goc:b) coordinate (B)
						+(a,0) coordinate (C); 
					 \path ($(S)!0.5!(C)$) coordinate (I);
					 \path ($(B)!0.5!(C)$) coordinate (J);
					\draw[dashed] (S)--(A) 
						 (D)--(A)--(B);
					\draw  (S)--(D)--(C)--(S)--(B)--(C)--cycle (I)--(J);
					\foreach  \x/\goc in {S/90,A/150,B/180,C/-60,D/0,I/30,J/-90}
					\draw[fill=white] (\x) node[shift={(\goc:7pt)}]{$\x$} circle  (1pt);
				\end{tikzpicture}
		\end{center}
		Ta  có $JI \parallel SB$, $SB\subset \left(SAB\right)$.\\
		Vậy $JI  \parallel (SAB)$.
	} 
\end{ex}

% nguồn: Ok_Toan11_TracNghiem_Lan2
\begin{ex}%[1H4H3-2]%[ma:a3aa19ae]
	Cho hình chóp $S.ABCD$ có đáy $ABCD$ là hình bình hành tâm $O$, $M$ là trung  điểm $SA$. Khẳng định nào sau đây là đúng?
	\choice 
		{\True $OM\parallel  \left(SCD\right)$}
		 {$OM\parallel \left(SBD\right)$}
		 {$OM\parallel \left(SAB\right)$}
		{$OM\parallel \left(SAD\right)$} 
	 \loigiai{
		\begin{center}
			\begin{tikzpicture}[declare function={a=2; b=4; h=4;},line join=round]
					 \path (0,0) coordinate (A)
						(-145:a) coordinate  (B)
						(b,0)  coordinate (D)
						($(B)!0.5!(D)$) coordinate (O) 
						($(O)+(0,h)$) coordinate (S) 
						 ($(S)!0.5!(A)$) coordinate (M);
					\path ($(D)-(A)+(B)$)  coordinate (C);
					\draw[dashed]  (S)--(A)--(B)
						 (A)--(D)--(B)
						 (A)--(C)
						 (M)--(O);
					\draw (B)--(C)--(D) 
						(B)--(S) 
						(D)--(S)--(C); 
					\foreach \t/\g in {A/150,B/-90,C/-90,D/0,S/90,O/-90,M/180}{ 
						\draw[fill=white] (\t)  circle (1pt) node[shift={(\g:7pt)},font=\scriptsize]{$ \t $};
					} 
				\end{tikzpicture}
		\end{center}
		Ta có $M$ là  trung điểm $SA$; $O$ là trung điểm $AC$ suy ra $OM$ là đường trung bình $SAC$.\\
		Suy ra  $\heva{& OM\parallel SC\left(SC\subset (SCD)\right)\\&OM\not\subset \left(SCD\right)} \Rightarrow OM\parallel \left(SCD\right)$.
	} 
\end{ex}

% nguồn: Ok_Toan11_TracNghiem_Lan2
\begin{ex}%[1H4H3-2]%[ma:54236230]
	Cho tứ diện $ABCD$. Gọi $G$ là trọng tâm tam giác $ABD$. $M$ là điểm trên cạnh $BC$ sao cho $MB=2MC$. Khi đó đường thẳng $MG$  song song với mặt phẳng nào dưới đây?
	 \choice
		 {\True $( ACD)$}
		 { $( BCD)$}
		 { $( ABD)$}
		{ $(  ABC)$}
	\loigiai{ 
		\begin{center}
			\begin{tikzpicture}[scale=1,font=\footnotesize,line join=round,line cap=round]
					\foreach  \x\y\t in {0/0/D,1.3/-1.6/B,4.5/0/C,1/3.5/A}
					\coordinate (\t) at  (\x,\y);
					\coordinate (E) at  ($(A)!1/2!(D)$);
					\coordinate  (M) at ($(C)!1/3!(B)$);
					\coordinate (G) at  ($(B)!2/3!(E)$);
					\coordinate (F) at ($(B)!1/2!(D)$); 
					\draw  (F)--(A)--(B)--(D)--(A)--(C)--(B)--(E);
					 \draw[dashed](E)--(C)--(D)
						(G)--(M); 
					\foreach \t/\g  in {A/90,C/0,B/-90,D/180,M/-80,G/180,E/150}
					\draw[fill=black] (\t) circle (1.0pt) 
						node[shift={(\g:7pt)}]{$\t$}; 
				\end{tikzpicture}
		\end{center}
		Gọi $E$ là  trung điểm $AD$. \\
		Xét tam giác $BCE$ có $\dfrac{BG}{BE}=\dfrac{BM}{BC}=\dfrac{2}{3}$ nên  suy ra $MG \parallel ( ACD)$.
	 }
\end{ex}

% nguồn: Ok_Toan11_TracNghiem_Lan2
\begin{ex}%[1H4H3-2]%[ma:10c946e2]
	Cho hình chóp $S.ABCD$ có đáy $ABCD$ là hình thoi tâm $O$. Gọi $I$ là trung điểm của $BC$, $K$ thuộc cạnh $SD$ sao cho $SK=\dfrac{1}{2} KD$, $M$ là giao điểm của của $BD$ và $AI$.  Khẳng định nào sau đây là đúng
	 \choice
		 {$MK\parallel \left(SCD\right)$}
		 {$MK\parallel \left(SBD\right)$}
		{$MK\parallel \left(ABCD\right)$} 
		{\True $MK\parallel  \left(SAB\right)$}
	 \loigiai{
		\begin{center}
			\begin{tikzpicture}[declare function={a=2; b=4; h=4;},line join=round]
					 \path (0,0) coordinate (A)
						(-145:a) coordinate  (B)
						 (b,0) coordinate (D)
						($(B)!0.5!(D)$)  coordinate (O)
						 ($(O)+(0,h)$) coordinate (S);
					 \path ($(D)-(A)+(B)$) coordinate (C);
					\path ($(C)!0.5!(B)$)  coordinate (I);
					\path  ($(S)!1/3!(D)$) coordinate (K);
					\path (intersection of A--I  and B--D) coordinate (M);
					 \draw[dashed] (S)--(A)--(B)
						 (A)--(D)--(B)
						(I)--(A)--(C) 
						(M)--(K); 
					\draw (B)--(C)--(D) 
						 (B)--(S)
						 (D)--(S)--(C);
					 \foreach \t/\g in {A/150,B/-90,C/-90,D/0,S/90,O/-90,I/-90,K/60,M/160}{
						\draw[fill=black] (\t)  circle (1pt) node[shift={(\g:7pt)},font=\scriptsize]{$ \t $};
					 }
				\end{tikzpicture}
		\end{center}
		\begin{itemize}
			\item  $MK\parallel \left(SCD\right)$ sai vì $MK\cap \left(SCD\right)=K$.
			\item $MK\parallel \left(SBD\right)$ sai vì $MK\subset  \left(SBD\right)$.
			\item $MK\parallel \left(ABCD\right)$  sai vì $MK\cap \left(ABCD\right)=M$.
			\item Ta có $M$ là trọng tâm tam giác  $ABC$ do đó $BM=\dfrac{2}{3} BO =\dfrac{1}{3} BD$.\\
			Suy ra $\dfrac{DK}{DS}=\dfrac{DM}{DB}=\dfrac{2}{3}  \Rightarrow MK\parallel SB$.\\
			Vậy  $MK\parallel \left(SAB\right)$.
		\end{itemize}
	 }
\end{ex}

% nguồn: Ok_Toan11_TracNghiem_Lan2
\begin{ex}%[1H4H3-2]%[ma:e7b8b4d5]
	Cho hai hình bình hành $ABCD$ và $ABEF$ không cùng nằm trong  một mặt phẳng. Gọi $O$, $O_1$ lần lượt là tâm của $ABCD$, $ABEF$. $M$ là trung điểm của $CD$. Khẳng định nào sau đây \textbf{sai}?
	\choice 
		{$OO_1 \parallel (  BEC) $}
		 {$OO_1 \parallel ( AFD) $}
		{$OO_1 \parallel ( EFM)  $}
		{\True  $MO_1$ cắt $( BEC) $}
	\loigiai{ 
		\begin{center}
			\begin{tikzpicture}[scale=1,line cap=round,line join=round,font=\footnotesize,>=stealth]
					\def\dai{4} \def\rong{2} \def\goc{60} 
					\path (0,0) coordinate (A)++(0:\dai) coordinate (B) 
						++(\goc:\rong) coordinate (C)++(180:\dai) coordinate  (D)
						($(A)+(-\goc+20:\rong)$) coordinate (F)++(0:\dai) coordinate  (E);
					\coordinate (O) at ($(A)!1/2!(C)$); 
					 \coordinate (O1) at ($(A)!1/2!(E)$);
					 \draw (A)--(F)--(E)--(C)--(D)--(A)
						 (D)--(F);
					 \draw [dashed] (A)--(B)--(C)--(A)--(E)--(B)--(D)
						 (B)--(F)
						 (O)--(O1);
					\foreach \t/\g in  {A/180,B/0,C/0,D/180,F/-90,E/-90,O/90,O1/-90}
					\draw[fill=black] (\t) circle  (1.0pt)
						 node[shift={(\g:9pt)}]{$\t$};
				\end{tikzpicture}
		\end{center}
		Xét tam giác $ACE$ có $O$, $O_1$ lần lượt  là trung điểm của $AC$, $AE$.\\
		Suy ra $OO_1$ là đường trung bình trong tam giác $ACE$ $ \Rightarrow  OO_1 \parallel EC$.\\
		Tương tự, $OO_1$ là đường trung bình của tam  giác $BFD$ nên $OO_1 \parallel FD$.\\
		Vậy $OO_1 \parallel ( BEC)$, $OO_1 \parallel ( AFD)$ và $OO_1 \parallel ( EFC)$. Chú ý rằng: $( EFC)=( EFM)  $.
	} 
\end{ex}

% nguồn: Ok_Toan11_TracNghiem_Lan2
\begin{ex}%[1H4H3-2]%[ma:0e0cb2e6]
	Cho tứ diện $ABCD.$ Gọi $G_1$, $G_2$ lần lượt là trọng tâm tam giác $BCD$ và $ACD$. Mệnh  đề nào sau đây \textbf{sai}?
	\choice 
		{$G_1G_2 \parallel (  ABD)$}
		{\True $G_1G_2=\dfrac{2}{3}AB$} 
		{$G_1G_2 \parallel (  ABC)$}
		{Ba đường $BG_1$,  $AG_2$, $CD$ đồng quy}
	\loigiai{ 
		\begin{center}
			\begin{tikzpicture}[scale=1,font=\footnotesize,line join=round,line cap=round]
					 \foreach \x\y\t in {0/0/B,1.3/-1.6/C,4.5/0/D,1/3.5/A}
					\coordinate (\t) at  (\x,\y);
					\coordinate (M) at  ($(D)!1/2!(C)$);
					\coordinate (G1) at  ($(M)!1/3!(B)$);
					\coordinate (G2)  at ($(M)!1/3!(A)$);
					\draw (M)--(A)--(B)--(C)--(A)--(D)--(C); 
					\draw[dashed]  (M)--(B)--(D)
						(G1)--(G2); 
					\foreach \t/\g in  {A/90,B/180,C/-90,D/0,G1/-90,G2/30,M/-60}
					\draw[fill=black]  (\t) circle (1.0pt)
						node[shift={(\g:9pt)}]{$\t$}; 
				\end{tikzpicture}
		\end{center}
		Gọi $M$ là trung điểm của $CD$.\\ 
		Do $G_1$, $G_2$ lần lượt là trọng tâm tam giác  $BCD$ và $ACD$ nên ta có
		$$\dfrac{MG_1}{MB}=\dfrac{MG_2}{MC}=\dfrac{1}{3} \Rightarrow  G_1G_2 \parallel AB,\dfrac{G_1G_2}{AB}=\dfrac{1}{3}.$$
		Do $G_1G_2 \parallel AB \Rightarrow G_1G_2 \parallel  ( ABD)$ và $G_1G_2 \parallel ( ABC)$.\\
		Dễ thấy, ba đường $BG_1,AG_2,CD$  đồng quy tại $M$. \\
		Và $\dfrac{MG_1}{MB}=\dfrac{MG_2}{MA}=\dfrac{1}{3}\Rightarrow  \dfrac{G_1G_2}{AB}=\dfrac{1}{3}$.\\
		Vậy mệnh  đề $G_1G_2=\dfrac{2}{3}AB$ là mệnh đề sai.
	} 
\end{ex}

% nguồn: Ok_Toan11_TracNghiem_Lan2
\begin{ex}%[1H4H3-2]%[ma:09a4a869]
	Cho tứ diện $ABCD$. Gọi $E$, $F$ lần lượt là trung điểm của các cạnh $AB$ và $AC$. Khi đó 
	\choice 
		{\True $EF\parallel (BCD$} 
		{$EF$ cắt  $(BCD$}
		 {$EF\parallel (ABD$}
		 {$EF\parallel (ABC$}
	\loigiai{ 
		\begin{center}
			\begin{tikzpicture}[scale=1,font=\footnotesize,line join=round,line cap=round,>=stealth,declare function={a=5; b=0.45*a; h=4;}]
					\path (0,0) coordinate  (B)
						(a,0) coordinate  (D)
						(-45:b) coordinate  (C)
						 ($(C)+(0,h)$) coordinate (A);
					\path ($(A)!0.5!(B)$) coordinate  (E);
					 \path ($(A)!0.5!(C)$) coordinate (F);
					\draw[dashed] (B)--(D); 
					\draw (A)--(C) 
						 (A)--(B)--(C)--(D)--cycle (E)--(F);
					\foreach \x/\goc in  {A/90,B/240,D/-60,C/-90,E/150,F/0}
					\draw[fill=white] (\x) node[shift={(\goc:7pt)},font=\scriptsize]{$\x$} circle (1pt); 
				\end{tikzpicture}
		\end{center}
		Do $E$, $F$ lần lượt là trung điểm của các cạnh $AB$ và $AC$  nên $EF\parallel BC$.\\
		Vì $BC\subset \left(BCD\right)$ và $EF\not\subset \left(BCD\right)$ nên $EF\parallel  (BCD)$.
	 }
\end{ex}

% nguồn: Ok_Toan11_TracNghiem_Lan2
\begin{ex}%[1H4H3-2]%[ma:cbc68acd]
	Cho tứ diện $ABCD$. Gọi $G$ là trong tâm của tam giác $ABD$, $Q$ thuộc cạnh $AB$ sao cho $AQ=2QB$ và $P$ là trung điểm của $AB$. Khẳng định nào sau  đây đúng?
	\choice 
		 {$GQ \parallel (ACD)$}
		{$GQ \parallel  (BCD)$}
		{$GQ$ cắt $(BCD)$} 
		{$Q$ thuộc mặt phẳng $(CDP)$} 
	\loigiai{ 
	\immini{Gọi $M$ là  trung điểm của $BD$.\\
		Vì $G$ là trọng tâm tam giác $ABD \Rightarrow \dfrac{AG}{AM}=\dfrac{2}{3}$.\\ 
		Điểm $Q \in  AB$ sao cho $AQ=2 QB \Leftrightarrow \dfrac{AQ}{AB}=\dfrac{2}{3}$.\\ Suy ra $\dfrac{AG}{AM}=\dfrac{AQ}{AB} \Rightarrow GQ \parallel BD$.\\
		Mặt khác $BD$ nằm trong mặt phẳng $(BCD)$ suy ra $GQ \parallel(BCD)$.}{\begin{tikzpicture}[scale=0.9,font=\footnotesize,line join=round,line cap=round,>=stealth]
		\coordinate  (B) at (0,0);
		\coordinate (D)  at (2,-2);
		\coordinate (C)  at (5,0);
		\coordinate (A)  at ($(B)+(60:4)$);
		\coordinate (Q)  at ($(A)!2/3!(B)$);
		\coordinate (P) at ($(B)!.5!(A)$); 
		\coordinate (M)  at ($(B)!.5!(D)$);
		\coordinate (G) at  ($(A)!2/3!(M)$);
		\draw (A)--(B)--(D)--(C)--(A)--(D)--(P) 
			 (A)--(M)
			 (Q)--(G);
		\draw [dashed] (B)--(C) 
			 (P)--(C);
		\foreach \x/\g in {A/90,B/180,C/0,D/-90,M/-110,P/120,G/30,Q/180} \fill (\x) circle (1pt) +(\g:3mm) node{$\x$}; 
	\end{tikzpicture}
		 }
	 }
\end{ex}

% nguồn: Ok_Toan11_TracNghiem_Lan2
\begin{ex}%[1H4H3-2]%[ma:63bf9e1e]
	Cho hình chóp $S.ABCD$ có  đáy là hình bình hành. Các điểm $I$, $J$ lần lượt là trọng tâm các tam giác và $SAB$ và $SAD$. Gọi $M$ là trung điểm $CD$. Chọn mệnh đề đúng trong các mệnh đề sau
	 \choice
		{\True  $IJ\parallel \left(SBD\right)$}
		{$IJ\parallel \left(SBM\right)$} 
		{$IJ\parallel \left(SBC\right)$} 
		{$IJ\parallel  \left(SCD\right)$}
	\loigiai{ 
	\immini{Gọi $N$, $P$ lần lượt là trung điểm của $AB$,  $AD$.\\
		Ta có $\dfrac{SI}{SN}=\dfrac{SJ}{SP}=\dfrac{1}{3} \Rightarrow IJ\parallel NP$ mà $NP\parallel  BD$ suy ra $IJ\parallel \left(SBD\right)$.
		}{\begin{tikzpicture}[scale=1,font=\footnotesize,line join=round,line cap=round,>=stealth,declare function={gocx=80; goc=-150; a=3; b=a/1.5; h=3;}]
		\path (0,0) coordinate (A)--+(gocx:h) coordinate (S) 
			 (a,0) coordinate (B)
			(goc:b) coordinate  (D)
			+(a,0)  coordinate (C);
		 \path ($(B)!0.5!(A)$) coordinate (N);
		 \path ($(D)!0.5!(A)$) coordinate (P);
		\path ($(D)!0.5!(C)$)  coordinate (M);
		\path ($(S)!2/3!(N)$) coordinate (I); 
		\path ($(S)!2/3!(P)$) coordinate  (J);
		\draw[dashed]  (S)--(A)
			(D)--(A)--(B) 
			(N)--(P) 
			 (B)--(D)
			(N)--(S)--(P)--(N) 
			 (I)--(J);
		 \draw (S)--(D)--(C)--(S)--(B)--(C)--cycle;
		 \foreach \x/\goc in {S/90,A/150,B/0,C/-60,D/180,N/60,P/150,I/0,J/180,M/-90}
		\draw[fill=white] (\x) node[shift={(\goc:7pt)}]{$\x$}  circle (1pt);
	\end{tikzpicture}
		 }
	 }
\end{ex}

% nguồn: Ok_Toan11_TracNghiem_Lan2
\begin{ex}%[1H4H3-2]%[ma:28cc57c1]
	Cho hình  chóp $S.ABCD$ có đáy $ABCD$ là hình thoi. Gọi $H$, $K$ lần lượt là trung điểm của $SA$, $CD$. Khẳng định nào sau đây đúng?
	 \choice
		 {\True $HK\parallel (SBC)$}
		 {$HK\parallel (SBD)$}
		{$HK\parallel (SAC)$} 
		 {$HK\parallel (SAD)$}
	\loigiai{ 
	\immini{Gọi $I$ là trung điểm $SB$, khi đó $IH$ là đường trung bình của tam giác $SAB$  nên $IH\parallel AB$ và $IH=\dfrac{1}{2}AB$.\\
		 Mặt khác $CK\parallel AB$ và $CK=\dfrac{1}{2}AB$.\\
		 Từ đó suy ra $IH\parallel CK$ và $IH=CK$ nên $IHKC$ là hình bình hành.\\
		 Suy ra $HK\parallel IC$.\\
		Khi đó ta  có $\heva{&HK\not\subset(SBC)\\& IC\subset(SBC)\\&HK \parallel IC}\Rightarrow HK \parallel (SBC)$.
		}{\begin{tikzpicture}[scale=0.8,font=\footnotesize,line join=round,line cap=round,>=stealth]
		\coordinate  (A) at (0,0);
		\coordinate (B) at  (-2,-2);
		\coordinate (D)  at (5,0);
		\coordinate (C) at ($(B)-(A)+(D)$); 
		\coordinate  (S) at ($(A)+(80:4)$);
		 \coordinate (H) at ($(S)!.5!(A)$);
		\coordinate (I)  at ($(S)!.5!(B)$);
		\coordinate (K) at ($(C)!.5!(D)$); 
		\draw  (S)--(B)--(C)--(D)--(S)--(C)--(I);
		\draw  [dashed] (S)--(A)--(D)
			 (B)--(A)
			 (K)--(H)--(I);
		\foreach \x/\g in  {A/180,B/-90,C/-90,S/90,D/0,H/20,I/180,K/-20} \fill (\x) circle (1pt) +(\g:3mm) node{$\x$};
	\end{tikzpicture}
		} 
	 }
\end{ex}

% nguồn: Ok_Toan11_TracNghiem_Lan2
\begin{ex}%[1H4H3-2]%[ma:8af9dbb4]
	Cho hình chóp $S.ABCD$ có đáy $ABCD$ là hình chữ nhật.  Gọi $M$, $N$ theo thứ tự là trọng tâm $\triangle SAB$, $\triangle SCD$, khi đó $MN$ song song với mặt phẳng
	 \choice
		 {$(SAC)$}
		{$(SBD)$} 
		{$(SAB)$} 
		 {\True $(ABCD)$}
	 \loigiai{
	\immini{Gọi  $E$, $F$ lần lượt là trung điểm $AB$ và $CD$.\\
		Ta có $\dfrac{SM}{SE}=\dfrac{SN}{SF}=\dfrac{2}{3}$ và $MN$, $EF$ đồng phẳng $\Rightarrow MN\parallel  EF$.\\
		Khi đó ta có $\heva{& MN \not \subset(ABCD)\\& EF \subset(ABCD)\\& MN\parallel EF} \Rightarrow MN\parallel (ABCD)$.}{\begin{tikzpicture}[scale=0.8,font=\footnotesize,line join=round,line cap=round,>=stealth]
		\coordinate (A)  at (0,0);
		 \coordinate (B) at (-2,-2);
		 \coordinate (D) at (4,0);
		\coordinate (C)  at ($(B)-(A)+(D)$);
		\coordinate (S)  at ($(A)+(100:4)$);
		\coordinate (E) at ($(A)!.5!(B)$); 
		 \coordinate (F) at ($(C)!.5!(D)$);
		\coordinate (M) at  ($(S)!2/3!(E)$);
		 \coordinate (N) at ($(S)!2/3!(F)$);
		 \draw (S)--(B)--(C)--(D)--(S)--(C)
			 (S)--(F);
		\draw [dashed] (S)--(A)--(D) 
			(B)--(A) 
			(S)--(E)--(F) 
			(M)--(N); 
		\foreach \x/\g in {A/180,B/-90,C/-90,S/90,D/0,E/180,F/-20,M/180,N/0} \fill (\x) circle (1pt) +(\g:3mm)  node{$\x$};
	\end{tikzpicture}
		 }
	 }
\end{ex}

% nguồn: Ok_Toan11_TracNghiem_Lan2
\begin{ex}%[1H4H3-2]%[ma:a1f16817]
	Cho hình chóp $S.ABCD$, có đáy $ABCD$ là hình bình hành tâm $O$. Gọi $H,K$ lần lượt là trung điểm của $BC,{\kern 1pt} CD$. Mệnh đề  nào dưới đây đúng?
	 \choice
		{\True $HK\parallel \left(SBD\right)$} 
		{$HK\parallel  \left(SAD\right)$}
		{$HK\parallel  \left(SCD\right)$}
		{$HK\parallel \left(SBC\right)$} 
	\loigiai{ 
	\immini{Ta  có $HK\not\subset \left(SBD\right)$.\\
		Ta thấy $HK$ là đường trung bình của tam giác $BCD$ nên $HK\parallel BD$ mà $BD\subset  \left(SBD\right)$.\\
		Do đó $HK\parallel \left(SBD\right)$. 
		}{\begin{tikzpicture}[scale=1,font=\footnotesize,line join=round,line cap=round,>=stealth,declare function={gocx=80; goc=-130; a=5; b=a/2; h=4;}]
		\path (0,0) coordinate (A)--+(gocx:h) coordinate (S) 
			(a,0) coordinate  (B)
			(goc:b) coordinate (D) 
			+(a,0) coordinate  (C);
		\path  ($(C)!0.5!(D)$) coordinate (K);
		 \path ($(C)!0.5!(B)$) coordinate (H);
		 \path ($(C)!0.5!(A)$) coordinate (O);
		\draw[dashed]  (S)--(A)
			 (D)--(A)--(B)
			 (A)--(C)
			(B)--(D) 
			(H)--(O)--(K)--(H); 
		\draw (S)--(D)--(C)--(S)--(B)--(C)--cycle; 
		 \foreach \x/\goc in {S/90,A/150,B/0,C/-60,D/-180,K/-90,H/-30,O/90}
		\draw[fill=white] (\x) node[shift={(\goc:7pt)}]{$\x$} circle (1pt); 
	\end{tikzpicture}
		 }
	 }
\end{ex}

% nguồn: Ok_Toan11_TracNghiem_Lan2
\begin{ex}%[1H4H3-2]%[ma:c2040289]
	Cho hình chóp $S.ABCD$ có đáy $ABCD$ là hình  bình hành. Đường thẳng $BC$ song song với mặt phẳng nào trong các mặt phẳng dưới đây?
	\choice 
		{\True $\left(SAD\right)$} 
		 {$\left(ABCD\right)$}
		 {$\left(SAC\right)$}
		 {$\left(SAB\right)$}
	 \loigiai{
		Vì $BC\parallel AD$ mà $AD\subset  (SAD)\Rightarrow BC\parallel (SAD)$.
	} 
\end{ex}

% nguồn: Ok_Toan11_TracNghiem_Lan2
\begin{ex}%[1H4H3-2]%[ma:681cb833]
	Cho hình hộp $ABCD.A'B'C'D'$. Gọi $M$ là điểm trên cạnh $AC$ sao cho $AM=3MC.$ Lấy $N$ trên cạnh $C'D$ sao cho $C'N=xC'D$.  Với giá trị nào của $ x$ thì $MN \parallel BD'$.
	 \choice
		{\True  $ x=\dfrac{2}{3}$}
		{ $ x=\dfrac{1}{3}$} 
		{  $ x=\dfrac{1}{4}$}
		{ $ x=\dfrac{1}{2}$} 
	 \loigiai{
		\begin{center}
			\begin{tikzpicture}[scale=1,font=\footnotesize,line join=round,line cap=round]
					 \def\h{4}
					\foreach \x\y\t in {0/0/A,-1/-1.1/B,2.6/-1.1/C} 
					\coordinate  (\t) at (\x,\y);
					\coordinate (D) at  ($(A)+(C)-(B)$);
					\coordinate (A') at  ($(A)+(0,2.5)$);
					\coordinate (B') at ($(B)+(0,2.5)$); 
					\coordinate  (C') at ($(C)+(0,2.5)$);
					\coordinate (D') at  ($(D)+(0,2.5)$);
					\coordinate (N) at ($(C')!3/4!(D)$); 
					\coordinate (M) at ($(A)!3/4!(C)$); 
					\draw  (B')--(A')--(D')--(C')--(B')--(B)--(C)--(D)--(D')
						 (C')--(C)--(D')
						 (D)--(C');
					\draw[dashed] (D')--(B)--(A)--(D)--(B) 
						 (C)--(A)--(A')
						(M)--(N); 
					\foreach \t/\g  in {A/170,B/-150,C/-70,D/0,A'/100,B'/170,C'/-130,D'/50,M/100,N/70}
					\draw[fill=black]  (\t) circle (1.0pt)
						 node[shift={(\g:9pt)}]{$\t$};
				\end{tikzpicture}
		\end{center}
		Từ giả thiết, ta có $M$ là trọng tâm tam giác $BCD$.\\ 
		Gọi $O$, $I$ lần lượt là  trung điểm của $AC$ và $DD'$. Khi đó, ta có $BD' \parallel ( IAC)$.\\
		Trong $( CDD'C')$ gọi $N'=CI \cap C'D$ khi đó $N'$ là trọng tâm tam giác $CDD'$.\\ 
		Do đó $\dfrac{CM}{CO}=\dfrac{2}{3}=\dfrac{CN'}{CI} \Rightarrow MN' \parallel OI$ mà $OI \parallel BD'$ nên  $MN \parallel BD'$. \\
		Vậy $N'\equiv N$  và $ x=\dfrac{2}{3}$.
	} 
\end{ex}

% nguồn: Ok_Toan11_TracNghiem_Lan2
\begin{ex}%[1H4H3-2]%[ma:a9fbe2b1]
	Cho hai hình bình hành  $ABCD$ và $ABEF$ nằm trong hai mặt phẳng khác nhau lần lượt có tâm $O$ và $O'$. Mệnh đề nào sau đây \textbf{sai}?
	\choice 
		{$OO'  \parallel (ADF)$}
		{$OO' \parallel  (BCE)$}
		{\True $OO'\parallel  (ACE)$}
		{$OO'\parallel  (DCEF)$}
	\loigiai{ 
	\immini{Do $OO' \subset (ACE)$ nên $OO'\parallel (ACE)$ sai.}{\begin{tikzpicture}[line cap=round,line join=round,scale=1.0]
		 \def\h{3}
		 \path
			 (0,0) coordinate (E)
			 (3.5,0) coordinate (C)
			(-50:2)  coordinate (B)
			($(E)+(90:\h)$) coordinate  (F)
			 ($(C)+(90:\h)$) coordinate (D)
			 ($(B)+(90:\h)$) coordinate (A);
		\coordinate (O') at ($(F)!.5!(B)$); 
		\coordinate  (O) at ($(C)!.5!(A)$);
		\draw[dashed]  (C)--(E)
			(O)--(O'); 
		\draw 
			 (F)--(A)--(D)--cycle
			 (F)--(E)--(B)--(C)--(D)
			(B)--(A) 
			 (F)--(B)--(D)
			(E)--(A)--(C); 
		\foreach \x/\g in 
			 {F/180,D/0,A/90,E/180,C/0,B/-40,O/0,O'/-180}
		 \fill(\x) circle (1pt)
			($(\x)+(\g:3mm)$)  node{$\x$};
	\end{tikzpicture}
		} 
	} 
\end{ex}

% nguồn: Ok_Toan11_TracNghiem_Lan2
\begin{ex}%[1H4H3-2]%[ma:81776bc8]
	Cho tứ diện $ABCD$, $G$ là trọng tâm $\triangle ABD$ và $M$ là điểm trên cạnh $BC$ sao cho  $BM=2MC$. Đường thẳng $MG$ song song với mặt phẳng
	 \choice
		{\True  $(ACD)$}
		 {$(ABC)$}
		{$(ABD)$} 
		{$(BCD)$} 
	\loigiai{ 
	\immini{Gọi $F$ là trung  điểm $AD$.\\
		Ta có $\dfrac{BG}{BF}=\dfrac{BM}{BC}=\dfrac{2}{3}  \Rightarrow MG\parallel CF$.\\
		Khi đó ta có $\heva{& MG \not \subset(ACD)\\& CF\subset(ACD) \\& MG\parallel CF}\Rightarrow MG\parallel (ACD)$.}{\begin{tikzpicture}[scale=0.8,font=\footnotesize,line join=round,line cap=round,>=stealth]
		\coordinate (B)  at (0,0);
		\coordinate (C)  at (2,-2);
		 \coordinate (D) at (5,0);
		\coordinate (A)  at ($(B)+(80:5)$);
		\coordinate (F) at ($(A)!.5!(D)$); 
		 \coordinate (G) at ($(B)!2/3!(F)$);
		 \coordinate (M) at ($(B)!2/3!(C)$);
		\draw (A)--(B)--(C)--(D)--(A)--(C)--(F); 
		\draw [dashed]  (B)--(D)
			(B)--(F) 
			(M)--(G); 
		\foreach \x/\g in {A/90,B/180,C/-90,D/0,F/30,M/-110,G/90} \fill (\x) circle  (1pt) +(\g:3mm) node{$\x$};
	\end{tikzpicture}
		 }
	} 
\end{ex}

% nguồn: Ok_Toan11_TracNghiem_Lan2
\begin{ex}%[1H4H3-2]%[ma:332a0a0a]
	Cho hình chóp $S.ABCD$ có đáy là hình bình hành. Gọi $I$ và $J$ lần lượt là trung điểm của $SC$ và $BC$. Chọn khẳng định đúng  trong các khẳng định sau
	\choice 
		{$JI\parallel (SAC)$} 
		{\True  $JI\parallel (SAB)$}
		{$JI\parallel (SBC)$} 
		{$JI\parallel (SAD)$} 
	 \loigiai{
	\immini{Ta có $IJ$ là đường trung bình của tam giác  $SBC$ nên $JI\parallel SB$.\\
		Khi đó ta có $\heva{&IJ\not\subset(SAB)\\& SB\subset(SAB)\\&IJ \parallel SB}\Rightarrow IJ \parallel (SAB)$.}{	\begin{tikzpicture}[scale=0.8,font=\footnotesize,line join=round,line cap=round,>=stealth]
		\coordinate  (A) at (0,0);
		 \coordinate (B) at (-3,-2);
		 \coordinate (D) at (5,0);
		 \coordinate (C) at ($(B)-(A)+(D)$);
		\coordinate (S)  at ($(A)+(70:5)$);
		\coordinate (I) at ($(S)!.5!(C)$); 
		\coordinate (J) at  ($(B)!.5!(C)$);
		\draw (S)--(B)--(C)--(D)--(S)--(C) 
			(I)--(J); 
		\draw [dashed]  (S)--(A)--(D)
			 (B)--(A);
		 \foreach \x/\g in {A/160,B/-90,C/-90,S/90,D/0,I/0,J/-90} \fill (\x) circle (1pt) +(\g:3mm) node{$\x$};
	\end{tikzpicture}
		} 
	 }
\end{ex}

% nguồn: Ok_Toan11_TracNghiem_Lan2
\begin{ex}%[1H4H3-2]%[ma:5cf2d40b]
	Cho hình chóp $S.ABCD$ có đáy $ABCD$ là hình bình hành. Gọi $M$, $N$ lần lượt là trung điểm của $SA$, $SB$. Mệnh đề nào sau đây là mệnh đề \textbf{đúng} 
	 \choice
		{$MN\parallel (SAC)$} 
		 {$MN\parallel (SAB)$}
		{$MN\parallel  (SBD)$}
		 {\True $MN\parallel (ACD)$}
	 \loigiai{
	\immini{  Ta có
		\begin{itemize}
			\item $MN$ cắt $(SAC)$  tại $M$.
			\item $MN$  chứa trong $(SAB)$.
			\item $MN$  cắt $(SBD)$ tại $N$.
			 \item Theo tính chất đường trung bình $MN\parallel AB\parallel CD$ mà $CD\subset (ACD)$ nên $MN\parallel (ACD)$.
		\end{itemize}
		}{\begin{tikzpicture}[scale=1,font=\footnotesize,line join=round,line cap=round,>=stealth,declare function={gocx=80; goc=-130; a=4; b=a/2; h=3;}]
		\path (0,0) coordinate (A)--+(gocx:h) coordinate (S) 
			(a,0)  coordinate (B)
			(goc:b)  coordinate (D)
			 +(a,0) coordinate (C);
		\path ($(S)!0.5!(A)$) coordinate (M); 
		\path ($(S)!0.5!(B)$) coordinate (N); 
		\draw[dashed] (S)--(A) 
			 (D)--(A)--(B)
			(M)--(N); 
		\draw  (S)--(D)--(C)--(S)--(B)--(C)--cycle;
		 \foreach \x/\goc in {S/90,A/150,B/0,C/-60,D/180,M/-45,N/30}
		\draw[fill=white] (\x) node[shift={(\goc:7pt)}]{$\x$}  circle (1pt);
	\end{tikzpicture}
		 }
	} 
\end{ex}

% nguồn: Ok_Toan11_TracNghiem_Lan2
\begin{ex}%[1H4H3-2]%[ma:f3fd0495]
	Cho tứ diện $ABCD$. Gọi $G, G'$ lần lượt là trọng tâm $\triangle ABD$ và $\triangle BCD$. Khẳng định  nào sau đây là \textbf{sai}?
	 \choice
		 {$GG' \parallel(ACD)$}
		 {\True $GG' \parallel BD$}
		 {$GG' \parallel(ABC)$}
		{$GG'  \parallel AC$}
	 \loigiai{
		\begin{center}
			\begin{tikzpicture}[line join = round,line cap = round,>=stealth,font=\footnotesize,scale=1,declare function={bc=3; ab=4.5; gocB=-30;}]
					\path (0,0) coordinate (B) 
						(gocB:bc)  coordinate (C)
						 (ab,0) coordinate (A)
						(50:4) coordinate  (D)
						($(D)!0.5!(A)$)  coordinate (N)
						($(D)!0.5!(C)$)  coordinate (M)
						 ($(B)!2/3!(N)$) coordinate (G)
						 ($(B)!2/3!(M)$) coordinate (G');
					\draw  (D)--(B)--(C)--(A)--(D)--(C)
						(B)--(M)--(N); 
					\draw[dashed]  (A)--(B)
						(G)--(G') 
						 (B)--(N);
					\foreach \x/\gm in {A/-90,B/-90,C/-90,D/90,N/0,M/-50,G/90,G'/-90} \draw[fill=black] (\x)  circle (1pt)
						 ($(\x)+(\gm:3.5mm)$) node[black]{$\x$};
				\end{tikzpicture}
		\end{center}
		Ta có $GG'$ cắt mặt phẳng ($ABD$) tại $G$. Do đó $GG'$ không thể song song được với $BD$ nằm trong mặt phẳng  $(ABD)$.
	 }
\end{ex}

% nguồn: Ok_Toan11_TracNghiem_Lan2
\begin{ex}%[1H4H3-2]%[ma:0b0820dc]
	\immini{Cho  hình lăng trụ $ABCD.A'B'C'D'$ có hai đáy là các hình bình hành. Các điếm $M$, $N$, $P$ lần lượt là trung điếm của cạnh $AD$, $BC$, $CC'$.
	\begin{itemize}
		\item [i)] $A'B' \parallel  (MNP)$.
		\item [ii)] $(MNP) \parallel  (BC'D')$
		\item [iii)] $(MNP)  \parallel (B'C'D')$.
		\item [iv)] $DD'$ cắt mặt  phẳng $(MNP)$.
	\end{itemize}
	Trong các mệnh đề trên có bao nhiêu mệnh đề \textbf{sai}?}{\begin{tikzpicture}[scale=0.7,font=\footnotesize,line join=round,line cap=round,>=stealth]
		 \def\bc{4}
		\def\ba{2} 
		 \def\gocB{35}
		\coordinate[label=below left:$B$] (B)  at (0,0);
		\coordinate[label=above left:$A$] (A) at (\gocB:\ba); 
		\coordinate[label=below:$C$] (C) at  (\bc,0);
		\coordinate[label=right:$D$] (D) at  ($(C)-(B)+(A)$);
		\coordinate[label=above left:$A'$] (A') at  ($(A)+(90:\bc)$);
		\coordinate[label=left:$B'$] (B') at  ($(B)-(A)+(A')$);
		\coordinate[label=below right:$C'$] (C') at  ($(C)-(A)+(A')$);
		 \coordinate[label=right:$D'$] (D') at ($(D)-(A)+(A')$);
		\coordinate[label=above right:$P$] (P)  at ($(C)!.5!(C')$);
		\coordinate[label=below:$N$] (N) at  ($(B)!.5!(C)$);
		\coordinate[label=below:$M$] (M)  at ($(D)!.5!(A)$);
		\draw  (B')--(B)--(C)--(D)--(D')--(A')--(B')--(C')--(D')
			 (C)--(C')
			(N)--(P); 
		 \draw[dashed] (A')--(A)--(D)
			(A)--(B) 
			(N)--(M)--(P); 
		\foreach \diem in {A,B,C,D,A',B',C',D',P,M,N} \fill (\diem) circle (1pt); 
	\end{tikzpicture}
	} 
	 \choice
		 {$4$}
		{$2$} 
		{$3$} 
		 {\True $1$}
	 \loigiai{
		Ta có $\heva{&A'B'\parallel AB\\&AB \parallel MN} \Rightarrow A'B' \parallel MN \Rightarrow  A'B' \parallel (MNP)$.\\
		Lại có $\heva{&MN\parallel C'D'\\&NP \parallel  BC'} \Rightarrow (MNP) \parallel (BC'D')$.\\
		Hơn  nữa $\heva{&(MNP) \cap (ABCD)=MN\\& (B'C'D') \parallel (ABCD)} \Rightarrow (MNP)$ cắt $(B'C'D')$.\\
		Và $\heva{&(MNP)\parallel (BC'D')\\&DD' \cap (BC'D')\equiv D'} \Rightarrow  DD' \cap (MNP) \equiv Q$.\\
		Vậy chỉ có mệnh đề  iii) \textbf{sai}.
	 }
\end{ex}

% nguồn: Ok_Toan11_TracNghiem_Lan2
\begin{ex}%[1H4H3-2]%[ma:135561f1]
	Cho hình chóp $S.ABCD$ có đáy $ABCD$ là hình bình hành. Gọi $M$ là một điểm trên cạnh $SA$, mặt phẳng $\left(\alpha \right)$ qua $M$ song song với $SB$ và $AC$. Mặt phẳng $\left(\alpha \right)$ cắt $AB$, $BC$, $SC$, $SD$, $BD$ lần lượt tại $N$, $E$, $F$, $I$,  $J$. Khẳng định nào sau đây là đúng?
	\choice 
		 {$MN\; \parallel\left(SCD\right)$}
		{$EF\; \parallel\left(SAD\right)$} 
		{$NF\;  \parallel\left(SAD\right)$}
		{\True $IJ\; \parallel\left(SAB\right)$} 
	\loigiai{ 
		\begin{center}
			\begin{tikzpicture}[declare function={a=2; b=4; h=4;},line join=round]
					 \path (0,0) coordinate (A)
						 (-145:a) coordinate (B)
						(b,0) coordinate  (D)
						 ($(B)!0.5!(D)$) coordinate (O)
						($(O)+(0,h)$)  coordinate (S);
					\path  ($(D)-(A)+(B)$) coordinate (C);
					\path ($(A)!0.4!(S)$)  coordinate (M);
					\path ($(A)!0.4!(B)$) coordinate  (N);
					\path ($(B)!0.6!(C)$) coordinate  (E);
					 \path ($(C)!0.4!(S)$) coordinate (F);
					\path ($(C)!0.4!(S)$)  coordinate (F);
					\path ($(D)!0.6!(S)$)  coordinate (I);
					\path (intersection of B--D and N--E)  coordinate (J);
					\draw[dashed] (S)--(A)--(B) 
						 (A)--(D)--(B)
						 (A)--(C)
						(E)--(N)--(M)--(I); 
					\draw  (B)--(C)--(D)
						(B)--(S) 
						 (D)--(S)--(C)
						 (E)--(F)--(I);
					\foreach \t/\g  in {A/30,B/-90,C/-90,D/0,S/90,M/180,N/160,E/-90,F/0,I/60,J/90}{
						\draw[fill=white]  (\t) circle (1pt) node[shift={(\g:7pt)},font=\scriptsize]{$ \t $};
					} 
				\end{tikzpicture}
		\end{center}
		Ta có: $\heva{& IJ=\left(\alpha \right)\cap \left(SBD\right) \\ &\left(\alpha \right)\parallel  SB\subset \left(SBD\right)} \Rightarrow \left(\alpha \right)\cap \left(SBD\right)=IJ\parallel SB$. \\
		Mà  $SB\subset (SAB)\Rightarrow IJ\parallel (SAB)$.
	 }
\end{ex}

% nguồn: Ok_Toan11_TracNghiem_Lan2
\begin{ex}%[1H4H3-2]%[ma:196e4d7f]
	Cho hình chóp $S.ABCD$ có đáy $ABCD$ là hình chữ nhật. Gọi $M$, $N$ theo thứ tự là trọng tâm tam giác $SAB$ và tam giác $SCD$. Khi đó $MN$ song song với mặt  phẳng
	 \choice
		 {$(SAC$}
		 {$(SBD$}
		 {$(SAB)$}
		{\True  $(ABCD)$}
	 \loigiai{
		\begin{center}
			\begin{tikzpicture}[declare function={a=2; b=4; h=4;},line join=round]
					\path (0,0)  coordinate (A)
						 (-145:a) coordinate (B)
						 (b,0) coordinate (D)
						($(B)!0.5!(D)$)  coordinate (O)
						($(O)+(0,h)$) coordinate (S); 
					 \path ($(D)-(A)+(B)$) coordinate (C);
					\path ($(A)!0.5!(B)$) coordinate (E); 
					\path  ($(C)!0.5!(D)$) coordinate (F);
					\path ($(S)!2/3!(E)$)  coordinate (M);
					 \path ($(S)!2/3!(F)$) coordinate (N);
					\draw[dashed]  (F)--(E)--(S)--(A)--(B)
						 (A)--(D)
						(M)--(N); 
					 \draw (B)--(C)--(D)
						(B)--(S)--(F) 
						 (D)--(S)--(C);
					\foreach \t/\g in  {A/50,B/-90,C/-90,D/0,S/90,E/-90,F/-30,M/-60,N/30}{
						 \draw[fill=white] (\t) circle (1pt) node[shift={(\g:7pt)},font=\scriptsize]{$ \t $};
					 }
				\end{tikzpicture}
		\end{center}
		Gọi $E$ và $F$ lần lượt là trung điểm  $AB$ và $C$
		Do  $M$, $N$ là trọng tâm tam giác $SAB$, $SCD$ nên $S$, $M$, $E$ thẳng hàng và $S$, $N$, $F$ thẳng hàng.\\
		 Xét tam giác $SEF$ có $\dfrac{SM}{SE}=\dfrac{2}{3}=\dfrac{SN}{SF}$ nên theo định lý Ta--let $\Rightarrow MN\parallel EF$.\\
		Mà $EF\subset \left(ABCD\right)$ nên $MN\parallel \left(ABCD\right)$. 
	} 
\end{ex}

% nguồn: Ok_Toan11_TracNghiem_Lan2
\begin{ex}%[1H4H3-2]%[ma:67bc6586]
	Cho tứ diện $ABCD$ có $M,N$ lần lượt là trung điểm của $AB,AC$. Đường  thẳng $MN$ song song với mặt phẳng nào sau đây?
	 \choice
		{$\left(ABC\right)$} 
		{\True  $\left(BCD\right)$}
		{$\left(ABD\right)$} 
		 {$\left(ACD\right)$}
	\loigiai{ 
		\begin{center}
			\begin{tikzpicture}[scale=1,font=\footnotesize,line join=round,line cap=round,>=stealth,declare function={a=5; b=0.45*a; h=4;}]
					 \path (0,0) coordinate (B)
						(a,0) coordinate (D) 
						(-45:b) coordinate (C) 
						 ($(C)+(0,h)$) coordinate (A);
					\path  ($(A)!0.5!(B)$) coordinate (M);
					\path ($(A)!0.5!(C)$) coordinate  (N);
					\draw[dashed] (B)--(D); 
					\draw (A)--(C) 
						(A)--(B)--(C)--(D)--cycle (M)--(N); 
					\foreach  \x/\goc in {A/90,B/240,D/-60,C/-90,M/150,N/0}
					 \draw[fill=white] (\x) node[shift={(\goc:7pt)},font=\scriptsize]{$\x$} circle (1pt);
				\end{tikzpicture}
		\end{center}
		Vì $MN$ là đường trung bình của tam giác $ABC$ nên $MN\parallel BC$.  Mà $BC\subset \left(BCD\right)$ nên $MN\parallel \left(BCD\right)$.
	} 
\end{ex}

% nguồn: Ok_Toan11_TracNghiem_Lan2
\begin{ex}%[1H4H3-2]%[ma:4171b645]
	Cho hình chóp $S.ABCD$ có đáy là hình bình hành. Các điểm $I$, $J$ lần lượt là trọng tâm các  tam giác $SAB$, $SAD$. Chọn mệnh đề đúng trong các mệnh đề sau
	 \choice
		{$IJ\parallel  (SCD$}
		{$IJ\parallel (SBM)$} 
		{$IJ\parallel  (SBC$}
		 {\True $IJ\parallel (SBD$}
	\loigiai{ 
		\begin{center}
			\begin{tikzpicture}[declare function={a=2; b=4; h=4;},line join=round]
					\path  (0,0) coordinate (B)
						 (35:a) coordinate (A)
						 (b,0) coordinate (C)
						($(C)-(B)+(A)$) coordinate  (D)
						 ($(A)!1/2!(B)$) coordinate (N)
						 ($(A)!1/2!(D)$) coordinate (P)
						($(N)+(90:h)$) coordinate  (S);
					\path ($(S)!2/3!(N)$)  coordinate (I);
					\path ($(S)!2/3!(P)$) coordinate (J); 
					\draw (B)--(C)--(D)--(S)--cycle  (S)--(C);
					 \draw[dashed] (A)--(D)--(B)
						(S)--(P)--(N)--(S)--(A)--(B) 
						 (I)--(J);
					\foreach \t/\g in  {A/150,B/-90,C/-90,D/0,S/90,N/180,P/40,I/180,J/60}{
						\draw[fill=black] (\t) circle (1pt) node[shift={(\g:7pt)},font=\scriptsize]{$ \t  $};
					} 
				\end{tikzpicture}
		\end{center}
		Gọi $N$, $P$ lần lượt là trung điểm  của cạnh $AB$, $AD$.\\
		Xét tam  giác $SNP$ có $\dfrac{SI}{SN}=\dfrac{SJ}{SP}=\dfrac{2}{3} \Rightarrow IJ\parallel NP$.\\
		Xét tam giác $ABD$ có $NP$ là đường trung bình trong tam  giác $\Rightarrow NP\parallel BD$.\\
		Suy ra $IJ\parallel  BD$.
		Ta có  $\heva{&IJ\not\subset (SBD) \\ & IJ\parallel BD\\& BD\subset (SBD)} \Rightarrow IJ\parallel (SBD)$.
	 }
\end{ex}

% nguồn: Ok_Toan11_TracNghiem_Lan2
\begin{ex}%[1H4H3-2]%[ma:ed91e1ee]
	 Cho hình chóp $S.ABCD$. Gọi $M$, $N$ lần lượt là trọng tâm tam giác $SAB$ và tam giác $SCD$. Khi đó $MN$ song song với mặt phắng
	\choice 
		 {$(SAC)$}
		{$(SBD)$} 
		{$(SAB)$} 
		 {\True $(ABCD)$}
	 \loigiai{
	\immini{Gọi $E$ và $F$ lần lượt là trung điểm của $AB$ và  $CD$.\\
		Do $M$, $N$ là trọng tâm $\triangle SAB$, $\triangle SCD$ nên $S$, $M$, $E$ thẳng hàng và  $S$, $N$, $F$ thẳng hàng.\\
		Xét $\triangle SEF$ có  $\dfrac{SM}{SE}=\dfrac{2}{3}=\dfrac{SN}{SF}$ nên theo định lý Ta-lét $\Rightarrow MN\parallel EF$.\\
		Mà $EF \subset(ABCD)$ nên $MN\parallel (ABCD)$.}{\begin{tikzpicture}[scale=0.8,font=\footnotesize,line join=round,line cap=round,>=stealth]
		 \coordinate (A) at (0,0);
		\coordinate (D)  at (5,0);
		\coordinate (C) at  (4,-3);
		\coordinate  (B) at (1,-2);
		\coordinate (S) at  ($(A)+(60:5)$);
		\coordinate (E) at  ($(A)!.5!(B)$);
		 \coordinate (F) at ($(C)!.5!(D)$);
		\coordinate (M) at ($(S)!2/3!(E)$); 
		\coordinate (N) at  ($(S)!2/3!(F)$);
		\draw (B)--(S)--(A)--(B)--(C)--(D)--(S)--(C) 
			(E)--(S)--(F); 
		\draw  [dashed] (A)--(D)
			(M)--(N) 
			 (E)--(F);
		\foreach \x/\g in {A/180,B/-90,C/-90,D/0,M/180,N/0,S/90,E/200,F/0} \fill (\x)  circle (1pt) +(\g:3mm) node{$\x$};
	\end{tikzpicture}
		 }
	 }
\end{ex}

% nguồn: Ok_Toan11_TracNghiem_Lan2
\begin{ex}%[1H4H3-2]%[ma:a6283d28]
	Cho hình chóp $SABCD$ có đáy là hình bình hành. $M$, $N$ lần lượt là trung điểm của $SC$ và $SD$. Mệnh đề nào sau đây là đúng? 
	 \choice
		{$MN\parallel \left(SBD\right)$} 
		{\True  $MN\parallel \left(SAB\right)$}
		{$MN\parallel \left(SAC\right)$} 
		 {$MN\parallel \left(SCD\right)$}
	\loigiai{ 
		\begin{center}
			\begin{tikzpicture}[declare function={a=2; b=4; h=4;},line join=round]
					\path  (0,0) coordinate (A)
						(-145:a) coordinate (B) 
						(b,0)  coordinate (D)
						($(B)!0.5!(D)$) coordinate (O) 
						($(O)+(0,h)$) coordinate  (S);
					\path ($(D)-(A)+(B)$)  coordinate (C);
					\path ($(S)!0.5!(D)$) coordinate (N); 
					\path ($(S)!0.5!(C)$) coordinate  (M);
					 \draw[dashed] (S)--(A)--(B)
						(A)--(D); 
					 \draw (B)--(C)--(D)
						(B)--(S) 
						 (D)--(S)--(C)
						(M)--(N); 
					\foreach \t/\g  in {A/150,B/-90,C/-90,D/0,S/90,M/200,N/30}{
						\draw[fill=white] (\t) circle (1pt) node[shift={(\g:7pt)},font=\scriptsize]{$ \t $}; 
					 }
				\end{tikzpicture}
		\end{center}
		Ta có $MN\parallel CD\Rightarrow MN\parallel  AB
		 \Rightarrow MN\parallel \left(SAB\right)$.
	 }
\end{ex}

% nguồn: Ok_Toan11_TracNghiem_Lan2
\begin{ex}%[1H4H3-2]%[ma:3e853921]
	Cho hình chóp tứ giác $S.ABCD$. Gọi $M$, $N$ lần lượt là trung điểm của $SA$  và $SC$. Mệnh đề nào sau đây đúng?
	\choice 
		{$MN\parallel (SAB)$} 
		{$MN\parallel  (SBC)$}
		{$MN\parallel (SBD)$} 
		{\True  $MN\parallel (ABCD)$}
	 \loigiai{
	\immini{Vì $MN$ là đường trung bình của tam giác $SAC$ nên $MN\parallel  AC$.\\
		Khi đó ta có $\heva{&MN\not\subset(ABCD)\\& AC\subset(SBD)\\&MN \parallel AC}\Rightarrow MN \parallel (ABCD)$.}{\begin{tikzpicture}[scale=0.8,font=\footnotesize,line join=round,line cap=round,>=stealth]
		\coordinate (A) at  (0,0);
		\coordinate (B) at  (5,0);
		 \coordinate (C) at (4,-2);
		\coordinate (D) at (2,-3); 
		\coordinate (S) at ($(A)+(60:4)$); 
		\coordinate (M) at ($(S)!.5!(A)$); 
		\coordinate (N) at ($(S)!.5!(C)$); 
		\draw  (D)--(S)--(A)--(D)--(C)--(B)--(S)--(C);
		\draw  [dashed] (A)--(B)
			(A)--(C) 
			 (M)--(N);
		 \foreach \x/\g in {A/180,B/0,C/0,D/-90,M/180,N/0,S/90} \fill (\x) circle (1pt) +(\g:3mm) node{$\x$};
	\end{tikzpicture}
		 }
	} 
\end{ex}

% nguồn: Ok_Toan11_TracNghiem_Lan2
\begin{ex}%[1H4H3-2]%[ma:5241f229]
	Cho tứ diện $ABCD$. Gọi $I$, $J$ lần lượt là trung điểm của các cạnh $AB$ và $BC$. Trong các mệnh đề sau, mệnh đề nào đúng 
	 \choice
		{$IJ\parallel (BCD)$} 
		{$IJ\parallel  (ABC)$}
		 {\True $IJ\parallel (ACD)$}
		{$IJ\parallel (ABD)$} 
	 \loigiai{
	 \immini{Xét tam giác $ABC$ có $I$, $J$ lần lượt là trung điểm của các cạnh $AB$ và $BC$ suy ra $IJ\parallel AC$.\\
		Khi đó ta có $\heva{& IJ \not \subset(ACD)\\& AC \subset(ACD)\\& IJ\parallel CD} \Rightarrow IJ\parallel (ACD)$.}{\begin{tikzpicture}[scale=0.8,font=\footnotesize,line join=round,line cap=round,>=stealth]
		 \coordinate (A) at (0,0);
		 \coordinate (B) at (3,-2);
		\coordinate (C) at  (5,0);
		\coordinate (D) at  ($(A)+(60:4)$);
		\coordinate (I) at ($(A)!.5!(B)$); 
		\coordinate (J) at ($(B)!.5!(C)$); 
		 \draw (D)--(A)--(B)--(C)--(D)--(B);
		 \draw [dashed] (A)--(C)
			 (I)--(J);
		\foreach \x/\g in {A/180,B/-90,C/0,D/90,I/210,J/-20} \fill (\x) circle (1pt) +(\g:3mm)  node{$\x$};
	\end{tikzpicture}
		 }
	 }
\end{ex}

% nguồn: Ok_TracNghiemLop11_Lan1
\begin{ex}%[1H4H3-2]%[ma:1f4089e4]
	Cho tứ diện $ABCD$ với $M, N$ lần lượt là trọng tâm các tam giác $ABD$, $ACD$. Xét các khẳng  định sau
	\begin{enumEX}[]{2}
		\item (I). $MN\parallel \left(ABC\right)$ 
		\item  (II). $MN\parallel \left(BCD\right)$
		\item (III). $MN\parallel  \left(ACD\right)$.
		\item (IV). $MN\parallel  \left(ABD\right)$.
	\end{enumEX}
	Các mệnh đề đúng là 
	 \choice
		{(I), (IV)} 
		{(II),  (III)}
		{(III),  (IV)}
		{\True (I), (II)} 
	\loigiai{ 
		\begin{center}
			\begin{tikzpicture}[font=\footnotesize,scale=.8]
					\def\a{4} 
					\path (0:0) coordinate  (B)
						++(0:\a) coordinate (C) 
						++(-120:\a/2)  coordinate (D)
						($(B)+(70:\a)$) coordinate  (A)
						($(B)!.5!(D)$)  coordinate (I)
						 ($(C)!.5!(D)$) coordinate (K)
						 ($(A)!2/3!(I)$) coordinate (M)
						($(A)!2/3!(K)$) coordinate  (N);
					\draw[dashed]  (B)--(C)
						(I)--(K) 
						(M)--(N); 
					\draw (B)--(A)--(D) 
						 (A)--(C)
						(B)--(D)--(C) 
						(I)--(A)--(K); 
					\foreach \x/\g in {A/90,B/180,C/0,D/-45,I/-135,K/-45,M/175,N/0} 
					\fill[black]  (\x) circle (1pt)
						 ($(\g:3mm)+(\x)$) node {$\x$};
				\end{tikzpicture}
		\end{center}
		Gọi $I,K$ lần lượt là trung điểm của $BD,DC$. 
		\begin{itemize}
			\item (II)--  Đúng. Xét tam giác $AIK$ có $\heva{&{MN\parallel IK} \\&{IK\subset (BCD)} \\&{MN\not\subset (BCD)}} \Rightarrow MN\parallel (BCD)$.
			\item (I)-- Đúng. $\heva{&{MN\parallel IK}  \\&{IK\parallel BC}} \Rightarrow MN\parallel BC$ và $MN\not\subset (ABC)$ do đó $MN\parallel (ABC)$.
			\item  (III)-- Đúng. Có $M\in (ABD),N\in (ACD)$.
			\item  (IV)-- Sai.
		\end{itemize}
	 }
\end{ex}

% nguồn: Ok_TracNghiemLop11_Lan1
\begin{ex}%[1H4H3-2]%[ma:6fb60453]
	Cho hình chóp $S.ABCD$ đáy $ABCD$ là hình chữ nhật tâm $O$. Gọi $M$ là trung  điểm của $OC$. Mặt phẳng $(\alpha)$ qua $M$ và $(\alpha)$ song song với $SA$ và $BD$. Thiết diện của hình chóp $S.ABCD$ và mặt phẳng $(\alpha)$ là hình gì?
	\choice 
		 {\True hình tam giác}
		{hình bình hành} 
		 {hình chữ nhật}
		{hình ngũ  giác}
	\loigiai{ 
	\immini{ 
		Ta  có
		\begin{itemize}
			\item $\heva{&M\in(\alpha)\cap(SAC)\\&SA\parallel (\alpha)}\Rightarrow (\alpha)\cap (SAC)=Mx\parallel  SA$.
			 \item $\heva{&M\in(\alpha)\cap(ABCD)\\&BD\parallel (\alpha)}\Rightarrow (\alpha)\cap (ABCD)=My\parallel BD$.
		\end{itemize}
		Gọi $E$, $F$, $G$ lần lượt là trung điểm cạnh $BC$, $CD$ và $SC$  thì
		\begin{itemize}
			\item  $(\alpha)\cap(ABCD)=EF$.
			 \item $(\alpha)\cap(SCD)=FG$.
			$(\alpha)\cap(SBC)=GE$. 
		\end{itemize}
		Vậy $(\alpha)$ cắt hình chóp $S.ABCD$  theo thiết diện là $\triangle EFG$.
		 }
		{\begin{tikzpicture}[scale=.6,font=\footnotesize]
		\def\a{4} 
		\path  (0:0) coordinate (A)
			 ++(0:\a) coordinate (D)
			++(-125:.65*\a) coordinate  (C)
			($(A)+(C)-(D)$) coordinate  (B)
			 ($(A)+(85:\a)$) coordinate (S)
			(intersection of A--C and  B--D) coordinate (O)
			($(O)!.5!(C)$) coordinate  (M)
			($(S)!.75!(C)$)  coordinate (G)
			($(B)!.5!(C)$) coordinate (E) 
			($(D)!.5!(C)$)  coordinate (F);
		 \draw[dashed] (B)--(A)--(D)
			(A)--(S) 
			 (A)--(C)
			 (B)--(D)
			(M)--(G) 
			 (E)--(F);
		\draw[thick] (B)--  (C)--(D)
			(E)--(G)--(F) 
			(B)--(S) 
			(C)--(S) 
			(D)--(S); 
		\foreach \x/\g in  {A/135,B/-135,C/-45,D/45,S/90,O/90,M/-90,E/-90,F/-45,G/45}
		\fill[black] (\x)  circle (1pt)
			($(\g:4mm)+(\x)$)  node {$\x$};
	\end{tikzpicture}
		} 
	} 
\end{ex}

% nguồn: Ok_TracNghiemLop11_Lan1
\begin{ex}%[1H4H3-2]%[ma:81fd8fbd]
	Cho hình chóp tứ giác $S.ABCD$. Gọi $M$ và lần lượt là  trung điểm của $SA$ và $SC$. Khẳng định nào sau đây đúng?
	\choice 
		{\True $MN\parallel  (ABCD)$}
		{$MN\parallel (SAB)$} 
		 {$MN\parallel (SCD)$}
		 {$MN\parallel (SBC)$}
	 \loigiai{
	\immini{Ta có $MN$ là đường trung bình của $\triangle SAC$  nên $MN\parallel AC$.\\
		 Mà $MN\not\subset (ABCD)$, $CD\subset (ABCD)$ nên $MN\parallel (ABCD)$. }
		{\begin{tikzpicture}
		\def\a{3} 
		\path  (0:0) coordinate (A)
			++(0:\a) coordinate  (D)
			++(-130:.6*\a) coordinate (C) 
			++(-170:.5*\a) coordinate  (B)
			($(A)+(70:\a)$) coordinate (S) 
			($(A)!.5!(S)$) coordinate  (M)
			($(C)!.5!(S)$) coordinate (N); 
		\draw[dashed] (C)--(A)--(D) 
			 (M)--(N);
		\draw  (A)--(B)--(C)--(D)
			 (A)--(S)
			 (B)--(S)
			 (C)--(S)
			 (D)--(S);
		\foreach  \x/\g in {A/180,B/-135,C/-45,D/0,S/90,M/170,N/0}
		 \fill[black] (\x) circle (1pt)
			($(\g:3mm)+(\x)$)  node {$\x$};
	\end{tikzpicture}
		 }
	} 
\end{ex}

% nguồn: Ok_TracNghiemLop11_Lan1
\begin{ex}%[1H4H3-2]%[ma:8179c728]
	Cho  tứ diện $ABCD$, $G$ là trọng tâm tam giác $ABD$. Trên đoạn $BC$ lấy điểm $M$ sao cho $MB=2MC$. Khẳng định nào sau đây đúng?
	 \choice
		 {\True $MG$ song song $(ACD)$}
		{$MG$ song  song $(ABD)$}
		 {$MG$ song song $(ACB)$}
		{$MG$ song  song $(BCD)$}
	\loigiai{ 
	 \immini{Gọi $N$ là trung điểm của $AD$.\\
		 Ta có
		\begin{itemize}
			\item Vì $G$ là trọng tâm $\triangle  ABD$ nên $\dfrac{BG}{BN}=\dfrac{2}{3}$.\tagEX{1}
			\item $\dfrac{BM}{BC}=\dfrac{2}{3}$  vì $MB=2MC$.\tagEX{2}
		\end{itemize}
		Từ (1) và (2) thoe định lí Tha-lès ta có  $MG\parallel CN$.\\
		Mà $CN\subset (ACD)$ nên  $MG\parallel (ACD)$. }
		{\begin{tikzpicture}[font=\footnotesize,scale=.7]
		 \def\a{4}
		\path (0:0) coordinate (B) 
			 ++(0:\a) coordinate (D)
			 ++(-135:.6*\a) coordinate (C)
			($(B)+(62:\a)$) coordinate (A) 
			($(B)!2/3!(C)$)  coordinate (M)
			($(D)!.5!(A)$)  coordinate (N)
			($(B)!2/3!(N)$) coordinate (G); 
		 \draw[dashed] (B)--(D)
			 (B)--(N)
			(G)--(M); 
		\draw (B)--(A)--(D) 
			 (A)--(C)--(N)
			 (B)--(C)--(D);
		 \foreach \x/\g in {A/90,B/180,C/-45,D/0,M/-135,N/45,G/120}
		\fill[black] (\x) circle  (1pt)
			($(\g:3mm)+(\x)$) node {$\x$}; 
	\end{tikzpicture}
		} 
	} 
\end{ex}

% nguồn: Ok_TracNghiemLop11_Lan1
\begin{ex}%[1H4H3-2]%[ma:a54ea94f]
	Cho  tứ diện $ABCD$. Gọi $G_1$ và $G_2$ lần lượt là trọng tâm các tam giác $BCD$ và $ACD$. Chọn khẳng định {\bf sai}?
	\choice 
		{$G_1  G_2 \parallel \left(ABD\right)$}
		{$G_1 G_2  \parallel \left(ABC\right)$}
		{$BG_1$, $AG_2$ và $CD$ đồng qui} 
		 {\True $G_1 G_2=\dfrac{2}{3} AB$}
	 \loigiai{
	\immini{Ta có $G_1$ và $G_2$ lần  lượt là trọng tâm các $\triangle BCD$ và $\triangle ACD$.\\
		Suy ra $BG_1$, $AG_2$ và $CD$ đồng qui tại $M$ với  $M$ là trung điểm $CD$.\\
		Vì $G_1 G_2 \parallel AB$ nên  $G_1 G_2 \parallel \left(ABD\right)$ và $G_1 G_2 \parallel \left(ABC\right)$.\\
		Lại có $\dfrac{G_1 G_2}{AB}=\dfrac{MG_1}{MB}=\dfrac{1}{3}$ $\Rightarrow$ $G_1 G_2=\dfrac{1}{3} AB$.} 
		{\begin{tikzpicture}[font=\footnotesize,scale=.7]
		\def\a{4} 
		\path  (0:0) coordinate (B)
			 ++(0:\a) coordinate (D)
			++(-145:.68*\a)  coordinate (C)
			 ($(B)+(70:.8*\a)$) coordinate (A)
			($(C)!.5!(D)$) coordinate  (M)
			($(M)!1/3!(B)$) coordinate (G_1) 
			 ($(M)!1/3!(A)$) coordinate (G_2);
		 \draw[dashed] (B)--(D)
			 (M)--(B)
			 (G_1)--(G_2);
		\draw (B)--(A)--(D) 
			(A)--(C) 
			 (B)--(C)--(D)
			 (A)--(M);
		\foreach \x/\g in {A/90,B/180,C/-45,D/0,M/-45,G_1/-45,G_2/5} 
		\fill[black] (\x) circle  (1pt)
			($(\g:4mm)+(\x)$)  node {$\x$};
	\end{tikzpicture}
		 }
	 }
\end{ex}

% nguồn: Ok_TracNghiemLop11_Lan1
\begin{ex}%[1H4H3-2]%[ma:ce502571]
	Cho tứ diện $ABCD$. Gọi $M$ là trọng tâm của $\triangle ABC$ và $N$ là điểm nằm  trên cạnh $AD$ sao cho $AN=2ND$. Khi đó ta có
	 \choice
		 {$MN$ cắt $BD$}
		{\True $MN\parallel (BCD)$} 
		 {$MN\parallel CD$}
		 {$AC$ cắt $BD$}
	\loigiai{ 
	\immini{Gọi $E$  là trung điểm $BC$.\\
		Trong $\triangle AED$, có $\dfrac{AM}{AE}=\dfrac{AN}{AD}=\dfrac{2}{3} \\ 
		 \Rightarrow MN\parallel ED\Rightarrow MN\parallel (BCD)$. }
		{\begin{tikzpicture}[font=\footnotesize,scale=.7]
		 \def\a{4}
		\path  (0:0) coordinate (B)
			 ++(0:\a) coordinate (D)
			++(-120:\a/2) coordinate (C) 
			($(B)+(70:\a)$)  coordinate (A)
			($(B)!.5!(C)$)  coordinate (E)
			 ($(A)!2/3!(E)$) coordinate (M)
			($(A)!2/3!(D)$) coordinate  (N);
		\draw[dashed] (B)--(D) 
			 (M)--(N)
			 (E)--(D);
		\draw (B)--(A)--(D) 
			 (A)--(C)
			(B)--(C)--(D) 
			 (A)--(E);
		\foreach \x/\g in  {A/90,B/180,C/-45,D/0,E/-90,M/150,N/45}
		\fill[black] (\x) circle (1pt) 
			($(\g:3mm)+(\x)$)  node {$\x$};
	\end{tikzpicture}
		 }
	} 
\end{ex}

% nguồn: Ok_TracNghiemLop11_Lan1
\begin{ex}%[1H4H3-2]%[ma:2403f9f4]
	Cho tứ diện $ABCD$. Gọi $G$ là trọng tâm của tam giác $ABD$, $Q$ thuộc cạnh $AB$ sao cho $AQ=2QB$, $P$ là trung điểm của  $CB$. Khẳng định nào sau đây đúng?
	 \choice
		{$PQ\parallel (BCD)$} 
		{\True  $GQ\parallel (BCD)$}
		{$PQ\parallel  (ACD)$}
		 {$Q\in (GDP)$}
	\loigiai{ 
	\immini{Ta có  $\dfrac{AG}{AM}=\dfrac{AQ}{AB}=\dfrac{2}{3} \Rightarrow GQ\parallel MB\subset (BCD)$ $\Rightarrow GQ\parallel (BCD)$.}
		{\begin{tikzpicture}[scale=.5,font=\footnotesize]
		\def\a{5} 
		\path (0:0) coordinate (D) 
			++(0:\a)  coordinate (A)
			++(-145:.8*\a)  coordinate (B)
			($(D)+(65:.9*\a)$) coordinate (C) 
			 ($(D)!.5!(B)$) coordinate (M)
			($(B)!.5!(A)$) coordinate (N) 
			($(A)!2/3!(B)$)  coordinate (Q)
			(intersection of D--N and A--M)  coordinate (G);
		\draw[dashed] (A)--(D)--(N) 
			(A)--(M) 
			(G)--(Q); 
		\draw  (D)--(C)--(A)
			 (B)--(C)
			 (A)--(B)--(D);
		\foreach \x/\g  in {A/0,B/-45,C/90,D/180,Q/-90,M/-135,N/-45,G/90}
		\fill[black] (\x) circle (1pt) 
			($(\g:4mm)+(\x)$)  node {$\x$};
	\end{tikzpicture}
		 }
	} 
\end{ex}

% nguồn: Ok_TracNghiemLop11_Lan1
\begin{ex}%[1H4H3-2]%[ma:172af32a]
	\immini[thm]{Cho hình tứ diện $ABCD$. Gọi $M$, $N$, $P$, $Q$, $I$ lần lượt là trung điểm của $BD$, $BC$,  $AC$, $AD$, $CD$. Khẳng định nào sau đây đúng?
	 \choice[2]
		 {$NI\parallel(BCD)$}
		 {\True $MN\parallel(ACD)$}
		 {$MN\parallel(BCD)$}
		 {$PQ\parallel(ACD)$}}
		{ 
		\begin{tikzpicture}[scale=1,font=\footnotesize,line join=round,line cap=round,>=stealth]
				\def\a{4} 
				\path  (0:0) coordinate (B)
					++(0:\a)  coordinate (D)
					++(-150:0.7*\a) coordinate  (C)
					($(B)+(80:\a)$) coordinate (A) 
					($(B)!0.5!(D)$) coordinate  (M)
					($(B)!0.5!(C)$)  coordinate (N)
					($(A)!0.5!(C)$) coordinate (P) 
					 ($(A)!0.5!(D)$) coordinate (Q)
					($(C)!0.5!(D)$)  coordinate (I);
				\draw[dashed] (B)--(D); 
				 \draw (B)--(A)--(D)
					(A)--(C) 
					(B)--(C)--(D); 
				\foreach \x/\g  in {A/90,B/180,C/-45,D/0,M/90,N/-120,P/45,Q/45,I/-45}
				\fill  (\x) circle (1pt)
					($(\g:3mm)+(\x)$)  node {$\x$};
			\end{tikzpicture}
		 }
	 \loigiai{
		Xét $\triangle  BCD$ có:\\
		$M$ là trung điểm  $BD$, $N$ là trung điểm $BC$.\\
		$\Rightarrow MN$ là đường trung bình  $\triangle BCD$.\\
		Khi đó, ta có  $\heva{&MN \parallel CD\\&CD \subset (ACD).}$\\
		 Vậy $CD \parallel (ACD)$.
	 }
\end{ex}

% nguồn: Ok_Toan11_TracNghiem_Lan2
\begin{ex}%[1H4V3-2]%[ma:f392aa04]
	Cho hình chóp $S.ABCD$ có  đáy $ABCD$ là hình bình hành tâm $O$. Gọi $M$ là trung điểm của $SB$. Đường thẳng $DM$ cắt mặt phẳng $(SAC)$ tại $N$. Mặt phẳng $(CDM)$ cắt $S A$ tại $K$. Khẳng định nào sau đây \textbf{sai}?
	\choice 
		{Ba  điểm $S$, $N$, $O$ thẳng hàng}
		 {Ba điểm $C$, $N$, $K$ thẳng hàng}
		{$KM\parallel  CD$}
		{\True $N$ là trung điểm  của đoạn thẳng $CK$}
	\loigiai{ 
		\begin{center}
			\begin{tikzpicture}[scale=1,font=\footnotesize,line join=round,line cap=round,>=stealth]
					 \coordinate (A) at (0,0);
					\coordinate (B) at  (4,0);
					\coordinate (C) at  (2,-2);
					\coordinate (D) at  (-2,-2);
					\coordinate (O) at (intersection of A--C  and B--D);
					\coordinate (S) at (0.5,3); 
					 \coordinate (M) at ($(S)!0.5!(B)$);
					\coordinate (K) at ($(S)!0.5!(A)$); 
					\coordinate (N) at (intersection  of D--M and S--O);
					\draw[dashed]  (A)--(B)
						 (A)--(C)
						 (B)--(D)
						 (D)--(A)
						 (A)--(B)
						(S)--(A) 
						 (S)--(O)
						(K)--(M) 
						(D)--(M); 
					\draw  (S)--(B)
						 (S)--(C)
						(S)--(D) 
						(D)--(C)--(B); 
					\fill (A) circle (1pt)  node[left] {$A$}
						 (B) circle (1pt) node[right] {$B$}
						(C) circle (1pt) node[below  right] {$C$}
						(D) circle (1pt) node[below  left] {$D$}
						(S) circle (1pt) node[above] {$S$} 
						(O) circle (1pt)  node[below] {$O$}
						(M)  circle (1pt) node[right] {$M$}
						(K)  circle (1pt) node[left] {$K$}
						(N) circle (1pt) node[above left=0.8pt and  0pt] {$N$};
				\end{tikzpicture}
		\end{center}
		Trong mặt phẳng  $(SBD)$, gọi $N=DM \cap SO$ mà $SO\subset(SAC)$ nên $DM \cap(SAC)=N$.\\
		Vậy ba điểm $S, N, O$ thẳng hàng.\\ 
		 Trong mặt phẳng $(SAC)$, $CN\cap SA=K$. Mà $CN\subset(CDM)$ nên $SA\cap(CDM)=K$.\\
		Vậy ba điểm $C$, $N$, $K$ thẳng hàng.\\ 
		Ta có $N$ là trọng tâm $\triangle SBD$ nên  $\dfrac{SN}{SO}=\dfrac{2}{3}$.\\
		Mà $\triangle SAC$ có đường trung tuyến $SO$  và $\dfrac{SN}{SO}=\dfrac{2}{3}$ nên $N$ là trọng tâm $\triangle SAC$.\\
		 Vậy $K$ là trung điểm của $SA$ do đó $KM\parallel AB\parallel CD$.\\
		Vì $N$ là trọng tâm $\triangle  SAC$ nên $\dfrac{CN}{CK}=\dfrac{2}{3}$.\\
		Do đó $N$ không phải là trung điểm  của đoạn thẳng $CK$.\\
		Ngoài ra, ta có thể lập luận giao tuyến của $(CDM)$ và $(SAB)$  là đường thẳng đi qua $M$ và song song với $AB$, $CD$, đường thẳng này cắt $SA$ tại $K$.\\
		Suy ra $KM\parallel AB\parallel  CD$ và $K$ là trung điểm của $SA$.
	 }
\end{ex}

% nguồn: Ok_Toan11_TracNghiem_Lan2
\begin{ex}%[1H4V3-2]%[ma:fb51c4fa]
	Cho hình chóp $S.ABCD$, có đáy $ABCD$ là hình bình hành tâm $O$. Gọi $H$, $K$ lần lượt là  trung điếm của $BC$, $CD$. Mệnh đề nào đưới đây \textbf{sai}?
	\choice 
		{$HK \parallel (SBD)$} 
		{$OK \parallel  (SAD)$}
		{$OH  \parallel (SAB)$}
		{\True $HK \parallel (SAB)$} 
	\loigiai{ 
	\immini{Ta  có $HK \not \subset(SBD)$.\\
		Ta thấy $HK$ là đường trung bình của $\triangle BCD$ nên $HK \parallel BD$ mà $BD  \subset(SBD)$.\\
		Do đó  $HK \parallel (SBD)$.\\
		Mà $OK \not  \subset(SAD)$.\\
		Ta thấy $OK$ là đường trung bình của $\triangle  ACD$ nên $OK \parallel AD$ mà $AD \subset(SAD)$.\\
		Do đó  $OK \parallel (SAD)$.\\
		Lại có $OH \not  \subset(SAB)$.\\
		Ta thấy $OH$ là đường trung  bình của $\triangle ABC$ nên $OH \parallel AB$ mà $AB \subset (SAB)$.\\
		 Do đó $OH \parallel (SAB)$.\\
		Trong mặt phẳng $(ABCD)$ ta thấy $AB \cap HK$ mà $AB \subset(SAB)$ nên $HK$  không sông song với ($SAB)$.}{
		\begin{tikzpicture}[scale=0.8,font=\footnotesize,line join=round,line cap=round,>=stealth]
				\coordinate (A) at  (0,0);
				\coordinate (B) at (-2,-2); 
				\coordinate (D)  at (4,0);
				 \coordinate (C) at ($(B)-(A)+(D)$);
				\coordinate (S) at  ($(A)+(100:4)$);
				\coordinate (H)  at ($(C)!.5!(B)$);
				\coordinate (K)  at ($(C)!.5!(D)$);
				 \coordinate (O) at ($(C)!.5!(A)$);
				\draw  (S)--(B)--(C)--(D)--(S)--(C);
				\draw [dashed] (S)--(A)--(D) 
					(B)--(A) 
					(O)--(H)--(K)--(O) 
					(A)--(C) 
					 (B)--(D);
				\foreach \x/\g in {A/180,B/-90,C/-90,S/90,D/0,H/-90,K/-20,O/60} \fill  (\x) circle (1pt) +(\g:3mm) node{$\x$};
			\end{tikzpicture}
		} 
	} 
\end{ex}

% nguồn: Ok_Toan11_TracNghiem_Lan2
\begin{ex}%[1H4V3-2]%[ma:853c575d]
	Cho  hình bình hành $ABCD$. Vẽ các tia $Ax$, $By$, $Cz$, $Dt$ song song, cùng hướng nhau và không nằm trong $( ABCD)$. Mặt phẳng $( \alpha)$ song song với $AB$, và cắt $Ax$, $By$, $Cz$, $Dt$ lần lượt tại các điểm $A'$, $B'$, $C'$, $D'$. Biết $O$ là tâm hình bình hành $ABCD$, $O'$ là giao điểm của $A'C'$ và $B'D'$. Khẳng định nào sau đây \textbf{sai}?
	\choice 
		 {$A'B'C'D'$ là hình bình hành}
		{$( AA'B'B)\parallel C'D'$} 
		 {\True $AA'=CC'$ và $BB'=DD'$}
		{ $OO'\parallel  AA'$}
	\loigiai{ 
		\begin{center}
			\begin{tikzpicture}[scale=1,line cap=round,line join=round,font=\footnotesize,>=stealth]
					\def\dai{4}  \def\rong{2.5} \def\goc{50}
					\path (0,0)  coordinate (B)++(0:\dai) coordinate (C)
						 ++(\goc:\rong) coordinate (D)++(180:\dai) coordinate (A)
						($(B)+(90:3.5)$) coordinate  (B')++(90:1.5) coordinate (y)
						($(A)+(90:3.5)$) coordinate (A')++(90:1.5) coordinate  (x)
						($(C)+(90:4.5)$) coordinate (C')++(90:1.0) coordinate (z) 
						($(D)+(90:4.5)$) coordinate (D')++(90:1.0) coordinate (t); 
					\coordinate (O) at  ($(A)!1/2!(C)$);
					 \coordinate (O') at ($(A')!1/2!(C')$);
					 \draw (B)--(B')--(y) node[left]{$y$}
						(C)--(C')--(z) node[below  right]{$z$}
						(D)--(D')--(t)  node[left]{$t$}
						 (B)--(C)--(D)
						(B')--(C')--(D')--(B') 
						 (x) node[left]{$x$}--(A')--(C')
						 (B')--(A')--(D');
					\draw  [dashed]
						(B)--(A)--(D)--(B) 
						 (A)--(C)
						 (A)--(A');
					\foreach  \t/\g in {A/160,B/-130,C/-60,D/0,B'/180,A'/160,D'/40,C'/-20,O'/90,O/-90}
					 \draw[fill=black] (\t) circle (1.0pt)
						 node[shift={(\g:9pt)}]{$\t$};
				\end{tikzpicture}
		\end{center}
		Ta có $\heva{& \alpha \parallel AB \\ & \alpha  \cap ( ABB'A')=A'B' } \Rightarrow A'B'\parallel AB\parallel CD$.\\
		$\Rightarrow \heva{& A'B'\parallel CD \\ &  \alpha \cap ( DD'C'C)=C'D' } \Rightarrow C'D'\parallel A'B'\Rightarrow C'D'\parallel ( AA'B'B)$.\\
		 Dễ thấy $C'D'\parallel A'B'\parallel AB\parallel CD$ theo câu $A$ mà $AA'\parallel BB'\parallel CC'\parallel DD'$.\\
		$ \Rightarrow AA'B'B$, $CC'D'D$, $ABCD$ là các hình  bình hành.\\
		$\Rightarrow A'B'\parallel C'D'$, $A'B'=C'D'$ suy  ra $A'B'C'D'$ là hình bình hành. \\
		Mặt khác: $O,O'$ lần lượt là trung điểm của $AC$, $A'C'$ nên $OO'$ là đường trung bình trong hình thang $AA'C'C$  do đó $OO'\parallel AA'$.
	} 
\end{ex}

% nguồn: Ok_Toan11_TracNghiem_Lan2
\begin{ex}%[1H4V3-2]%[ma:2dc7af2d]
	Cho lăng  trụ $ABC.A'B'C'$. Gọi $M$, $N$ lần lượt là trung điểm $AA'$ và $B'C'$. Khi đó đường thẳng $AB'$ song song với mặt phẳng
	\choice 
		{$(A'MN)$} 
		{$(C'MN)$} 
		 {\True $(A'CN)$}
		{$(CMN)$} 
	 \loigiai{
	\immini{Gọi $H$,  $K$ lần lượt là trung điếm của $A'B'$, $ A'C$.\\
		Ta có $HM$ là đường trung bình $\triangle A'B'A \Rightarrow  HM \parallel A B'$ \quad $(1)$.\\
		Lại có $HN$, $MK$ lần lượt  là đường trung bình $\triangle A'B'C'$, $\triangle A'AC$.\\
		$\heva{&HN \parallel A'C', HN=\dfrac{1}{2}A'C'\\&MK \parallel AC,  MK=\dfrac{1}{2}AC}$ mà $\heva{&A'C' \parallel AC \\&A'C'=AC}$ nên $\heva{&HN \parallel MK\\&HN=MK.}$\\ $\Rightarrow HNKM$ là hình hình hành.\\
		$\Rightarrow HM \parallel NK$ \quad $(2)$.\\ 
		Từ $(1)$ và $(2)$ suy ra $AB' \parallel NK \Rightarrow AB'\parallel (A'NC)$.}{\begin{tikzpicture}[line cap=round,line join=round,scale=1.0]
		 \def\h{3}
		\path 
			(0,0) coordinate (A') 
			 (3.5,0) coordinate (B')
			(-50:2)  coordinate (C')
			($(A')+(90:\h)$)  coordinate (A)
			 ($(B')+(90:\h)$) coordinate (B)
			 ($(C')+(90:\h)$) coordinate (C)
			($(A)!0.5!(A')$) coordinate  (M)
			($(B')!0.5!(C')$)  coordinate (N)
			 ($(B')!0.5!(A')$) coordinate (H)
			($(A')!0.5!(C)$) coordinate  (K);
		\draw[dashed] (A')--(B')--(A) 
			 (K)--(N)--(H)--(M);
		\draw 
			(A)--(B)--(C)--cycle 
			(A)--(A')--(C')--(B')--(B) 
			 (C)--(C')
			 (A)--(C')--(B)
			 (M)--(K);
		\foreach \x/\g in 
			 {A/180,B/0,C/-20,A'/180,B'/0,C'/-40,M/180,N/-30,H/90,K/-160}
		\fill(\x) circle (1pt) 
			($(\x)+(\g:3mm)$) node{$\x$}; 
	\end{tikzpicture}
		 }
	 }
\end{ex}

% nguồn: Ok_TracNghiemLop11_Lan1
\begin{ex}%[1H4V3-2]%[ma:ea916569]
	Cho lăng trụ đứng $ ABC.A'B'C'$. Gọi $M$, $N$ lần lượt là trung điểm của $ A'B'$ và $ CC'$. Khi đó $ CB'$  song song với
	\choice 
		{$AM$} 
		{$A'N$} 
		{\True  $\left(AC'M\right)$}
		 {$\left(BC'M\right)$}
	 \loigiai{
		\begin{center}
			\begin{tikzpicture}[line join = round,line cap = round,font=\footnotesize,>=stealth,scale=1]
					 \coordinate (B') at (0,0);
					\coordinate (A') at  (-2,1);
					\coordinate (C')  at (4,1);
					\coordinate  (B) at (0,4);
					\coordinate (A)  at (-2,5);
					\coordinate  (C) at (4,5);
					 \coordinate (M) at ($(A')!0.5!(B')$);
					\coordinate (N) at  ($(C)!0.5!(C')$);
					 \coordinate (I) at (intersection of A--C' and A'--C);
					\draw  (A')--(B')--(C')--(C)
						(A)--(B)--(C)--(A)--(A') 
						(B)--(B'); 
					\draw[dashed]  (A)--(C')--(A')--(C)
						 (I)--(M);
					\foreach \i/\g in {A'/-90,B'/-90,C'/-90,A/90,B/90,C/90,M/-90,N/0,I/90}{\draw[fill](\i) circle  (1pt)
						 ($(\i)+(\g:3mm)$) node[scale=1]{$\i$};}
				\end{tikzpicture}
		\end{center}
		Gọi $ I$ là trung điểm của $ A'C$. Ta có $ MI\parallel B'C$ và $ MI\subset\left(AC'M\right)$. Do đó $  CB'\parallel\left(AC'M\right)$.
	} 
\end{ex}


%%%%%%%%%%%%%%%% ĐÚNG/SAI %%%%%%%%%%%%%%%%

% nguồn: Ok_dungsai_lop11_Lan1
\begin{ex}%[1H4N3-2]%[ma:773930ae]
 Cho hình chóp $S.ABCD$, đáy $ABCD$ là hình bình hành có $O$ là giao  điểm của hai đường chéo. Cho $M$ là trung điểm của $SC$. Xét tính đúng sai của các khẳng định sau
 \choiceTF 
  {Đường thẳng $OM$ là giao tuyến  của hai mặt phẳng $(SAC)$ và $(SBD)$}
  {\True Đường thẳng $BC$ song song  với mặt phẳng $(AMD)$}
  {\True  Đường thẳng $OM$ song song với mặt phẳng $(SAD)$}
  {Giao tuyến của hai mặt phẳng $(OMD)$ và $(SAD)$ là đường thẳng qua $D$ và song  song với $SA$}
  \loigiai{
  \begin{center}
   \begin{tikzpicture}[line join=round,line cap=round,>=stealth,scale=0.9,font=\footnotesize]
     \tikzset{every  node/.style={outer sep=-0.5mm}}
     \foreach \x/\y/\n in {0/0/A,5/0/D} \coordinate (\n) at (\x,\y); 
     \coordinate (B) at  ($(A)+(-2,-2)$);
      \coordinate (C) at ($(D)+(-2,-2)$);
     \coordinate (O)  at ($(A)!0.5!(C)$);
     \coordinate (S) at  ($(O)+(0,6)$);
     \coordinate (M) at ($(S)!0.5!(C)$); 
     \draw (S)--(B)--(C)--(D)--(S) 
       (C)--(S)--(B)
       (D)--(S)
       (M)--(D);
     \draw[dashed]  (B)--(A)--(D)--(B)
       (C)--(O)--(A)--(S)
      (M)--(O)--(S) 
      (A)--(M); 
     \foreach \t/\g in {A/180,B/180,C/0,D/0,S/90,O/-90,M/0}{ 
      \fill (\t) circle  (1pt) node[shift={(\g:10pt)}]{$ \t $};}
    \end{tikzpicture}
  \end{center}
  \begin{itemchoice}
   \itemch \textbf{Sai}. Ta có $\heva{&O \in AC\\&O \in BD} \Rightarrow  \heva{&O \in (SAC)\\&O \in (SBD)} \Rightarrow O \in (SAC) \cap (SBD) \quad(1)$.\\
   Hơn nữa dễ  thấy $\heva{&S \in (SAC)\\&S \in (SBD)} \Rightarrow S \in (SAC) \cap (SBD) \quad(2)$.\\
   Từ $(1)$ và $(2)$ ta được $\heva{&O \in (SAC) \cap (SBD)\\&S \in  (SAC) \cap (SBD)}$ nên $(SAC) \cap (SBD) = OS \neq OM$.
   \itemch \textbf{Đúng}. Ta có đáy $ABCD$ là hình bình hành nên  $BC \parallel AD$.\\
   Do $BC$ không nằm trong mặt phẳng $(AMD)$ và $BC \parallel AD$, với $AD$ nằm trong $(AMD)$, suy ra $BC  \parallel (AMD)$.
   \itemch \textbf{Đúng}. Ta có $OM$ là đường trung bình của  tam giác $SAC$, suy ra được $OM \parallel SA$.\\
   $OM$ không nằm trong mặt phẳng $(SAD)$ và $OM \parallel SA$ với $SA$ nằm trong $(SAD)$,  suy ra $OM \parallel (SAD)$.
   \itemch \textbf{Sai}. Ta có hai mặt phẳng $(OMD)$ và $(SAD)$ có điểm chung $D$ và lần lượt đi  qua hai đường thẳng song song $OM$ và $SA$.\\
   Suy ra  giao tuyến của hai mặt phẳng $(OMD)$ và $(SAD)$ là đường thẳng đi qua $D$ và song song với $SA$ và $OM$.
  \end{itemchoice}
  }
\end{ex}

% nguồn: Ok_dungsai_lop11_Lan1
\begin{ex}%[1H4N3-2]%[ma:af71b459]
 Cho hai hình bình hành $ABCD$ và $ABEF$ không  cùng nằm trong một mặt phẳng. Gọi $O$, $O'$ lần lượt là tâm của hình bình hành $ABCD$, $ABEF$ và $M$ là trung điểm của $CD$.
 \choiceTF 
  {\True $OO'$ là đường trung  bình của tam giác $ACE$}
  {\True  $OO'\parallel CE$}
  {\True $OO' \parallel (BEC)$} 
  {Số giao điểm  của đường thẳng $OO'$ và mặt phẳng $(CFM)$ là $1$}
 \loigiai{ 
  \begin{center}
   \begin{tikzpicture}[scale=1,font=\footnotesize,line join=round,line cap=round,>=stealth]
     \coordinate (A)  at (0,0);
     \coordinate (B) at  (4,0);
     \coordinate (D)  at (1.7,1.7);
     \coordinate (C)  at ($(B)+(D)-(A)$);
     \coordinate (F) at (1.3,-1.7); 
     \coordinate (E) at ($(B)+(F)-(A)$); 
      \coordinate (O) at ($(A)!1/2!(C)$);
     \coordinate (O')  at ($(A)!1/2!(E)$);
     \coordinate (M)  at ($(D)!1/2!(C)$);
      \draw (A)--(D)--(C)--(E)--(F)--(A)
      (D)--(F); 
      \draw[dashed] (A)--(E)--(B)--(C)
       (A)--(B)--(F)
      (A)--(C) 
       (B)--(D)
       (O)--(O');
     \foreach \x/\g in{A/180,B/0,C/45,D/135,M/90,F/-135,E/-45,O/90,O'/-90}  \fill (\x) circle (1pt)+(\g:.35) node{$\x$};
    \end{tikzpicture}
  \end{center}
  \begin{itemchoice}
   \itemch \textbf{Đúng.}\\ 
   Vì $O$, $O'$ lần lượt là tâm của hình bình hành $ABCD$, $ABEF$ nên $O$, $O'$ lần lượt là trung điểm của $AC$  và $AE$ nên $OO'$ là đường trung bình của tam giác $ACE$.
    \itemch \textbf{Đúng.}\\
   Vì $OO'$ là  đường trung bình của tam giác $ACE$ nên $OO'\parallel CE$.
   \itemch \textbf{Đúng.}\\ 
   Vì  $OO'\parallel CE \subset (BCE)\Rightarrow OO' \parallel (BEC)$.
   \itemch \textbf{Sai.}\\ 
   Vì $OO'\parallel CE \subset (CDFE)\equiv(CFM)\Rightarrow OO' \parallel (CFM)$. 
  \end{itemchoice}
  }
\end{ex}

% nguồn: Ok_dungsai_lop11_Lan1
\begin{ex}%[1H4N3-2]%[ma:152aafd8]
 \immini{Cho tứ diện $ABCD$. gọi $M$, $N$, $P$ lần lượt là trung  điểm của $AB$, $AC$, $AD$. Xét tính đúng sai của các khẳng định sau
  \choiceTF
  {Đường thẳng $MN$ song  song với mặt phẳng $\left(ABC\right)$}
  {\True Đường thẳng $MN$ song song với  mặt phẳng $\left(BCD\right)$}
  {\True Đường thẳng $BC$ song song với mặt  phẳng $\left(MNP\right)$ }
  {\True  Đường thẳng $BD$ song song với mặt phẳng $\left(MNP\right)$ }
  }{ 
  \begin{tikzpicture}[scale=0.9,line join=round,line cap=round]
    \coordinate (A) at (1,3); 
    \coordinate (B)  at (0,0);
    \coordinate (C) at  (2,-2);
    \coordinate  (D) at (4.5,0);
    \coordinate (M) at  ($(A)!0.5!(B)$);
    \coordinate  (N) at ($(A)!0.5!(C)$);
    \coordinate (P) at ($(A)!0.5!(D)$); 
    \draw (A) node[above]{$A$}--(B) node[left]{$B$}--(C) node[below]{$C$}--(D) node[right]{$D$}--(A)--(C); 
     \draw[dashed] (B)--(D)
      (M)--(P);
    \draw (M) node[left]{$M$}--(N) node[below right]{$N$}--(P)  node[above]{$P$};
    \foreach \point in  {A,B,C,D,M,N,P}
    \fill  [black] (\point) circle (1pt);
   \end{tikzpicture}
  } 
 \loigiai{ 
  \begin{itemchoice}
   \itemch \textbf{Sai.}\\ 
   Ta có $MN \subset  \left(ABC\right)$ nên $MN$ không song song với $\left(ABC\right)$
   \itemch  \textbf{Đúng.}\\
   Ta có $MN \parallel BC$ mà $BC \subset \left(BCD\right)$  nên $MN \parallel \left(BCD\right)$
   \itemch \textbf{Đúng.}\\ 
   Ta có $BC  \parallel MN$ mà $MN \subset \left(MNP\right)$ nên $BC \parallel \left(MNP\right)$
   \itemch \textbf{Đúng.}\\ 
   Ta có $BD \parallel MP$ mà $MP \subset \left(MNP\right)$  nên $BD \parallel \left(MNP\right)$
  \end{itemchoice}
 } 
\end{ex}

% nguồn: Ok_dungsai_lop11_Lan1
\begin{ex}%[1H4H3-2]%[ma:3aef4780]
 Cho tứ diện $ABCD$. Gọi $G$, $K$ là trọng tâm tam giác $BCD$, $ABD$. $N$  là điểm thuộc cạnh $AD$ sao cho $ND=2 NA$. Xét tính đúng sai các phát biểu sau
  \choiceTF
  {\True $NG \parallel (ABC)$} 
  {$(NGK)\parallel(ACD)$} 
   {\True $GK\parallel(ABC)$}
  {\True  $(KNG)\parallel(ABC)$}
 \loigiai{\begin{center}
  \begin{tikzpicture}[scale=1,font=\footnotesize,line join=round,line cap=round,>=stealth]
     \def\cnhat{5}
    \def\chai{2} 
     \def\cba{4}
     \def\goc{50}
    \coordinate[label=left:$B $] (B) at (0,0); 
    \coordinate[label=right:$D $] (D) at (\cnhat,0); 
    \coordinate[label=below left:$C$] (C) at  (-\goc:\chai);
    \coordinate[label=above:$A$] (A)  at (70:\cba);
    \path 
     ($(A)!0.33!(D)$) coordinate[label=right:$N  $] (N)
     ($(C)!0.5!(B)$) coordinate[label=left:$I  $] (I)
      ($(D)!0.66!(I)$) coordinate[label=below:$G $] (G)
     ($(A)!0.5!(B)$) coordinate[label=left:$J $] (J) 
     ($(D)!0.66!(J)$) coordinate[label=left:$K$]  (K);
     \draw (A)--(I)
     (B)--(C)--(D)--(A)--cycle (A)--(C); 
     \draw[dashed] (B)--(D)--(I)
     (N)--(G)--(K)--(N); 
    \foreach \diem  in {B,C,D,A,N,G,I,K,J}\fill (\diem) circle (1pt);
   \end{tikzpicture}
 \end{center}
 \begin{itemchoice}
   \itemch Gọi $I$ là trung điểm của $BC$.\\
  Tam giác $ADI$ có $\dfrac{DN}{DA}=\dfrac{DG}{DI}$ nên theo định  lí Talet, ta có $NG \parallel AI$.\\
  Vì $\heva{&NG \not\subset (ABC)\\&NG \parallel AI\\&AI\subset (ABC)}$  nên $NG\parallel(ABC)$.
  \itemch  Ta có $(NGK)\cap (ACD)=N$ nên $(NGK)$ không song song $(ACD)$.
  \itemch Ta có $\dfrac{DG}{DI}=\dfrac{DK}{DJ}=\dfrac{2}{3}$ nên  $GK\parallel IJ$, suy ra $GK\parallel (ABC)$.
  \itemch Ta có  $\heva{&KG\parallel (ABC)\\&NG\parallel (ABC)}\Rightarrow (KNG)\parallel(ABC)$.
 \end{itemchoice}
} 
\end{ex}

% nguồn: Ok_dungsai_lop11_Lan1
\begin{ex}%[1H4H3-2]%[ma:61428b13]
 \immini[thm]{Cho hình chóp $S.ABCD$, như hình vẽ, có đáy $ABCD$ là hình bình hành. Gọi $O$ giao điểm của $AC$ và $BD$; điểm $M$, $N$ lần lượt là trung điểm của $SA$ và $SC$.  Khi đó
  \choiceTF
  {\True Gọi $P$ là giao điểm của $SO$ và $MN$. Khi đó điểm $P$ nằm trên mặt  phẳng $\left(SBD\right)$}
  {\True Đường thẳng $AC$ song song  với $\left(DMN\right)$}
  {Giao tuyến của $(SAD)$ và  $(SBC)$ sẽ song song với đường thẳng $AB$}
  {\True Đường thẳng $MN$ song song với  mặt phẳng $(ABCD)$}
  } 
  { 
  \begin{tikzpicture}[scale=0.7,font=\footnotesize,line join=round,line cap=round,>=stealth]
    \path 
     (0,0) coordinate  (A)
      ++(-140:3) coordinate (B)
     ++(0:6) coordinate  (C)
      ($(A)+(C)-(B)$) coordinate (D)
      ($(A)+(1,3)$) coordinate (S)
     ($(A)!0.5!(C)$) coordinate  (O)
      ($(S)!0.5!(A)$) coordinate (M)
      ($(S)!0.5!(C)$) coordinate (N)
      ($(M)!0.5!(N)$) coordinate (P);
    \foreach \i in{B,C,D}{\draw (S)--(\i);}; 
    \draw (B)--(C)--(D); 
     \draw[dashed] (S)--(A)--(B)
      (A)--(D)
      (A)--(C)
      (B)--(D)
      (S)--(O)
     (M)--(N); 
    \foreach \i/\g in  {A/-90,B/-90,C/-90,D/0,S/90,M/30,N/0,O/-90,P/-150}
    \fill[black] (\i) circle (1pt)+(\g:4mm) node[scale=1]{$\i$}; 
   \end{tikzpicture}
   }
  \loigiai{
  \begin{center}
   \begin{tikzpicture}[scale=0.7,font=\footnotesize,line join=round,line cap=round,>=stealth]
      \path
      (0,0)  coordinate (A)
      ++(-140:3)  coordinate (B)
      ++(0:6)  coordinate (C)
      ($(A)+(C)-(B)$) coordinate  (D)
      ($(A)+(1,3)$)  coordinate (S)
      ($(S)-(1,0)$) coordinate  (Sx)
      ($(S)+(4,0)$) coordinate (Sxx) 
      ($(A)!0.5!(C)$) coordinate  (O)
      ($(S)!0.5!(A)$)  coordinate (M)
      ($(S)!0.5!(C)$)  coordinate (N)
      ($(M)!0.5!(N)$)  coordinate (P);
     \foreach  \i in{B,C,D}{\draw (S)--(\i);};
     \draw (B)--(C)--(D) 
       (Sx)--(Sxx);
     \draw[dashed] (S)--(A)--(B) 
      (A)--(D) 
       (A)--(C)
       (B)--(D)
       (S)--(O)
      (M)--(N); 
     \foreach  \i/\g in {A/-90,B/-90,C/-90,D/0,S/90,M/30,N/0,O/-90,P/-150}
     \fill[black]  (\i) circle (1pt)+(\g:4mm) node[scale=1]{$\i$};
     \fill[black] (Sxx)  circle (0pt) + (90:4mm) node[scale=1]{$x$};
    \end{tikzpicture}
  \end{center}
  \begin{itemchoice}
   \itemch  Đúng. $P=MN\cap SO \Rightarrow P \in SO \subset (SBD) \Rightarrow P \in (SBD)$.
   \itemch Đúng. Ta có $AC \parallel  MN \subset (DMN)$ nên $AC \parallel (DMN)$.
   \itemch Sai. $(SAD) \cap (SBC) = Sx \parallel AD \parallel BC$ nên  $Sx$ không song song với $AB$.
   \itemch  Đúng. $MN$ là đường trung bình của tam giác $SAC$ nên $MN \parallel AC \subset (ABCD)$. Suy ra $MN \parallel (ABCD)$.
  \end{itemchoice}
  }
\end{ex}

% nguồn: Ok_dungsai_lop11_Lan1
\begin{ex}%[1H4H3-2]%[ma:b9beef77]
 Cho hình chóp $S.ABCD$ có đáy là  hình bình hành. Các điểm $I$, $J$ lần lượt là trọng tâm các tam giác $SAB$ và $SAD$. Gọi $M$ là trung điểm của $CD$.
 \choiceTF 
  {\True $IJ\parallel  (ABCD)$}
  {\True  $IJ\parallel (SBD)$}
   {$IJ\parallel (SAC)$}
  {$IM\parallel  (SAB)$}
  \loigiai{
  \begin{center}
   \begin{tikzpicture}[scale=1,font=\footnotesize,line join=round,line cap=round,>=stealth]
     \def\a{6} 
     \path (0:0) coordinate  (A)
      ++(0:\a)  coordinate (D)
       ++(-130:\a/2) coordinate (C)
      ($(A)+(C)-(D)$) coordinate (B) 
      ($(A)+(80:\a)$) coordinate (S) 
      ($(C)!1/2!(D)$) coordinate (M) 
      ($(A)!1/2!(B)$) coordinate  (E)
      ($(A)!1/2!(D)$) coordinate (F) 
      ($(S)!2/3!(E)$) coordinate (I) 
      ($(S)!2/3!(F)$) coordinate (J) 
      (intersection of  E--F and A--C) coordinate (G)
      (intersection of I--J and S--G) coordinate (H); 
     \draw[dashed]  (B)--(A)--(D)
       (A)--(S)
       (F)--(S)--(E)
       (I)--(J)
       (E)--(F)
       (B)--(D)
      (A)--(C) 
      (I)--(M) 
      (S)--(G); 
     \draw (B)--  (C)--(D)
       (B)--(S)
      (C)--(S) 
       (D)--(S);
     \foreach \x/\g in  {A/135,B/-135,C/-45,D/45,S/90,M/0,I/-65,J/160,E/-90,F/145,G/-90,H/-50}
     \fill[black] (\x) circle (1pt) 
       ($(\g:3mm)+(\x)$) node {$\x$};
    \end{tikzpicture}
  \end{center}
  \begin{itemchoice}
   \itemch {\bf Đúng}.\\ 
   Gọi $E$, $F$ lần lượt là trung điểm của $AB$,  $AD$. Ta có $\dfrac{SI}{SE}=\dfrac{SJ}{SF}=\dfrac{2}{3}$ (vì $I$, $J$ lần lượt là trọng tâm của tam giác $SAB$, $SAD$). \\
   Suy  ra $IJ\parallel EF$.\\
   Khi đó  $\heva{&IJ\parallel EF\subset (ABCD)\\&IJ\not\subset (ABCD)}\Rightarrow IJ\parallel (ABCD)$.
   \itemch {\bf Đúng}.\\ 
   Ta có $EF\parallel BD$ (vì $EF$ là đường trung bình tam giác  $ABD$).\\ Mà $IJ\parallel EF$ nên $IJ\parallel BD$.\\ Khi đó $\heva{&IJ\parallel BD\subset (SBD)\\&IJ\not\subset (SBD)}\Rightarrow IJ\parallel (SBD)$.
    \itemch {\bf Sai}.\\
   Trong mặt phẳng $(ABCD)$, gọi $G$ là giao điểm của $EF$ và  $AC$.\\
   Ta có $SG=(SEF)\cap (SAC)$.\\ 
   Trong $(SEF)$, kẻ  $IJ\cap SG=H$. Vậy $IJ\cap (SAC)=H$.
   \itemch{\bf  Sai}.\\
    Ta có $IM\cap (SAB)=I$.
  \end{itemchoice}
 } 
\end{ex}

% nguồn: Ok_dungsai_lop11_Lan1
\begin{ex}%[1H4H3-2]%[ma:5de8a40b]
 Cho hình chóp $S.ABCD$ có đáy $ABCD$ là hình bình hành tâm $O$. Gọi $M$, $N$ lần lượt là trung điểm các  cạnh $AB$ và $CD$, $P$ là trung điểm cạnh $SA$.
 \choiceTF 
  {\True $PO  \subset (MNP)$}
  {\True $MN \parallel  (SBC)$}
  {\True $MN  \parallel (SAD)$}
  {$SB$ cắt mặt phẳng $(MNP)$} 
 \loigiai{ 
  \begin{center}
   \begin{tikzpicture} [>=stealth,line join=round,line cap=round,font=\footnotesize,scale=1]
     \coordinate (A)  at (-1.3,0);
     \coordinate (D)  at (5,0);
     \coordinate (C) at  (3.1,-1.6);
      \coordinate (B) at ($(A)-(D)+(C)$);
     \draw (B)--(C)--(D); 
     \draw[dashed]  (B)--(A)--(D);
     \coordinate  (S) at (-0.3,4.7);
      \draw (B)--(S);
     \draw  (C)--(S);
     \draw  (S)--(D);
     \draw[dashed]  (S)--(A);
     \draw  (S) node[above] {$S$}
      (A) node[above  left=-2pt] {$A$}
       (B) node[below] {$B$}
      (C) node[below=1pt]  {$C$}
       (D) node[above right=-2pt] {$D$};
     \coordinate (N)  at ($(C)!0.5!(D)$);
      \draw (N) node[right] {$N$};
     \coordinate (M) at ($(A)!0.5!(B)$); 
      \draw (M) node [left] {$M$};
     \draw[dashed]  (A)--(C);
      \draw[dashed] (B)--(D);
     \coordinate (O)  at (intersection of A--C and B--D);
     \draw (O) node[above] {$O$}; 
     \coordinate  (P) at ($(S)!0.5!(A)$);
     \draw (P) node  [above right] {$P$};
     \draw[dashed] (M)--(N); 
     \draw[dashed]  (P)--(N);
      \draw[dashed] (M)--(P);
    \end{tikzpicture}
  \end{center}
  \begin{itemchoice}
   \itemch Đúng vì $O \in  MN\subset (MNP)$ nên $PO\subset (MNP)$.
   \itemch Đúng vì $MN$  là đường trung bình của hình bình hành $ABCD$ nên $MN \parallel BC$, mà $BC\subset(SBC) \Rightarrow MN \parallel(SBC)$.
   \itemch Đúng. Tương tự $MN \parallel AD$, $AD \subset(SAD) \Rightarrow MN  \parallel(S A D)$.
   \itemch Ta có $MP$ là đường trung bình của tam giác $SAB$ nên $SB \parallel MP$, mà $M P \subset(MNP)$ nên $SB  \parallel(MNP)$.
  \end{itemchoice}
  }
\end{ex}

% nguồn: Ok_dungsai_lop11_Lan1
\begin{ex}%[1H4H3-2]%[ma:69798065]
 Cho hình chóp $S.ABCD$ có đáy $ABCD$ là hình bình hành. Lấy điểm  $M$ trên cạnh $AD$ sao cho
 $AD=3AM$. Gọi $G$, $N$ theo thứ tự  là trọng tâm các tam giác $SAB$ và $ABC$. Khi đó
 \choiceTF 
  {Giao tuyến của hai mặt phẳng $(SAB)$ và $(SCD)$ là đường thẳng đi qua $S$ và song song với  $AC$, $BD$}
   {$\dfrac{DN}{DB}=\dfrac{1}{3}$}
  {\True $MN$ song song với mặt phẳng  $(SCD)$}
   {\True $NG$ song song với mặt phẳng $(SAC)$}
  \loigiai{
  \begin{center}
   \begin{tikzpicture}[scale=1,line join=round,font=\footnotesize,line cap=round,>=stealth,thick]
     \coordinate  (A) at (0.5,0);
     \coordinate (B) at  (-2,-2);
     \coordinate (D) at  (6,0);
      \coordinate (C) at ($(B)+(D)-(A)$);
     \coordinate (S)  at ($(A)+(1,5)$);
      \path ($(A)!1/3!(D)$) coordinate (M);
      \path ($(A)!0.5!(B)$) coordinate (K);
     \path ($(A)!0.5!(C)$)  coordinate (O);
     \path  ($(S)!2/3!(K)$) coordinate (G);
     \path ($(C)!2/3!(K)$)  coordinate (N);
      \draw (S)--(B)
      (S)--(C) 
      (S)--(D) 
      (B)--(C)--(D); 
      \draw[dashed,thin](S)--(A)
      (A)--(B) 
      (A)--(D) 
       (S)--(K)--(C)--(A)
      (B)--(D) 
       (G)--(N)--(M);
     \foreach \i/\g in {S/90,A/45,B/-90,C/-90,D/0,M/90,K/-90,G/0,N/-90,O/90}{\draw[fill=black](\i)  circle (1pt)
      ($(\i)+(\g:3mm)$)  node[scale=1]{$\i$};}
    \end{tikzpicture}
  \end{center}
  \begin{itemchoice}
    \itemch{\bf Sai.}\\
   Giao tuyến của hai  mặt phẳng $(SAB)$ và $(SCD)$ là đường thẳng đi qua $S$ và song song với $AB$, $CD$.
   \itemch{\bf Sai.}\\ 
    Gọi $O=AC\cap BD$, khi đó $\dfrac{BN}{BO}=\dfrac{2}{3}\Rightarrow \dfrac{BN}{BD}=\dfrac{1}{3}\Rightarrow \dfrac{DN}{DB}=\dfrac{2}{3}$.
    \itemch{\bf Đúng.}\\
   Ta có $\heva{&\dfrac{DN}{DB}=\dfrac{2}{3}\quad\text{(cmt)}\\&\dfrac{DM}{DA}=\dfrac{2}{3}\quad\text{(gt)}}\Rightarrow MN\parallel  AB\parallel CD$.\\
   Mà $CD\subset (SCD)\Rightarrow MN\parallel(SCD)$. 
    \itemch{\bf Đúng.}\\
    Gọi $K$ là trung điểm $AB$.\\
   Ta có $\heva{&\dfrac{KG}{KS}=\dfrac{1}{3}\quad(\text{do G là trọng tâm}\,\,\triangle  SAB)\\&\dfrac{KN}{KC}=\dfrac{1}{3}\quad(\text{do N là trọng tâm}\,\,\triangle ABC)}\Rightarrow NG\parallel SC$.\\
   Mà $SC\subset (SAC)\Rightarrow NG\parallel  (SAC)$.
  \end{itemchoice}
 } 
\end{ex}

% nguồn: Ok_dungsai_lop11_Lan1
\begin{ex}%[1H4H3-2]%[ma:eb967de6]
 Cho lăng trụ tam giác $ ABC.A'B'C'$ có  $ I$, $K$, $G$ lần lượt là trọng tâm các tam giác $ ABC$, $A'B'C'$, $ACC'$. Các mệnh đề sau đúng hay sai?
 \choiceTF 
  {\True$  BB'\parallel \left(ACC'A'\right)$}
  {\True$\left(ABC\right)\parallel  \left(A'B'C'\right)$}
  {$ IG$ cắt $\left(BCC'B'\right)$} 
  {\True$ (IKG)\parallel \left(BCC'B'\right)$} 
 \loigiai{ 
  \begin{center}
   \begin{tikzpicture}[line join = round,line cap = round,font=\footnotesize,>=stealth,scale=1]
      \coordinate (B') at (0,0);
     \coordinate (A') at  (-2,1);
     \coordinate  (C') at (4,1);
     \coordinate (B) at (0,4); 
     \coordinate  (A) at (-2,5);
      \coordinate (C) at (4,5);
     \coordinate  (M) at ($(C)!0.5!(B)$);
     \coordinate (M') at  ($(C')!0.5!(B')$);
     \coordinate (N)  at ($(C)!0.5!(C')$);
     \coordinate (I) at ($(A)!2/3!(M)$); 
     \coordinate (K)  at ($(A')!2/3!(M')$);
     \coordinate (G)  at ($(A)!2/3!(N)$);
     \draw  (A')--(B')--(C')--(C)
       (A)--(B)--(C)--(A)--(A')
       (B)--(B')
      (A)--(M)--(N) 
       (M)--(M');
     \draw[dashed,thin] (A)--(C')--(A') 
      (A)--(N) 
       (I)--(G)
       (A')--(M')
      (I)--(K); 
     \foreach \i/\g in {A'/-90,B'/-90,C'/-90,A/90,B/90,C/90,M/90,N/0,I/90,K/-40,G/-130,M'/-90}{\draw[fill=black](\i)  circle (1pt)
       ($(\i)+(\g:3mm)$) node[scale=1]{$\i$};}
    \end{tikzpicture}
  \end{center}
  \begin{itemchoice}
   \itemch Ta có: $\heva{&BB'\parallel AA'\\&AA'\subset\left(ACC'A'\right)} 
   \Rightarrow  BB'\parallel \left(ACC'A'\right)$.\\
   Vậy mệnh đề $  BB'\parallel \left(ACC'A'\right)$ đúng.
   \itemch Ta có: $\left(ABC\right)\parallel \left(A'B'C'\right)$ (do $\left(ABC\right),\left(A'B'C'\right)$ là hai mặt phẳng chứa hai đáy của lăng trụ nên song song với nhau).\\ 
   Vậy mệnh đề $\left(ABC\right)\parallel \left(A'B'C'\right)$  đúng.
   \itemch Gọi  $ M$, $M'$ lần lượt là trung điểm của $ BC$, $B'C'$.
   Gọi $  N$ là trung điểm của $ CC'$, tam giác $ AMN$ có\\
   $\dfrac{AI}{AM}=\dfrac{AG}{AN}=\dfrac{2}{3}{\rm}$ (tính chất trọng  tâm).\\
   Suy ra $ IG\parallel MN$ mà $ MN\subset\left(BCC'B'\right)$ nên $ IG\parallel \left(BCC'B'\right)$.\hfill  $(1)$\\
   Vậy mệnh đề $ IG$ cắt $\left(BCC'B'\right)$  sai.
   \itemch $ MM'$ là đường  trung bình của hình bình hành $ BCC'B'$ nên\\
    $\heva{&MM'\parallel BB'\\&MM'=BB'}
   \Rightarrow\heva{&MM'\parallel  AA'\\&MM'=AA'}
    \Rightarrow AMM'A'$ là hình bình hành.\\
   Vì $ I$, $K$ theo  thứ tự là trọng tâm các tam giác $ ABC$, $A'B'C'$ nên\\
   $ IM=KM'=\dfrac{1}{3}A'M'=\dfrac{1}{3}AM$, mà $ IM\parallel KM'$ nên $ IKM'M$  là hình bình hành.\\
   Suy ra $ IK\parallel MM',MM'\subset\left(BCC'B'\right)\Rightarrow  IK\parallel \left(BCC'B'\right)$.\hfill $(2)$\\
   Từ $(1)$ và $(2)$ suy ra $ (IKG)\parallel \left(BCC'B'\right)$.  Vậy mệnh đề $ (IKG)\parallel \left(BCC'B'\right)$ đúng.
  \end{itemchoice}
 } 
\end{ex}

% nguồn: Ok_dungsai_lop11_Lan1
\begin{ex}%[1H4H3-2]%[ma:5699dfc1]
 Cho hình chóp $S.ABCD$ có đáy $ABCD$ là hình bình hành. Gọi $M$, $N$, $P$, $Q$ lần lượt là trung điểm của các cạnh bên $SA$, $SB$, $SC$,  $SD$.
 \choiceTF 
  {\True $MN \parallel  AB$}
   {$PQ \parallel CB$}
  {$MN \parallel  (SAD)$}
  {\True $PQ \parallel (SAB)$} 
  \loigiai{
  \begin{center}
   \begin{tikzpicture}[line join=round,line cap=round,thick]
     \coordinate (A) at  (0,0);
     \coordinate (B)  at (2,-2);
     \coordinate (D) at  (5,0);
     \coordinate  (C) at ($(B)+(D)-(A)$);
     \coordinate  (O) at ($(A)!0.5!(C)$);
      \coordinate (S) at ($(O)+(0.2,6)$);
      \foreach \i/\j/\k in {M/S/A,N/S/B,P/S/C,Q/S/D}{
      \coordinate (\i) at  ($(\j)!0.5!(\k)$);
     } 
     \foreach \i/\j in  {A/B,B/C,S/A,S/B,S/C,M/N,N/P}{
      \draw (\i)--(\j); 
     } 
     \foreach \i/\j in  {A/D,D/C,S/D,M/Q,Q/P}{
       \draw[dashed](\i)--(\j);
     } 
     \foreach  \i/\g in {S/90,A/180,B/-90,C/-90,D/0,M/160,N/-45,P/0,Q/45}{\draw[fill=black](\i) circle (1pt)
      ($(\i)+(\g:3mm)$) node[scale=1]{$\i$};} 
    \end{tikzpicture}
  \end{center}
  \begin{itemchoice}
    \itemch \textbf{Đúng.}\\
   Xét tam giác $SAB$ có $M$ và $N$ lần lượt là trung điểm của các cạnh $SA$ và $SB$ nên $MN$ là đường trung bình của  tam giác $SAB$, suy ra $MN\parallel AB$.
    \itemch \textbf{Sai.}\\
    Xét tam giác $SCD$ có $P$ và $Q$ lần lượt là trung điểm của các cạnh $SC$ và $SD$ nên $PQ$ là đường trung bình của tam giác $SCD$, suy ra $PQ\parallel CD$.
   \itemch  \textbf{Sai.}\\
    Do $M$ là trung điểm $SA$ nên suy ra $M\in SA$.\\
   Mặt  khác $SA \subset (SAD)$ nên suy ra $M \in (SAB)$.\\
    Suy ra $MN$ cắt $(SAB)$.
    \itemch \textbf{Đúng.}\\
   Do $PQ\parallel CD$ mà $CD \parallel AB \subset (SAB)$ nên suy ra $PQ \parallel (SAB)$. 
  \end{itemchoice}
  }
\end{ex}

% nguồn: Ok_dungsai_lop11_Lan1
\begin{ex}%[1H4H3-2]%[ma:18e5a9cb]
 Cho hình chóp $S.ABCD$ có đáy $ABCD$ là hình bình hành tâm $O$. Gọi $M$, $N$ lần lượt là trung điểm của các cạnh $AB$ và  $CD$, $P$ là trung điểm cạnh $SA$. Khi đó các khẳng định sau đây đúng hay sai?
 \choiceTF 
  {\True $MN\parallel  (SBC)$}
  {\True $MN \parallel (SAD)$} 
  {$SB$  cắt mặt phẳng ($MNP)$}
  {$SC$ cắt  mặt phẳng ($MNP)$}
  \loigiai{
  \begin{center}
   \begin{tikzpicture}[scale=1,line join = round,line cap = round]
      \def\a{4}
      \path (0:0) coordinate (A)
      ++(0:\a) coordinate (D) 
      ++(-130:\a/2) coordinate (C) 
      ($(A)+(C)-(D)$) coordinate (B) 
       ($(A)+(80:\a)$) coordinate (S)
      ($(A)!.5!(B)$)  coordinate (M)
       ($(C)!.5!(D)$) coordinate (N)
      ($(S)!.5!(A)$) coordinate  (P)
       (intersection of A--C and B--D) coordinate (O);
      \draw[dashed] (B)--(A)--(D)
      (A)--(S) 
      (M)--(N)--(P)--cycle 
      (S)--(O) 
      (P)--(O) 
      (A)--(C) 
       (B)--(D);
     \draw (B)--  (C)--(D)
      (B)--(S) 
       (C)--(S)
       (D)--(S);
     \foreach \x/\g in {A/45,B/-135,C/-45,D/45,S/90,M/-70,N/0,P/30,O/-90} 
     \fill[black] (\x) circle  (1pt)
      ($(\g:3.5mm)+(\x)$) node  {$\x$};
    \end{tikzpicture}
  \end{center}
  Ta có $M$ là trung điểm $AB$, $N$ là trung điểm $CD$, $P$ là trung điểm $SA$ và  $O$ là trung điểm $AC$.\\
  Suy ra $MN\parallel BC$, $MN\parallel AD$, $MP\parallel SP$  và $OP\parallel SC$.
  \begin{itemchoice}
   \itemch {\bf Đúng.} \\ 
   Vì $\heva{&MN\parallel BC\\& BC\subset  (SBC)} \Rightarrow MN\parallel (SBC)$.
   \itemch {\bf Đúng.}\\ 
   Vì $\heva{&MN\parallel AD\\& AD\subset (SAD)}  \Rightarrow MN\parallel (SAD)$.
   \itemch {\bf  Sai.}\\
   Vì $\heva{&SB\parallel MP\\& MP\subset (MNP)} \Rightarrow SB\parallel (MNP)$.  Do đó $SB$ không cắt mặt phẳng ($MNP)$.
    \itemch {\bf Sai.}\\
   Vì $\heva{&OP\parallel  SC\\& OP\subset (MNP)} \Rightarrow SC\parallel (MNP)$. Do đó $SC$ không cắt mặt phẳng ($MNP)$.
  \end{itemchoice}
 } 
\end{ex}

% nguồn: Ok_dungsai_lop11_Lan1
\begin{ex}%[1H4H3-2]%[ma:89d43d60]
 Cho hình chóp $S.ABCD$ có đáy là hình  bình hành; điểm $M$ là trung điểm $SD$. Gọi $I$ là giao điểm của $BM$ và $(SAC)$; $J$ là giao điểm của $SA$ và $(BCM)$.
 \choiceTF 
   {\True $I$ là trọng tâm $\triangle SAC$}
   {$BM=3BI$}
   {\True $J$ là trung điểm của $SA$}
  {\True $MJ\parallel BC\parallel AD$} 
  \loigiai{
  \begin{center}
   \begin{tikzpicture}[scale=1]
     \path (0,0) coordinate (A)++(0:4) coordinate (B)++(-130:2)  coordinate (C)
       ($(A)!.5!(C)$) coordinate (O)++(105:3.5) coordinate (S)
      ($(O)!-1!(B)$) coordinate  (D)
      ($(S)!.5!(D)$)  coordinate (M)
      (intersection of  B--M and S--O) coordinate (I)
       (intersection of C--I and S--A) coordinate (J);
     \draw (D)--(C)--(B)--(S)--cycle 
      (S)--(C); 
      \draw[dashed] (D)--(A)--(B)--cycle
       (S)--(A)--(C)--(D)
      (S)--(O) 
       (M)--(B)
      (C)--(J)--(M); 
     \foreach \a/\b in {A/170,B/0,C/-30,D/200,O/-90,S/90,M/170,I/50,J/0}{ 
      \fill[draw=black] (\a) circle (.7pt) 
      ($(\a)+(\b:2mm)$)  node[scale=.7]{$\a$};
      }
    \end{tikzpicture}
  \end{center}
  Trong mặt phẳng $(SBD)$,  gọi $I=BM\cap SO$. \\
  Ta có $\heva{&I\in BM \\ &I\in SO\subset (SAC)}\Rightarrow I=BM\cap (SAC)$. \\ 
  Do $I$ là giao điểm của hai đường  trung tuyến $SO$ và $BM$ nên $I$ là trọng tâm $\triangle SBD$.\\
  Suy ra $BM=\dfrac{3}{2}\cdot BI$ và $SO=\dfrac{3}{2}\cdot SI\Rightarrow  I$ là trọng tâm $\triangle SAC$.\\
  Trong mặt phẳng  $(SAC)$, kéo dài $CI$ cắt $SA$ tại $J$, suy ra $J$ là trung điểm $SA$.\\
  Do $M, J$ là trung điểm của $SD, SA$ nên $MJ\parallel AD\parallel BC$. 
  \begin{itemchoice}
   \itemch $I$ là  trọng tâm $\triangle SAC$.
   \itemch  $BM=\dfrac{3}{2}\cdot BI$.
   \itemch $J$  là trung điểm $SA$.
   \itemch $MJ\parallel  AD\parallel BC$.
  \end{itemchoice}
 } 
\end{ex}

% nguồn: Ok_dungsai_lop11_Lan1
\begin{ex}%[1H4H3-2]%[ma:a39f2fa5]
 Cho hình bình hành  $ABCD$ và $ABEF$ nẳm ở hai mặt phẳng khác nhau. Gọi $M$ là trọng tâm $\triangle ABE$. Gọi $(P)$ là mặt phẳng đi qua $M$ và song song với mặt $(ADF)$. Lấy $N$ là giao điểm của $(P)$ và $AC$. Khi đó
 \choiceTF 
   {$EFDC$ là hình thang}
   {\True $FD \parallel EC$}
  {\True $(ADF) \parallel(BCE)$} 
  {$\dfrac{AN}{NC}=3$} 
 \loigiai{ \begin{center}
  \begin{tikzpicture}[line join=round,line cap=round,>=stealth,thick]
    \coordinate (B)  at (-4,-3);
    \coordinate  (C) at (0,-3);
    \coordinate  (A) at (-1.5,-1.2);
    \coordinate (D) at (2.5,-1.2); 
    \coordinate (F)  at (1,2.8);
     \coordinate (E) at (-1.5,1);
    \coordinate (J)  at (-0.8,1.5);
     \coordinate (K) at (0.7,-2.5);
    \coordinate  (I) at (-3.3,-2.5);
    \coordinate  (S) at ($(A)!0.5!(B)$);
    \coordinate (M) at (intersection  of J--I and E--S);
    \coordinate (N) at (intersection of  I--K and A--C);
    \draw (E)--(F) 
      (E)--(B)
     (E)--(C) 
      (B)--(C)
     (J)--(K) 
     (F)--(D) 
     (C)--(D); 
    \draw[dashed](E)--(S) 
     (J)--(I) 
     (I)--(K) 
     (A)--(C) 
     (E)--(A) 
      (F)--(A)
      (A)--(D)
      (A)--(B);
    \foreach \i in {E,J,F}{\draw (\i)  node[scale=0.9,above]{$\i$};}
    \foreach  \i in {C,K,D}{\draw (\i) node[scale=0.9,right]{$\i$};}
    \foreach \i in  {M,I,A,B}{\draw (\i) node[scale=0.9,left]{$\i$};}
    \foreach \i  in {N,S}{\draw (\i) node[scale=0.9,below]{$\i$};}
    \foreach  \i in {A,B,C,D,E,F,J,M,I,S,N,K}{\draw[fill=](\i) circle (1pt);}
   \end{tikzpicture}
 \end{center}
 \begin{itemchoice}
  \itemch  \textbf{Sai}.
  Ta có $\heva{&EF \parallel CD\parallel AB \\& EF=CD=AB} \Rightarrow  EFDC$ là hình bình hành.
  \itemch  \textbf{Đúng}. Do $EFDC$ là hình bình hành $\Rightarrow FD \parallel EC$.
  \itemch \textbf{Đúng}. Ta có $\heva{&AD \parallel BC,\, AF \parallel BE \\ &AD, AF \subset(ADF),\, AD \cap AF=A \\  &BC,\, BE \subset(BEC),\, BC \cap BE=B}\Rightarrow(ADF)\parallel(BCE)$.
  \itemch \textbf{Sai}.  \\
  Vẽ $(P)$ chứa  $M$ và $(P)\parallel(ADF)$ cắt $AB$, $AC$, $CD$, $EF$ lần lượt tại $I$, $N$, $K$, $J$.\\
  Ta có $\dfrac{AI}{BI}=\dfrac{AN}{NC}\,(IN  \parallel BC)$.\\
  Ta có $\dfrac{EJ}{IS}=\dfrac{ME}{MS}=2\,  (IS\parallel JE)$,
  $BI=EJ$ (tứ giác  $BIJE$ là hình bình hành).\\
  $\Rightarrow  \dfrac{BI}{IS}=2 \Rightarrow \dfrac{BI}{2}=\dfrac{IS}{1}=\dfrac{BI+IS}{2+1}=\dfrac{BS}{3}$.\\
  $\Rightarrow  BI=\dfrac{2}{3}BS$, $IS=\dfrac{1}{3} BS$.\\
  Ta có $AI=AS+AI=BS+\dfrac{1}{3}  BS=\dfrac{4}{3}BS$.\\
  $\Rightarrow  \dfrac{AI}{BI}=\dfrac{\dfrac{4}{3}BS}{\dfrac{2}{3}BS}=2\Rightarrow \dfrac{AN}{NC}=2$.
\end{itemchoice} }
\end{ex}

% nguồn: Ok_dungsai_lop11_Lan1
\begin{ex}%[1H4H3-2]%[ma:bd7c97a9]
 Cho  hình chóp $S.ABCD$ có đáy là hình bình hành tâm $O$. Gọi $M$, $N$ lần lượt là trung điểm của $SA$ và $SD$. Khi đó
 \choiceTF 
  {\True $MN \parallel(SBC)$} 
  {\True $(OMN)  \parallel(SBC)$}
  {Gọi  $E$ là trung điểm đoạn $AB$ và $F$ là một điểm thuộc đoạn $ON$. Khi đó $EF$ cắt với mặt phẳng $(SBC)$}
  {Gọi $G$ là một điểm trên mặt phẳng $(ABCD)$ cách đều $AB$ và  $CD$. Khi đó $GN$ cắt $(SAB)$}
 \loigiai{ 
  \begin{center}
   \begin{tikzpicture}[line join=round,line cap=round,>=stealth,thick]
     \coordinate  (A) at (-2,1);
      \coordinate (B) at (-4,-1);
      \coordinate (C) at (1,-1);
     \coordinate (D) at (2,1); 
      \coordinate (S) at (-1.8,5);
     \coordinate (M) at  ($(A)!0.5!(S)$);
     \coordinate  (N) at ($(D)!0.5!(S)$);
     \coordinate (E) at ($(A)!0.5!(B)$); 
      \coordinate (F) at (-0.15,2);
      \coordinate (I) at ($(C)!0.5!(B)$);
     \coordinate (H)  at ($(A)!0.5!(D)$);
     \coordinate (O) at (intersection of A--C  and B--D);
     \coordinate (K)  at ($(O)!0.5!(D)$);
     \draw (S)--(B) 
       (S)--(C)
      (S)--(D) 
      (B)--(C) 
       (C)--(D);
     \coordinate (G) at (intersection of  O--H and A--K);
      \draw[dashed](S)--(A)
       (A)--(B)
       (A)--(C)
      (A)--(D) 
      (B)--(D) 
      (M)--(N) 
       (E)--(F)
      (M)--(O) 
      (N)--(O) 
       (M)--(E)
      (E)--(O) 
       (I)--(H)
       (N)--(G);
     \foreach  \i in {S}{\draw (\i) node[scale=0.9,above]{$\i$};}
     \foreach  \i in {A,D,F,N,M}{\draw (\i) node[scale=0.9,above right]{$\i$};}
     \foreach \i  in {B,C,O,E,I}{\draw (\i) node[scale=0.9,below]{$\i$};}
     \foreach  \i in {S,A,B,C,D,O,E,F,M,N}{\draw[fill=](\i) circle (1pt);}
     \foreach \i in  {G}{\draw (\i) node[scale=0.9,below right]{$\i$};}
      \foreach \i in {S,A,B,C,D,O,E,F,M,N,G,I}{\draw[fill=](\i) circle (1pt);}
    \end{tikzpicture}
  \end{center}
  \begin{itemchoice}
   \itemch \textbf{Đúng}. Vì $MN$ là đường trung bình của tam giác $SAD$ nên $$MN \parallel AD \Rightarrow MN \parallel  BC \Rightarrow MN \parallel(SBC).\quad(1)$$
   \itemch \textbf{Đúng}. Ta có $O$, $N$ theo thứ tự là trung điểm của $BD$, $SD$ nên $ON$ là đường trung bình của  tam giác $SBD \Rightarrow ON \parallel SB \Rightarrow ON \parallel(SBC). \quad(2)$\\ Từ $(1)$ và $(2)$ suy ra $(OMN) \parallel(SBC)$.
   \itemch \textbf{Sai}. Ta có $OE$ là đường trung bình của tam giác  $ABD$ nên $OE \parallel AD \Rightarrow OE \parallel MN$.\\
   Do đó $E  \in(OMN)$.\\
   Mặt khác $F \in ON,\, ON  \subset(OMN) \Rightarrow F \in(OMN)$.\\
   Ta có $\heva{&EF \subset(OMN)\\&(OMN) \parallel(SBC)} \Rightarrow EF \parallel(SBC)$. 
    \itemch \textbf{Sai}.
   Vì $G$ thuộc mặt phẳng $(ABCD)$ và  cách đều $AB$, $CD$ nên $G$ thuộc đường trung bình của hình bình hành $ABCD$.\\
   Gọi  $I$ là trung điểm $BC$ thì $I$, $O$, $G$ thẳng hàng.\\
   Ta có $OI$  là đường trung bình của $\triangle ABC$ nên $OI \parallel AB \Rightarrow OI \parallel(SAB)\quad(3)$.\\
   Tương tự, ta có $ON \parallel SB  \Rightarrow ON \parallel(SAB). \quad(4)$.\\
   Từ $(3)$, $(4)$ suy ra $(OIN) \parallel(SAB)$ mà $GN \subset(OIN)$ nên $GN  \parallel(SAB)$.
 \end{itemchoice}}
\end{ex}

% nguồn: Ok_dungsai_lop11_Lan1
\begin{ex}%[1H4H3-2]%[ma:8e14b07b]
 Cho  hình chóp $S.ABCD$ có đáy $ABCD$ là hình bình hành. Gọi $M$, $N$, $P$ lần lượt là trung điểm của $AB$, $CD$ và $SA$.
  \choiceTF
  {$MP  \parallel SB$}
   {$NP \parallel SD$}
  {\True $MN\parallel (SAD)$} 
  {\True $SB\parallel  (MNP)$}
  \loigiai{
 \immini{ 
  \begin{itemchoice}
    \itemch $\triangle ABS$ có $MP$ là đường trung bình nên $MP \parallel SB$.
   \itemch $MP$ và $SD$ không đồng phẳng. 
   \itemch Vì $MN\parallel AD$ nên \\ 
    $\heva{&MN\parallel AD\\&AD\subset (SAD)\\&MN\not\subset (SAD)}\Rightarrow MN\parallel (SAD)$.
   \itemch Vì  $SB\parallel MP$ nên\\ $\heva{&SB\parallel MP\\&MP\subset (MNP)\\&SB\not\subset (MNP)}\Rightarrow SB\parallel (MNP)$.
  \end{itemchoice}}{\vspace{2cm}\begin{tikzpicture}[>=stealth,line join=round,line cap = round,scale=0.7,font=\footnotesize]
  \coordinate (A)  at (0,0);
  \coordinate  (D) at (4,0);
   \coordinate (B) at (-2,-2);
   \coordinate (C) at ($(D)-(A)+(B)$);
  \coordinate (S) at  ($(A)+(0,4)$);
  \coordinate (M) at ($(A)!0.5!(B)$); 
  \coordinate (N) at  ($(D)!0.5!(C)$);
  \coordinate  (P) at ($(S)!0.5!(A)$);
  \draw[dashed] (A) node[left]{$A$}--(B)  node[below]{$B$};
   \draw (B)--(C) node[right]{$C$}--(D) node[right]{$D$};
   \draw[dashed](D)--(A)--(S) node[above]{$S$};
  \draw (S)--(B); \draw (S)--(C); \draw (S)--(D); 
  \draw[dashed](M) node[left]{$M$}--(N) node[right]{$N$}--(P)  node[right]{$P$}--(M);
 \end{tikzpicture}
   }
  }
\end{ex}

% nguồn: Ok_dungsai_lop11_Lan1
\begin{ex}%[1H4H3-2]%[ma:2f5a1356]
 Cho tứ diện $ABCD$ có điểm $G$ là trọng tâm tam giác $ABD$ và điểm $M$ thuộc cạnh $BC$ sao cho $MB=2MC$. Khẳng  định nào dưới đây là đúng?
  \choiceTF
   {$MG $ cắt $AC$}
  {$MG \parallel AB  $}
  {\True $MG  \parallel (ACD) $}
   {$(BMG)\cap (ACD)=MG$}
  \loigiai{
  \begin{center}
   \begin{tikzpicture}[scale=.7,>=stealth,font=\footnotesize,line join=round,line cap=round]
     \foreach \x/\y/\t in {0/0/B,1.3/-1.6/C,4.5/0/D,1/3.5/A}{ 
      \path (\x,\y) coordinate  (\t);
     } 
     \path 
      (barycentric cs:A=1,B=1,D=1)  coordinate (G)
      ($(B)!2/3!(C)$) coordinate  (M)
      ($(A)!.5!(D)$)  coordinate (I);
      \draw[dashed] (D)--(B)--(I)
       (M)--(G);
      \draw (A)--(B)--(C)--(D)--(A)--(C)--(I);
     \foreach \t/\g  in {A/90,B/180,C/-90,D/0,G/90,M/190,I/60}{
      \draw[fill=black] (\t) circle (1pt) node[shift={(\g:7pt)},font=\scriptsize]{$ \t $}; 
      }
    \end{tikzpicture}
  \end{center}
  \begin{itemchoice}
   \itemch{\bf  Sai.}\\
   Vì $MG$ và $AC$  không đồng phẳng nên $MG$ không cắt $AC$.
   \itemch{\bf Sai.}\\ 
   Vì $MG$ và $AB$ không đồng  phẳng nên $MG$ không song $AB$.
    \itemch{\bf Đúng.}\\
   Gọi $I$ là trung điểm $AD$. Vì $G$ là trọng tâm tam giác $ABD$  nên $\dfrac{BG}{BI}=\dfrac{2}{3}$.\\
   Lại có $\dfrac{BM}{BC}=\dfrac{2}{3}$  nên $MG\parallel CI$ (định lí Ta-lét đảo).\\
   Mà $CI \subset (ACD)$ nên  $MG\parallel (ACD)$.
   \itemch{\bf Sai.}\\ 
    Ta có $(BMG)\equiv (BIC)$ nên $(BMG)\cap (ACD)=CI$.
  \end{itemchoice}
 } 
\end{ex}

% nguồn: Ok_dungsai_lop11_Lan1
\begin{ex}%[1H4H3-2]%[ma:03ae1da5]
 Cho hình chóp  $S.A B C D$ có đáy $A B C D$ là hình bình hành tâm $O$. Gọi $M, N, P$ lần lượt là trung điểm các cạnh $A B, C D$ và $SA$. Các khẳng định sau \textbf{đúng} hay \textbf{sai}?
 \choiceTF 
  {$MP \parallel  SC $}
  {\True $OP \parallel SC$} 
  {\True $M N \parallel(S  B C)$}
  {\True $M N \parallel(S  A D)$}
  \loigiai{
  \begin{center}
   \begin{tikzpicture}[>=stealth,line join=round,line cap=round,font=\footnotesize,scale=0.8]
      \path
       (0,0) coordinate (A)
      (-1.5,-2)  coordinate (B)
      (5,0) coordinate (D) 
      ($(B)+(D)-(A)$) coordinate  (C)
       ($(A)+(0.5,4)$) coordinate (S)
      ($(S)!1/2!(A)$)  coordinate (P)
      ($(A)!1/2!(B)$) coordinate  (M)
      ($(C)!1/2!(D)$) coordinate (N) 
      (intersection  of A--C and B--D) coordinate (O);
      \draw[dashed] (S)--(A)--(B)--(D)--(A)--(C)
       (S)--(O)--(P)--(M)--(N)--(P);
      \draw (S)--(B)--(C)--(S)--(D)--(C);;
     \foreach \x/\y  in {A/30,B/180,C/-10,D/10,O/-90,M/-70,N/-30,S/90,P/20}
     \draw[fill=black] (\x) circle (1.1pt) +  (\y:0.5cm) node {$\x$};
    \end{tikzpicture}
  \end{center}
  \begin{itemchoice}
   \itemch Đúng. Vì $OP$ là đường trung bình của tam giác $S A  C$ nên $OP \parallel SC$
   \itemch Đúng. Vì $M N$ là đường trung bình của hình bình hành $A B C D$ nên $M  N \parallel B C$, mà $B C \subset(S B C) \Rightarrow M N \parallel(S B C)$.
   \itemch Đúng. Vì $M N \parallel A D$, $A D \subset(S A D) \Rightarrow  M N \parallel(S A D)$.
  \end{itemchoice}
  }
\end{ex}

% nguồn: Ok_dungsai_lop11_Lan1
\begin{ex}%[1H4H3-2]%[ma:1cc052db]
 Cho hình chóp tứ giác $S.ABCD$ và lấy một điểm $E$ thuộc cạnh $SA$ của hình chóp ($E$ khác $S$, $A$). Trong mặt phẳng $(ABCD)$ vẽ một đường thẳng $d$ cắt các cạnh $CB$,  $CD$ lần lượt tại $M$, $N$ và cắt các tia $AB$, $AD$ lần lượt tại $P$, $Q$, đường thẳng $EP$ cắt cạnh $SB$ tại một điểm $H$.
  \choiceTF
  {\True $H$ là giao điểm của đường thẳng $SB$ và mp$(E,d)$} 
  {Các điểm $S$, $A$,  $D$, $E$, $Q$ cùng thuộc mặt phẳng $(SAB)$}
  {\True  Đường thẳng $MN$ nằm trong mặt phẳng mp$(E,d)$}
   {$MN$ là giao tuyến của hai mặt phẳng $(SAB)$ và mp$(E,d)$}
  \loigiai{
  \begin{center}
   \begin{tikzpicture}[line join=round,line cap=round,thick,scale=1]
     \tikzset{every  node/.style={scale=0.9}}
     \coordinate  (S) at (1,4);
      \coordinate (A) at (0,0);
     \coordinate (B) at  (1,-2);
     \coordinate (C) at (3,-3); 
     \coordinate  (D) at (5,0);
     \coordinate  (E) at ($(S)!0.4!(A)$);
      \coordinate (P) at ($(A)!1.3!(B)$);
     \coordinate (Q) at ($(A)!1.4!(D)$); 
      \coordinate (M) at (intersection of P--Q and B--C);
     \coordinate (N) at (intersection  of P--Q and C--D);
     \coordinate (H) at (intersection of  E--P and S--B);
      \draw (S)--(A)
       (S)--(B)
       (S)--(C)
       (S)--(D)
      (P)--(B)--(C) 
      (A)--(B)--(C)--(D)--(Q) 
      ($(P)!-0.1!(Q)$)--(M) 
      (N)--($(Q)!-0.1!(P)$) node[below]{$d$} 
      (E)--(P); 
     \draw[dashed,thin](A)--(D) 
      (M)--(N); 
     \foreach \i/\g in {S/90,A/180,B/180,C/-90,D/90,M/-90,N/-70,P/-90,Q/-90,E/180,H/-180}{\draw[fill=black](\i)  circle (1pt)
      ($(\i)+(\g:3.5mm)$) node[scale=1]{$\i$};} 
    \end{tikzpicture}
  \end{center}
  \begin{itemchoice}
   \itemch \textbf{Đúng.}\\ 
   Vì $E$ thuộc cạnh $SA$ nên  $E$ thuộc mặt phẳng $(SAB)$.\\
   Vì $P$ thuộc đường  thẳng $AB$ nên $P$ thuộc mặt phẳng $(SAB)$.\\
   Như vậy, các điểm $S$, $A$, $B$, $E$, $P$ cùng thuộc mặt phẳng  (SAB).\\
   Trong tam giác  $SAB$, đường thẳng $EP$ cắt cạnh $SB$ tại một điểm $H$.\\
   Do $P$ thuộc đường thẳng $d$ nên $EP$ nằm trong mp$(E,d)$ và $H$ thuộc $EP$, do đó  $H$ thuộc mp$(E,d)$.\\
   Vậy $H$  là giao điểm của đường thẳng $SB$ và mp$(E,d)$.
    \itemch \textbf{Sai}\\
   Vì $E$ thuộc cạnh  $SA$ nên $E$ thuộc mặt phẳng $(SAD)$.\\
   Vì $Q$ thuộc đường  thẳng $AD$ nên $Q$ thuộc mặt phẳng $(SAD)$.\\
   Như vậy, các điểm $S$, $A$, $D$, $E$, $Q$ cùng thuộc mặt  phẳng $(SAD)$.
   \itemch  \textbf{Đúng.}\\
   Đường thẳng  $d$ cắt các cạnh $CB$, $CD$ lần lượt tại $M$, $N$.\\
   Do đó $M$, $N$ thuộc $d$, mà $d$  nằm trong mp$(E,d)$ nên đường thẳng $MN$ cũng nằm trong mp$(E,d)$.
   \itemch  \textbf{Sai.}\\
   Ta có đường  thẳng $MN$ cũng nằm trong mp$(E,d)$.\\
   Ta lại có  $M$ thuộc $CB$ nằm trong mặt phẳng $(ABCD)$ nên $M$ thuộc mặt phẳng $(ABCD)$.\\
   Tương tự $N$ thuộc  $CD$ nằm trong mặt phẳng $(ABCD)$ nên $N$ thuộc mặt phẳng $(ABCD)$.\\
   Do đó đường  thẳng $MN$ nằm trong mặt phẳng $(ABCD)$.\\
   Vậy $MN$ là giao tuyến của hai mặt phẳng  $(ABCD)$ và mp$(E,d)$.
  \end{itemchoice}
 } 
\end{ex}

% nguồn: Ok_dungsai_lop11_Lan2
\begin{ex}%[1H4H3-2]%[ma:f21573ce]
	Cho hình chóp $S. ABCD$ có đáy $ABCD$ là hình bình hành tâm $O$. Gọi $M$, $N$ lần lượt là trung điểm các cạnh $AB$ và $CD$, $P$ là trung điểm cạnh  $SA$. Khi đó:
	 \choiceTF
		 {\True $MN\parallel(SBC)$}
		{\True $MN\parallel(SAD)$} 
		{$SB$ cắt  với mặt phẳng $(MNP)$}
		 {$SC$ cắt với mặt phẳng $(MNP)$}
	 \loigiai{
		\begin{center}
			\begin{tikzpicture}[scale=1,font=\footnotesize,line join=round,line cap=round,>=stealth]
					\def\canhAD{4}; \def\canhBA{2}; \def\gocBAD{-140}; \def\h{3};  \def\xdinhS{1};
					\coordinate  (A) at (0,0);
					\coordinate  (B) at ($(A)+(\gocBAD:\canhBA)$);
					 \coordinate (C) at ($(B)+(0:\canhAD)$);
					\coordinate (D) at  ($(A)+(0:\canhAD)$);
					\coordinate  (S) at ($(A)+(\xdinhS,\h)$);
					\coordinate (O)  at (intersection of A--C and B--D);
					\path 
						($(A)!.5!(B)$)  coordinate (M)
						($(A)!.5!(S)$) coordinate (P) 
						 ($(D)!.5!(C)$) coordinate (N);
					\draw (B)--(S)--(C)--cycle (S)--(D)--(C); 
					 \draw[dashed] (A)--(D)
						(A)--(C) 
						 (B)--(D)
						 (S)--(A)--(B)
						 (M)--(P)--(N)--cycle
						 (P)--(O);
					 \foreach \x/\y in {A/180,B/-90,C/-90,D/0,S/90,O/225,M/-80,N/0,P/45}{\fill (\x) circle (1pt)
						 ($(\x)+(\y:0.3cm)$) node{$\x$};}
				\end{tikzpicture}
		\end{center}
		\begin{itemchoice}
			\itemch {\bf  Đúng}.\\
			Vì $MN$ là  đường trung bình của hình bình hành $ABCD$ nên $MN\parallel BC$. \\
			 Mà $BC\subset(SBC)\Rightarrow MN\parallel(SBC)$.
			\itemch {\bf  Đúng}.\\
			Vì $MN$ là đường trung bình của hình bình hành  $ABCD$ nên $MN\parallel AD$. \\
			 Mà $AD\subset(SAD)\Rightarrow MN\parallel(SAD)$.
			\itemch {\bf Sai}.\\ 
			Ta có $MP$  là đường trung bình của tam giác $SAB$ nên $SB\parallel MP$. \\
			Mà $MP\subset(MNP)$ nên $SB\parallel(MNP)$. 
			 \itemch {\bf Sai}.\\
			$OP$ là đường  trung bình của tam giác $SAC$ nên $SC\parallel OP$. \\
			Mà  $OP\subset(MNP)$ nên $SC\parallel(MNP)$.
		\end{itemchoice}
	} 
\end{ex}

% nguồn: Ok_dungsai_lop11_Lan2
\begin{ex}%[1H4H3-2]%[1H4H4-2]%[1H4H2-4]%[ma:2e22f776]
	Cho hình chóp $S.ABCD$ có đáy là hình bình hành tâm $O$. Gọi $M$, $N$ lần lượt là trung điểm của $SA$ và $SD$.  Khi đó:
	\choiceTF 
		{\True  $MN\parallel(SBC)$}
		 {\True $(OMN)\parallel(SBC)$}
		{Gọi $E$ là trung điểm đoạn $AB$ và $F$ là một điểm thuộc đoạn $ON$. Khi đó $EF$ cắt với mặt  phẳng $(SBC)$}
		{Gọi $G$ là một điểm trên mặt phẳng $(ABCD)$ cách đều $AB$ và $CD$.  Khi đó $GN$ cắt $(SAB)$}
	\loigiai{ 
		\begin{center}
			\begin{tikzpicture}[font=\scriptsize]
					 \def\a{4}
					\path 
						(0:0)  coordinate (A)
						(225:2.75)  coordinate (B)
						++(0:4) coordinate (C) 
						($(A)+(C)-(B)$) coordinate  (D)
						($(A)+(80:3)$) coordinate (S) 
						($(S)!1/2!(A)$) coordinate  (M)
						 ($(S)!1/2!(D)$) coordinate (N)
						($(A)!1/2!(B)$)  coordinate (E)
						 (intersection of A--C and B--D) coordinate (O)
						($(O)!.75!(N)$) coordinate  (F);
					 \draw[dashed,thin] (A)--(B)--(D)--cycle (C)--(A)--(S)
						(M)--(N)--(O)--(E)--cycle (O)--(M) 
						(E)--(F); 
					\draw[] (B)--(C)--(D) 
						(B)--(S) 
						 (C)--(S)
						(D)--(S); 
					\foreach \x/\y in  {A/135,B/-135,C/-45,D/45,S/90,O/-90,M/30,N/80,E/-90,F/0}
					\fill[black] (\x) circle  (1pt)
						 ($(\y:3mm)+(\x)$) node {$\x$};
				\end{tikzpicture}
		\end{center}
		\begin{itemchoice}
			\itemch Đúng.\\ 
			Vì $MN$ là  đường trung bình của tam giác $SAD$\\
			nên $MN\parallel AD\Rightarrow MN\parallel BC\Rightarrow MN\parallel  (SBC)$.\quad (1)
			 \itemch Đúng.\\
			Tương tự, ta có $O$, $N$ theo thứ tự là trung điểm của  $BD$, $SD$ nên $ON$ là đường trung bình của tam giác $SBD \Rightarrow ON\parallel SB\Rightarrow ON\parallel (SBC)$.\quad (2)\\
			Từ $(1)$ và $(2)$ suy ra $(OMN)\parallel  (SBC)$.
			\itemch Sai.\\ 
			Ta có $OE$ là đường trung bình của tam giác $ABD$ nên $OE\parallel AD  \Rightarrow OE\parallel MN$.\\
			Do đó $E\in (OMN)$.\\ 
			Mặt khác $F\in ON,ON\subset (OMN)\Rightarrow  F\in (OMN)$.\\
			Ta có $\heva{&EF\subset  (OMN)\\&(OMN)\parallel (SBC)}$ $\Rightarrow EF\parallel (SBC)$.
			 \itemch Sai.
			\begin{center}
				\begin{tikzpicture}[font=\scriptsize]
						\def\a{4} 
						 \path
						 (0:0) coordinate (A)
						(225:2.75) coordinate  (B)
						++(0:4) coordinate (C) 
						($(A)+(C)-(B)$)  coordinate (D)
						($(A)+(80:3)$) coordinate (S) 
						 ($(S)!1/2!(A)$) coordinate (M)
						($(S)!1/2!(D)$) coordinate (N) 
						 ($(A)!1/2!(B)$) coordinate (E)
						(intersection of A--C and B--D)  coordinate (O)
						($(O)!.75!(N)$) coordinate  (F)
						($(B)!1/2!(C)$) coordinate  (I)
						($(A)!1/2!(D)$) coordinate (K) 
						($(I)!.75!(K)$) coordinate  (G);
						 \draw[dashed,thin] (A)--(B)--(D)--cycle (C)--(A)--(S)
						(M)--(N)--(O)--(E)--cycle (O)--(M) 
						 (E)--(F)
						 (I)--(K)
						(N)--(G); 
						\draw[]  (B)--(C)--(D)
						 (B)--(S)
						 (C)--(S)
						 (D)--(S);
						 \foreach \x/\y in {A/135,B/-135,C/-45,D/45,S/90,O/-90,M/30,N/80,E/-90,F/0,I/-90,G/-60}
						\fill[black] (\x) circle (1pt) 
						 ($(\y:3mm)+(\x)$) node {$\x$};
					\end{tikzpicture}
			\end{center}
			Vì $G$ thuộc mặt phẳng $(ABCD)$ và cách đều $AB$, $CD$ nên $G$ thuộc đường trung bình của hình bình hành $ABCD$ (ứng  với hai cạnh $AB$, $CD$).\\
			Gọi $I$ là trung điểm  $BC$ thì $I$, $O$, $G$ thẳng hàng.\\
			Ta có $OI$ là đường trung bình của $\Delta ABC$ nên  $OI\parallel AB\Rightarrow OI\parallel (SAB)$.\quad (3)\\
			Tương  tự, ta có $ON\parallel SB\Rightarrow ON\parallel (SAB)$.\quad (4)\\
			 Từ (3), (4) suy ra $(OIN)\parallel (SAB)$ mà $NG\subset (OIN)$ nên $NG\parallel (SAB)$.
		\end{itemchoice}
	 }
\end{ex}

% nguồn: Ok_dungsai_lop11_Lan2
\begin{ex}%[1H4H3-2]%[1H4H4-2]%[ma:c8464009]
	Cho hình chóp $S.ABCD$ có đáy  $ABCD$ là hình bình hành. Gọi $H$, $I$, $K$ lần lượt là trung điểm của $SA$, $SB$, $SC$. Gọi $M$ là giao điểm của $AI$ và $KD$, $N$ là giao điểm của $DH$ và $CI$. Khi đó:
	 \choiceTF
		 {\True $HI\parallel(ABCD)$}
		 {\True $(HIK)\parallel(ABCD)$}
		{$SM$ và $HI$ chéo  nhau}
		 {$(SMN)$ cắt $(HIK)$}
	 \loigiai{
		\begin{center}
			\begin{tikzpicture}[font=\footnotesize]
					 \def\a{4}
					\path 
						 (0:0) coordinate (A)
						(0:\a)  coordinate (B)
						++(-135:\a/2) coordinate (C) 
						 ($(A)+(C)-(B)$) coordinate (D)
						($(A)+(65:3)$)  coordinate (S)
						 ($(S)!1/2!(A)$) coordinate (H)
						 ($(S)!1/2!(B)$) coordinate (I)
						 ($(S)!1/2!(C)$) coordinate (K)
						(intersection of A--I and K--D) coordinate  (M)
						(intersection of I--C and  H--D) coordinate (N)
						(intersection of  D--M and C--N) coordinate (m)
						(intersection of D--M  and B--S) coordinate (n)
						(intersection of A--C and B--D)  coordinate (O);
					\draw[dashed,thin] (B)--(A)--(D) 
						(A)--(S) 
						(H)--(I)--(K)--cycle  (A)--(M)
						(D)--(N) 
						(m)--(N) 
						 (n)--(S);
					\draw[] (B)--  (C)--(D)
						 (C)--(S)
						 (D)--(S)
						(S)--(N)--(M)--cycle (D)--(M) 
						 (B)--(n)
						(C)--(m); 
					\foreach \x/\y  in {A/-90,B/0,C/-90,D/-90,S/90,H/-60,I/60,K/-90,M/0,N/90}
					\fill[black] (\x) circle  (1pt)
						 ($(\y:3mm)+(\x)$) node {$\x$};
				\end{tikzpicture}
		\end{center}
		\begin{itemchoice}
			\itemch Đúng.\\ 
			Vì $HI$  là đường trung bình của tam giác $SAB$ nên $HI\parallel AB$,\\
			mà $AB\subset (ABCD)\Rightarrow  HI\parallel (ABCD)$.\quad (1)\\
			\itemch Đúng.\\ 
			 Tương tự ta có $KI\parallel BC$, $BC\subset (ABCD) \Rightarrow KI\parallel (ABCD)$.\quad (2)\\
			Mặt  khác $HI\subset (HKI)$, $KI\subset (HKI)$, $HI\cap KI=I$.\qquad (3)\\
			Từ (1),  (2) và (3) suy ra $(HIK)\parallel (ABCD)$.
			\itemch Sai.\\ 
			Vì $\heva{&M\in AI,\,AI\subset (SAB)\\&M\in  DK,\,DK\subset (SCD)}$ $\Rightarrow M\in (SAB)\cap (SCD)$\\
			$\Rightarrow SM=(SAB)\cap  (SCD)$.\\
			\itemch Sai.\\ 
			Khi dó  $\heva{&(SAB)\cap(SCD)=SM\\ &AB\subset(SAB),\,CD\subset(SCD)\\ &AB\parallel CD\\}$
			$\Rightarrow SM\parallel AB\parallel CD\Rightarrow SM\parallel  HI$.\quad (1)\\
			 Vì $\heva{
			 &N\in DH,\, DH\subset (SAD) \\
			&N\in CI,\, CI\subset (SBC) } $ $\Rightarrow  N\in (SAD)\cap (SBC)$\\
			 $\Rightarrow SN=(SAD)\cap (SBC).$\\
			Khi đó,  ta có
			$\heva{&(SAD)\cap (SBC)=SN\\ &AD\subset(SAD),\,BC\subset(SBC)\\ &AD\parallel BC}$ $\Rightarrow SN\parallel AD\parallel BC\Rightarrow SN\parallel  KI$.\quad (2)\\
			Mặt khác ba điểm  $S$, $M$, $N$ không thẳng hàng.\quad (3)\\
			 Từ (1), (2), (3) suy ra $(SMN)\parallel (HIK)$.
		\end{itemchoice}
	} 
\end{ex}

% nguồn: Ok_dungsai_lop11_Lan1
\begin{ex}%[1H4V3-2]%[ma:c0cea7c8]
 Cho hình chóp $S.ABCD$  có đáy $ABCD$ là hình bình hành. Lấy điểm $M$ trên cạnh $AD$ sao cho $AD=3AM$. Gọi $G,N$ theo thứ tự là trọng tâm các tam giác $SAB,ABC$. Khi đó, các mệnh đề sau đúng hay sai?
  \choiceTF
  {Giao tuyến của  hai mặt phẳng $(SAB)$ và $(SCD.$ là đường thẳng đi qua $S$ và song song với $AC,BD$}
   {$\dfrac{DN}{DB}=\dfrac{1}{3}$}
  {\True $MN$ song  song với mặt phẳng $(SCD)$}
  {$NG$  cắt với mặt phẳng $(SAC)$}
  \loigiai{
  \begin{center}
   \begin{tikzpicture}[font=\footnotesize,scale=.7]
     \def\a{4} 
     \path (0:0) coordinate (A) 
      ++(0:\a) coordinate  (D)
      ++(-136:.7*\a) coordinate  (C)
      ($(A)+(C)-(D)$) coordinate  (B)
      ($(A)+(80:\a)$)  coordinate (S)
      ($(A)!1/3!(D)$)  coordinate (M)
      ($(A)!.5!(B)$) coordinate  (P)
      (intersection of A--C and B--D) coordinate  (O)
      (intersection of  P--C and B--D) coordinate (N)
      ($(S)!2/3!(P)$)  coordinate (G)
       ($(B)-(A)+(S)$) coordinate (Sx);
     \draw[dashed]  (B)--(A)--(D)
      (A)--(S) 
      (S)--(P)--(C) 
       (A)--(C)
      (B)--(D) 
      (M)--(N)--(G); 
      \draw (B)-- (C)--(D)
      (B)--(S) 
      (C)--(S) 
      (D)--(S) 
      (Sx)--(S)--([turn]0:1.5cm) node[right]{$x$}; 
     \foreach \x/\g  in {A/45,B/-135,C/-45,D/45,S/90,M/95,N/-60,P/180,O/-90,G/0}
     \fill[black] (\x)  circle (1pt)
      ($(\g:3.5mm)+(\x)$) node {$\x$}; 
    \end{tikzpicture}
  \end{center}
  \begin{itemchoice}
   \itemch Sai. Ta có $\heva{&S\in (SAB)\cap (SCD)\\&AB\parallel  CD\\&AB\subset (SAB),CD\subset (SCD)}\\
   \Rightarrow (SAB)\cap (SCD)=Sx$, (với $Sx$ qua $S$ và $Sx\parallel  AB\parallel CD$).
   \itemch Sai. Gọi $O$  là tâm hình bình hành $ABCD$.\\
   Vì $N$  là trọng tâm của $\triangle ABC$ nên $BN=\dfrac{2}{3} BO=\dfrac{2}{3} \cdot \dfrac{1}{2} BD=\dfrac{1}{3} BD\Rightarrow \dfrac{DN}{DB}=\dfrac{2}{3}$.
   \itemch Đúng. Mặt  khác, ta có $AD=3AM\Rightarrow \dfrac{DM}{DA}=\dfrac{2}{3}$.\\
   Xét tam giác $ADB$, ta có $\dfrac{DM}{DA}=\dfrac{DN}{DB}=\dfrac{2}{3}$ nên $MN\parallel AB\Rightarrow MN\parallel  CD$.\\
   Mà $CD\subset (SCD)\Rightarrow MN\parallel (SCD)$. 
    \itemch Sai. Chứng minh $NG$ song song $(SAC)$.\\
   Gọi $P$ là trung điểm  $AB$.\\
   Tam  giác $SPC$ có
   $\dfrac{PG}{PS}=\dfrac{PN}{PC}=\dfrac{1}{3}$ (tính chất trọng tâm)\\ $\Rightarrow NG\parallel SC,SC\subset  (SAC.\Rightarrow NG\parallel (SAC)$.
  \end{itemchoice}
 } 
\end{ex}

% nguồn: Ok_dungsai_lop11_Lan1
\begin{ex}%[1H4V3-2]%[ma:9a119e40]
 Cho hai hình bình hành $ABCD$ và $ABEF$ không cùng nằm trong một mặt phẳng và có tâm lần lượt là $O$ và $O'$. Gọi $M,N$ lần lượt là hai điểm trên các cạnh $AE,BD$ sao cho $AM=\dfrac{1}{3} AE$, $BN=\dfrac{1}{3} BD$. Khi đó, các  mệnh đề sau đúng hay sai?
 \choiceTF 
  {\True $OO'$ song song với mặt phẳng $(ADF)$} 
   {$OO'$ cắt mặt phẳng $(BCE)$}
  {$\dfrac{BN}{BD}=\dfrac{2}{3}$} 
  {\True $MN$ song song với mặt phẳng $(CDFE)$} 
 \loigiai{ 
  \begin{center}
   \begin{tikzpicture}
      \def\a{5}
      \path (0:0) coordinate (A)
      ++(0:\a) coordinate (B) 
      ++(-120:\a/2) coordinate (C) 
      ($(C)-(B)+(A)$)  coordinate (D)
       ($(A)+(105:.6*\a)$) coordinate (F)
       ($(F)-(A)+(B)$) coordinate (E)
       ($(A)!1/3!(E)$) coordinate (M)
      ($(B)!1/3!(D)$) coordinate  (N)
      (intersection of A--C and B--D)  coordinate (O)
      (intersection of A--E and B--F)  coordinate (O')
       (intersection of A--N and C--D) coordinate (I)
       (intersection of A--N and B--C) coordinate (Nt);
     \draw[dashed] (B)--(A)--(D)--cycle  (C)--(A)--(F)--(B)
      (A)--(E) 
      (O)--(O') 
       (A)--(Nt)
       (M)--(N);
     \draw (D)--(C)--(B)--(E)--(F)--cycle  (E)--(C)
      (Nt)--(I)--(C); 
     \foreach \x/\g  in {A/180,B/0,C/-45,D/180,E/45,F/135,O/-90,O'/90,M/135,N/-90,I/0}
     \fill[black]  (\x) circle (1pt)
      ($(\g:3mm)+(\x)$) node  {$\x$};
    \end{tikzpicture}
  \end{center}
  \begin{itemchoice}
   \itemch  Đúng. Chứng minh $OO'\parallel(ADF)$.\\
   Ta có $OO'$ là  đường trung bình của tam giác $BDF$ nên $OO' \parallel DF$.\\
   Mà $DF\subset (ADF)$ suy ra $OO' \parallel  (ADF)$.
    \itemch Sai. Chứng minh $OO\parallel(BCE)$.\\
   Ta có $OO'$ là  đường trung bình của tam giác $ACE$ nên $OO' \parallel CE$.\\
   Mà $CE\subset (BCE)$ suy ra $OO^{{'}} \parallel  (BCE)$.
   \itemch Sai.  Chứng minh $MN\parallel(CDFE)$.\\
    Trong mặt phẳng $(ABCD)$, gọi $I=AN\cap CD$. Do $AB\parallel CD$ nên $\dfrac{AN}{AI}=\dfrac{BN}{BD}=\dfrac{1}{3}$.\\
   Mặt khác $\dfrac{AM}{AE}=\dfrac{1}{3} \Rightarrow \dfrac{AN}{AI}=\dfrac{AM}{AE} \Rightarrow  MN\parallel IE$.\\
   Mà $IE\subset (CDFE)$, suy ra $MN\parallel  (CDFE)$.
   \itemch Đúng. Chứng  minh $MN\parallel(CDFE)$.\\
   Trong mặt phẳng $(ABCD)$,  gọi $I=AN\cap CD$. Do $AB\parallel CD$ nên $\dfrac{AN}{AI}=\dfrac{BN}{BD}=\dfrac{1}{3}$.\\
   Mặt khác $\dfrac{AM}{AE}=\dfrac{1}{3} \Rightarrow \dfrac{AN}{AI}=\dfrac{AM}{AE} \Rightarrow MN\parallel  IE$.\\
   Mà $IE\subset (CDFE)$,  suy ra $MN\parallel (CDFE)$.
  \end{itemchoice}
  }
\end{ex}

% nguồn: Ok_dungsai_lop11_Lan1
\begin{ex}%[1H4V3-2]%[ma:3a0a9657]
 \immini{Cho hình chóp $S.ABCD$ đáy $ABCD$ là hình bình hành. Gọi $I$, $J$ lần lượt là trọng tâm  của tam giác $SAB$ và $SCD$; $E$, $F$ lần lượt là trung điểm của $AB$ và $CD$. Khi đó}
 {\begin{tikzpicture}[scale=.75]
   \def\a{4}
   \path (0:0) coordinate (A)
   ++(0:\a) coordinate  (D)
   ++(-130:.55*\a) coordinate  (C)
   ($(A)+(C)-(D)$) coordinate (B) 
   ($(A)+(110:\a)$)  coordinate (S)
   ($(A)!.5!(B)$) coordinate  (E)
   ($(C)!.5!(D)$) coordinate (F) 
    ($(S)!2/3!(E)$) coordinate (I)
   ($(S)!2/3!(F)$) coordinate  (J)
   (intersection of A--C and  B--D) coordinate (O);
  \draw[dashed]  (B)--(A)--(D)
    (A)--(S)
    (S)--(E)--(F)
    (I)--(J);
  \draw (B)--  (C)--(D)
    (B)--(S)
   (C)--(S) 
   (D)--(S) 
   (S)--(F); 
   \foreach \x/\g in {A/45,B/-135,C/-45,D/45,S/90,E/-80,F/-45,I/80,J/70}
  \fill[black] (\x) circle (1pt) 
    ($(\g:3mm)+(\x)$) node {$\x$};
 \end{tikzpicture}
 } 
  \choiceTF
  {\True $\dfrac{SJ}{SF}=\dfrac{2}{3}$} 
  {\True $IJ\parallel (ABCD)$} 
  {\True  $BC$ song song với mặt phẳng $(SAD),(SEF)$}
  {$BC$ cắt mặt phẳng  $(AIJ)$}
  \loigiai{
  \begin{itemchoice}
   \itemch {\bf Đúng.}\\ 
   Ta có $J$ là trong  tâm $\Delta SCD$ suy ra $\dfrac{SI}{SE}=\dfrac{2}{3}$.
   \itemch {\bf Đúng.}\\ 
   Do $I$, $J$ lần lượt là trọng  tâm của tam giác $SAB$ và $SCD$ nên
   \[\dfrac{SI}{SE}=\dfrac{SJ}{SF}=\dfrac{2}{3} \Rightarrow  IJ\parallel EF.
   \] 
    Mà $ EF\subset (ABCD)\Rightarrow IJ\parallel (ABCD)$.
    \itemch {\bf Đúng.}\\
   Ta  có
   \begin{itemize}
    \item $BC\parallel AD,AD\subset  (SAD)\Rightarrow BC\parallel (SAD)$.
    \item $\heva{&BC\parallel EF\\&EF\subset  (SEF)}\Rightarrow BC\parallel (SEF)$.
   \end{itemize}
   \itemch {\bf  Sai.}\\
    Ta có $\heva{&IJ\parallel EF\\&EF\parallel BC}\Rightarrow BC\parallel IJ$.\\
   Mà $IJ\subset (AIJ)\Rightarrow  BC\parallel (AIJ)$.
  \end{itemchoice}
  }
\end{ex}

% nguồn: Ok_dungsai_lop11_Lan1
\begin{ex}%[1H4V3-2]%[ma:b4876465]
 Cho hình chóp $S.ABCD$ đáy $ABCD$ là hình bình hành. Gọi $I,J$ lần  lượt là trọng tâm của tam giác $SAB$ và $SCD;E,F$ lần lượt là trung điểm của $AB$ và $CD$. Khi đó, các mệnh đề sau đúng hay sai?
 \choiceTF 
  {\True $\dfrac{SJ}{SF}=\dfrac{2}{3}$} 
   {\True $IJ\parallel (ABCD)$}
  {\True $BC$ song song  với mặt phẳng $(SAD),(SEF)$}
  {$BC$  cắt mặt phẳng $(AIJ)$}
 \loigiai{ 
  \begin{center}
   \begin{tikzpicture}[font=\footnotesize,scale=.7]
     \def\a{4} 
     \path  (0:0) coordinate (A)
      ++(0:\a)  coordinate (D)
      ++(-135:.6*\a) coordinate (C) 
      ($(A)+(C)-(D)$)  coordinate (B)
      ($(A)+(80:\a)$) coordinate  (S)
       ($(A)!.5!(B)$) coordinate (E)
       ($(C)!.5!(D)$) coordinate (F)
      ($(S)!2/3!(E)$) coordinate  (I)
       ($(S)!2/3!(F)$) coordinate (J);
     \draw[dashed] (B)--(A)--(D) 
       (A)--(S)
       (S)--(E)--(F)
       (I)--(J);
     \draw (B)-- (C)--(D) 
      (B)--(S) 
       (C)--(S)
       (D)--(S)
      (S)--(F); 
     \foreach \x/\g in  {A/-45,B/-135,C/-45,D/45,S/90,E/-45,F/-45,I/-45,J/45}
     \fill[black] (\x)  circle (1pt)
       ($(\g:3mm)+(\x)$) node {$\x$};
    \end{tikzpicture}
  \end{center}
  \begin{itemchoice}
   \itemch Đúng. Do $I,J$ lần lượt là trọng tâm của tam giác  $SAB$ và $SCD$ nên \[\dfrac{SI}{SE}=\dfrac{SJ}{SF}=\dfrac{2}{3}\quad(*).\]
   \itemch Đúng.  Ta có $\heva{&IJ\parallel EF\text{ do } (*)\\&EF\subset (ABCD)} \Rightarrow IJ\parallel (ABCD)$.
   \itemch Đúng. Vì $\heva{&BC\parallel AD\\&AD\subset  (SAD)}\Rightarrow BC\parallel (SAD)$.\\
   Vì $EF$ là đường trung bình của hình bình hành  $ABCD$ nên
   \[\heva{&BC\parallel  EF\\&EF\subset (SEF)}\Rightarrow BC\parallel (SEF).\]
   \itemch Sai. Ta có $\heva{&IJ\parallel EF\\&EF\parallel BC}\Rightarrow BC\parallel  IJ$. Mà $IJ\subset (AIJ)\Rightarrow BC\parallel (AIJ)$.
  \end{itemchoice}
 } 
\end{ex}

% nguồn: Ok_dungsai_lop11_Lan2
\begin{ex}%[1H4V3-2]%[ma:ecadfedb]
	Cho hình chóp $S.ABCD$ có đáy $ABCD$ là hình bình hành. Lấy điểm $M$ trên cạnh $AD$ sao cho $AD=3AM$. Gọi $G,N$ theo thứ tự là trọng tâm các tam  giác $SAB$, $ABC$. Khi đó
	\choiceTF 
		{Giao tuyến của hai mặt phẳng $(SAB)$ và $(SCD)$  là đường thẳng đi qua $S$ và song song với $AC$, $BD$}
		{$\dfrac{DN}{DB}=\dfrac{1}{3}$} 
		 {\True $MN$ song song với mặt phẳng $(SCD)$}
		{$NG$  cắt với mặt phẳng $(SAC)$}
	\loigiai{ 
		\begin{itemchoice}
			 \itemch{\bf Sai.}\\
			Ta  có $\heva{&S\in (SAB)\cap (SCD)\\&AB\parallel CD\\&AB\subset (SAB),CD\subset (SCD)}\Rightarrow (SAB)\cap (SCD)=Sx\parallel AB\parallel CD$.
			\begin{center}
				\begin{tikzpicture}
						 \def\a{4}
						\path (0:0)  coordinate (A)
						++(0:\a)  coordinate (D)
						 ++(-140:.6*\a) coordinate (C)
						($(A)+(C)-(D)$)  coordinate (B)
						($(A)!.5!(B)$) coordinate (P) 
						($(A)!1/3!(D)$)  coordinate (M)
						($(A)+(80:\a)$)  coordinate (S)
						($(P)!1/3!(S)$) coordinate (G) 
						($(B)+(S)-(A)$) node[above]{$x$} coordinate (Sx) 
						($(Sx)!1.4!(S)$) coordinate  (S')
						(intersection  of P--C and B--D) coordinate (N)
						(intersection of A--C and B--D) coordinate  (O);
						 \draw[dashed] (B)--(A)--(D)
						 (A)--(S)
						 (A)--(C)
						 (B)--(D)
						 (S)--(P)--(C)
						 (G)--(N)--(M);
						 \draw (B)-- (C)--(D)
						(B)--(S) 
						(C)--(S) 
						 (D)--(S)
						 (Sx)--(S');
						\foreach \x/\g in {A/45,B/-135,C/-45,D/0,S/90,O/75,P/180,M/90,N/-90,G/0} 
						 \fill[black] (\x) circle (1pt)
						($(\g:3mm)+(\x)$) node {$\x$}; 
					\end{tikzpicture}
			\end{center}
			\itemch{\bf  Sai.}\\
			Gọi  $O$ là tâm hình bình hành $ABCD$.\\
			Vì $N$ là trọng tâm của $\triangle ABC$ nên $BN=\dfrac{2}{3} BO=\dfrac{2}{3} \cdot \dfrac{1}{2} BD=\dfrac{1}{3}  BD\Rightarrow \dfrac{DN}{DB}=\dfrac{2}{3}$.
			\itemch{\bf Đúng.}\\ 
			Ta có  $AD=3AM\Rightarrow \dfrac{DM}{DA}=\dfrac{2}{3}$.\\
			Xét tam giác $ADB$, ta có  $\dfrac{DM}{DA}=\dfrac{DN}{DB}=\dfrac{2}{3}$ nên $MN\parallel AB\Rightarrow MN\parallel CD$.\\
			Mà $CD\subset  (SCD)\Rightarrow MN\parallel (SCD)$.
			 \itemch{\bf Sai.}\\
			Gọi $P$ là trung điểm $AB$.  Tam giác $SPC$ có $\dfrac{PG}{PS}=\dfrac{PN}{PC}=\dfrac{1}{3}$ (tính chất trọng tâm)\\ $\Rightarrow NG\parallel SC,SC\subset (SAC)\Rightarrow NG\parallel (SAC)$.
		\end{itemchoice}
	} 
\end{ex}

% nguồn: Ok_dungsai_lop11_Lan2
\begin{ex}%[1H4V3-2]%[ma:811214d9]
	Cho hình chóp $S.ABCD$ đáy $ABCD$ là hình bình hành. Gọi $I$, $J$ lần lượt là trọng tâm của tam giác $SAB$ và $SCD$; $E$, $F$  lần lượt là trung điểm của $AB$ và $CD$. Khi đó
	 \choiceTF
		{\True  $\dfrac{SJ}{SF}=\dfrac{2}{3}$}
		{\True  $IJ\parallel (ABCD)$}
		{\True $BC$ song song với mặt phẳng $(SAD)$, $(SEF)$} 
		{$BC$ cắt mặt phẳng  $(AIJ)$}
	\loigiai{ 
		\begin{center}
			\begin{tikzpicture}
					\def\a{4} 
					\path  (0:0) coordinate (A)
						++(0:\a) coordinate  (D)
						 ++(-140:.6*\a) coordinate (C)
						 ($(A)!0.5!(C)+(0,1.1*\a)$) coordinate (S)
						($(A)+(C)-(D)$) coordinate (B) 
						($(A)!.5!(B)$)  coordinate (E)
						($(C)!.5!(D)$) coordinate (F) 
						($(S)!2/3!(E)$) coordinate (I) 
						($(S)!2/3!(F)$)  coordinate (J);
					\draw[dashed]  (B)--(A)--(D)
						(A)--(S) 
						(S)--(E)--(F) 
						(I)--(J); 
					 \draw (B)-- (C)--(D)
						 (B)--(S)
						 (C)--(S)
						 (S)--(D)
						(F)--(S); 
					\foreach \x/\g  in {A/-90,B/-135,C/-45,D/0,S/90,E/-90,F/-45,I/-55,J/5}
					 \fill[black] (\x) circle (1pt)
						($(\g:3mm)+(\x)$) node  {$\x$};
				\end{tikzpicture}
		\end{center}
		\begin{itemchoice}
			\itemch{\bf Đúng.}\\ 
			Do $I$, $J$  lần lượt là trọng tâm của tam giác $SAB$ và $SCD$ nên $\dfrac{SI}{SE}=\dfrac{SJ}{SF}=\dfrac{2}{3} \Rightarrow IJ\parallel EF$.\\
			Mà $EF\subset  (ABCD)\Rightarrow IJ\parallel (ABCD)$.
			\itemch{\bf  Đúng.}
			 \itemch{\bf Đúng.}
			Vì $BC\parallel AD$, $AD\subset  (SAD)\Rightarrow BC\parallel (SAD)$.\\
			Vì $EF$ là  đường trung bình của hình bình hành $ABCD$ nên\\ $BC\parallel EF,EF\subset (SEF)\Rightarrow BC\parallel (SEF)$.\\
			Ta có $IJ\parallel EF,EF\parallel BC\Rightarrow BC\parallel IJ$ mà $IJ\subset (AIJ)\Rightarrow BC\parallel (AIJ)$. 
			 \itemch{\bf Sai.}\\
		\end{itemchoice}
	 }
\end{ex}

% nguồn: Ok_dungsai_lop11_Lan2
\begin{ex}%[1H4V3-2]%[ma:01761c39]
	Cho hình chóp $S.ABCD$ đáy $ABCD$ là hình bình hành. Gọi $I$, $J$ lần lượt là trọng tâm của tam giác $SAB$ và $SCD$; $E$, $F$ lần lượt  là trung điểm của $AB$ và $CD$. Khi đó
	\choiceTF 
		{\True $\dfrac{SJ}{SF}=\dfrac{2}{3}$} 
		{\True  $IJ\parallel(ABCD)$}
		{\True $BC$ song song với mặt phẳng $(SAD)$,  $(SEF)$}
		{$BC$ cắt mặt phẳng $(AIJ)$} 
	\loigiai{ 
		\begin{center}
			\begin{tikzpicture}[scale=1,font=\footnotesize,line join=round,line cap=round,>=stealth]
					\def\canhAD{4}; \def\canhBA{2}; \def\gocBAD{-150}; \def\h{3}; \def\xdinhS{1}; 
					\coordinate (A) at  (0,0);
					\coordinate (B) at  ($(A)+(\gocBAD:\canhBA)$);
					 \coordinate (C) at ($(B)+(0:\canhAD)$);
					\coordinate  (D) at ($(A)+(0:\canhAD)$);
					\coordinate  (S) at ($(A)+(\xdinhS,\h)$);
					 \path
						($(A)!.5!(B)$)  coordinate (E)
						($(C)!.5!(D)$) coordinate (F) 
						($(S)!2/3!(E)$) coordinate (I) 
						($(S)!2/3!(F)$) coordinate (J); 
					 \draw (B)--(S)--(C)--cycle (S)--(D)--(C);
					\draw[dashed]  (A)--(D)
						 (S)--(A)--(B)
						(S)--(E)--(F)--cycle (I)--(J); 
					\foreach  \x/\y in {A/-90,B/-90,C/-90,D/0,S/90,E/-90,F/0,I/-45,J/0}{\fill (\x) circle (1pt)
						($(\x)+(\y:0.3cm)$) node{$\x$};} 
				\end{tikzpicture}
		\end{center}
		\begin{itemchoice}
			 \itemch {\bf Đúng}.\\
			Do $I$, $J$  lần lượt là trọng tâm của tam giác $SAB$ và $SCD$ nên
			 $\dfrac{SI}{SE}=\dfrac{SJ}{SF}=\dfrac{2}{3}$.
			\itemch {\bf Đúng}.\\ 
			Suy ra $IJ \parallel EF$ mà  $EF\subset(ABCD) \Rightarrow IJ \parallel (ABCD)$.
			\itemch {\bf  Đúng}.\\
			Vì $BC\parallel AD$, $AD\subset(SAD)\Rightarrow BC\parallel (SAD)$.\\ 
			 Vì $EF$ là đường trung bình của hình bình hành $ABCD$ nên
			$BC \parallel EF$,  $EF \subset(SEF)$.\\
			Suy ra $BC  \parallel (SEF)$.
			\itemch  {\bf Sai}.\\
			Ta có $IJ\parallel EF$, $EF\parallel BC\Rightarrow BC\parallel  IJ$ mà $IJ\subset(AIJ)\Rightarrow BC\parallel (AIJ)$.\\
		\end{itemchoice}
	 }
\end{ex}

% nguồn: Ok_dungsai_lop11_Lan2
\begin{ex}%[1H4V3-2]%[ma:aa593763]
	Cho hai hình bình hành $ABCD$ và $ABEF$ không cùng  nằm trong một mặt phẳng và có tâm lần lượt là $O$ và $O'$. Gọi $M$, $N$ lần lượt là hai điểm trên các cạnh $AE$, $BD$ sao cho $AM=\dfrac{1}{3} AE$, $BN=\dfrac{1}{3} BD$. Khi đó
	 \choiceTF
		{\True $OO'$ song song  với mặt phẳng $(ADF)$}
		{$OO'$  cắt mặt phẳng $(BCE)$}
		{$\dfrac{BN}{BD}=\dfrac{2}{3}$} 
		{\True  $MN$ song song với mặt phẳng $(CDFE)$}
	\loigiai{ 
		\begin{center}
			\begin{tikzpicture}
					 \def\a{3.5}
					 \path (0:0) coordinate (A)
						++(0:\a) coordinate (B) 
						++(-130:.55*\a) coordinate (C) 
						($(A)+(C)-(B)$) coordinate (D) 
						 ($(A)+(110:.75*\a)$) coordinate (F)
						($(F)+(B)-(A)$) coordinate (E) 
						($2*(C)-(D)$) coordinate (I) 
						($(B)!.5!(C)$)  coordinate (M')
						($(A)!1/3!(E)$) coordinate  (M)
						 ($(B)!1/3!(D)$) coordinate (N)
						(intersection of  A--C and B--D) coordinate (O)
						 (intersection of A--E and B--F) coordinate (O');
					\draw[dashed] (B)--(A)--(D) 
						(A)--(F)--(B) 
						 (A)--(E)
						(B)--(D) 
						 (O)--(O')
						 (A)--(C)
						(A)--(M') 
						(M)--(N); 
					\draw  (C)--(D)--(F)--(E)--cycle (E)--(B)--(C)
						(M')--(I)--(C); 
					\foreach \x/\g  in {A/180,B/0,C/-90,D/180,E/0,F/180,M/170,N/-90,O/-90,O'/70,I/0}
					 \fill[black] (\x) circle (1pt)
						($(\g:4mm)+(\x)$) node  {$\x$};
				\end{tikzpicture}
		\end{center}
		\begin{itemchoice}
			\itemch{\bf Đúng.}\\ 
			Ta có $OO'$ là  đường trung bình của tam giác $BDF$ nên $OO' \parallel DF$.\\
			Mà $DF\subset (ADF)$  suy ra $OO' \parallel (ADF)$.
			\itemch{\bf  Sai.}\\
			Ta có $OO'$ là đường trung bình của tam giác $ACE$ nên  $OO' \parallel CE$.\\
			Mà $CE\subset (BCE)$ suy ra $OO' \parallel (BCE)$. 
			\itemch{\bf  Sai.}\\
			Trong mặt phẳng $(ABCD)$, gọi  $I=AN\cap CD$.\\
			 Do $AB\parallel CD$ nên $\dfrac{AN}{AI}=\dfrac{BN}{BD}=\dfrac{1}{3}$.\\
			\itemch{\bf  Đúng.}\\
			Ta có $\dfrac{AM}{AE}=\dfrac{1}{3} \Rightarrow \dfrac{AN}{AI}=\dfrac{AM}{AE} \Rightarrow MN\parallel IE$, mà  $IE\subset (CDFE)$, suy ra $MN\parallel (CDFE)$.
		\end{itemchoice}
	} 
\end{ex}

% nguồn: Ok_dungsai_lop11_Lan2
\begin{ex}%[1H4V3-2]%[ma:7a081417]
	 Cho hình chóp $S.ABCD$ có đáy $ABCD$ là hình chữ nhật. Gọi $G$ là trọng tâm tam giác $SAD$ và $E$ là điểm trên cạnh $DC$ sao cho $DC=3DE$, $I$ là trung điểm $AD$. Khi đó
	 \choiceTF
		{\True $OI$ song  song với mặt phẳng $(SAB)$}
		{\True $OI$ song  song với mặt phẳng $(SCD)$}
		{$IE$ song song với  $AC$}
		{\True $GE\parallel (SBC)$} 
	\loigiai{ 
		\begin{center}
			\begin{tikzpicture}
					 \def\a{4}
					\path (0:0) coordinate (A) 
						++(0:\a) coordinate (D) 
						++(-120:.5*\a) coordinate (C) 
						($(A)+(C)-(D)$) coordinate (B) 
						($(A)+(78:\a)$)  coordinate (S)
						($(A)!.5!(D)$)  coordinate (I)
						 ($(B)!.5!(C)$) coordinate (K)
						 ($(S)!2/3!(I)$) coordinate (G)
						 ($(S)!2/3!(K)$) coordinate (G')
						($(D)!1/3!(C)$)  coordinate (E)
						(intersection of A--C and B--D)  coordinate (O);
					 \draw[dashed] (B)--(A)--(D)
						 (A)--(S)
						 (A)--(C)
						(B)--(D) 
						 (S)--(I)--(K)
						 (E)--(G)--(G')--(C);
					\draw  (B)-- (C)--(D)
						 (B)--(S)
						 (C)--(S)
						 (D)--(S)
						 (S)--(K);
					\foreach \x/\g in  {A/180,B/-135,C/-45,D/0,S/90,O/-80,I/-60,K/-90,G/170,G'/165,E/-45}
					\fill[black]  (\x) circle (1pt)
						($(\g:3mm)+(\x)$) node  {$\x$};
				\end{tikzpicture}
		\end{center}
		\begin{itemchoice}
			 \itemch{\bf Đúng.}\\
			Ta có $\heva{&OI\not \subset (SAB),AB\subset  (SAB)\\&OI\parallel AB}\Rightarrow OI\parallel (SAB)$.
			\itemch{\bf  Đúng.}\\
			Ta có  $\heva{&OI\not \subset (SCD),CD\subset (SCD)\\&OI\parallel CD}\Rightarrow OI\parallel (SCD)$.
			\itemch{\bf  Sai.}\\
			Vì  $\dfrac{DI}{DA}=\dfrac{1}{2} \ne \dfrac{1}{3}=\dfrac{DE}{DC}$ nên $IE$ không song song với $AC$.
			\itemch{\bf  Đúng.}\\
			Trong hình chữ nhật $ABCD$, gọi $P=IE\cap BC\Rightarrow P=IE\cap (SBC)$.\\ 
			Gọi $K$ là  trung điểm của $BC$, $G'$ là trọng tâm tam giác $SBC$.\\
			Khi đó $\dfrac{SG'}{SK}=\dfrac{SG}{SI}=\dfrac{G'  G}{KI}=\dfrac{2}{3}$, suy ra $G'G\parallel KI\parallel CE$ và $\Rightarrow G'G=\dfrac{2}{3} KI=\dfrac{2}{3} CD=CE$.\\
			Do dó tứ giác $G'GEC$ là hình bình hành, suy ra $CG'\parallel CE\Rightarrow CG\parallel  (SBC)$.
		\end{itemchoice}
	 }
\end{ex}

% nguồn: Ok_dungsai_lop11_Lan2
\begin{ex}%[1H4V3-2]%[ma:6e563c3e]
	Cho hai hình bình hành $ABCD$ và $ABEF$ không cùng nằm trong  một mặt phẳng và có tâm lần lượt là $O$ và $O'$. Gọi $M$, $N$ lần lượt là hai điểm trên các cạnh $AE$, $BD$ sao cho $AM=\dfrac{1}{3} AE$, $BN=\dfrac{1}{3} BD$. Khi đó
	\begin{center}
		\begin{tikzpicture}[scale=1,font=\footnotesize,line join=round,line cap=round,>=stealth]
				\def\canhAB{4}; \def\canhDA{2}; \def\gocDAB{-130}; \def\h{3}; \def\xdinhF{-1}; 
				\coordinate  (A) at (0,0);
				\coordinate (D) at ($(A)+(\gocDAB:\canhDA)$); 
				\coordinate  (C) at ($(D)+(0:\canhAB)$);
				 \coordinate (B) at ($(A)+(0:\canhAB)$);
				\coordinate (F) at ($(A)+(\xdinhF,\h)$); 
				\coordinate  (O) at (intersection of A--C and D--B);
				 \path
					($(B)+(F)-(A)$) coordinate  (E)
					 ($(A)!.5!(E)$) coordinate (O')
					($(A)!1/3!(E)$)  coordinate (M)
					($(B)!1/3!(D)$) coordinate (N) 
					(intersection of A--N and C--B) coordinate (1) 
					(intersection of A--N and  D--C) coordinate (I);
				\draw (D)--(F)--(E)--(C)--cycle (E)--(B)--(C) 
					 (1)--(I)--(C);
				\draw[dashed] (A)--(B) 
					(A)--(C) 
					 (D)--(B)
					 (F)--(A)--(D)
					(A)--(E) 
					(M)--(N) 
					(O)--(O') 
					 (A)--(1);
				\foreach \x/\y in  {A/180,D/-90,C/-90,B/0,F/90,O/225,E/90,O'/90,M/90,N/90,I/0}{\fill (\x) circle (1pt)
					($(\x)+(\y:0.3cm)$)  node{$\x$};}
			\end{tikzpicture}
	\end{center}
	 \choiceTF
		{\True $OO'$ song song với mặt phẳng $(ADF)$} 
		{$OO'$ cắt  mặt phẳng $(BCE)$}
		{$\dfrac{BN}{BD}=\dfrac{2}{3}$} 
		{\True $MN$ song song với mặt phẳng $(CDFE)$} 
	\loigiai{ 
		\begin{itemchoice}
			\itemch  {\bf Đúng}.\\
			Ta có $OO'$ là đường trung bình của tam giác $BDF$ nên $OO'\parallel  DF$. \\
			Mà $DF\subset(ADF)$  suy ra $OO'\parallel(ADF)$.
			 \itemch {\bf Sai}.\\
			Tương tự, $OO'$ là đường trung bình của tam giác $ACE$ nên $OO'\parallel CE$.  \\
			Mà $CE\subset(BCE)$  suy ra $OO'\parallel(BCE)$.
			 \itemch {\bf Sai}.\\
			Trong mặt phẳng $(ABCD)$, gọi  $I=AN\cap CD$.\\
			Do  $AB\parallel CD$ nên $\dfrac{AN}{AI}=\dfrac{BN}{BD}=\dfrac{1}{3}$.
			\itemch {\bf  Đúng}.\\
			Mặt khác: $\dfrac{AM}{AE}=\dfrac{1}{3}\Rightarrow\dfrac{AN}{AI}=\dfrac{AM}{AE}\Rightarrow MN\parallel  IE$, mà $IE\subset(CDFE)$, suy ra $MN\parallel(CDFE)$.
		\end{itemchoice}
	} 
\end{ex}


%%%%%%%%%%%%%%%% TRẢ LỜI NGẮN %%%%%%%%%%%%%%%%

% nguồn: Ok_TraLoiNgan_Lop11_Lan1
\begin{ex}%[1H4H3-2]%[1H4V3-4]%[1H4H2-2]%[ma:1daf7369]
	Cho tứ  diện $ABCD$ có $M$ và $N$ lần lượt là trung điểm $AB$ và $BC$. $P$ là điểm thuộc cạnh $CD$ sao cho $PD=2PC$. Gọi $Q$ là giao điểm của đường thẳng $AD$ và mặt phẳng $\left( MNP\right) $. Tỉ số $\dfrac {AQ}{AD}$ bằng
	\shortans{0,33} 
	 \loigiai{
	\immini{ 
		Ta có $ M, N$ lần lượt là trung điểm $AB,  BC$ nên $MN$ là đường trung bình tam giác $ABC$. Do đó $MN\parallel AC$.\\
		Ta có \mbox{$\heva{&MN\parallel AC\\&AC\subset \left( ACD\right)\\&MN \not\subset \left( ACD\right) }\Rightarrow MN \parallel\left(  ACD\right) $} .
		 }{
		\begin{tikzpicture}[scale=0.8,font=\footnotesize,line join=round,line cap=round,>=stealth]
				\coordinate (B) at (-4,0); \draw[fill=black] (B)  circle (1pt);
				\coordinate  (D) at (3,0); \draw[fill=black] (D) circle (1pt);
				\coordinate (C) at (0,-2); \draw[fill=black] (C) circle  (1pt);
				\coordinate  (A) at (0.5,6); \draw[fill=black] (A) circle (1pt);
				\coordinate (M) at (-1.75,3);  \draw[fill=black] (M) circle (1pt);
				\coordinate (N)  at (-2,-1); \draw[fill=black] (N) circle (1pt);
				\coordinate  (P) at (1.08,-1.28); \draw[fill=black] (P) circle (1pt);
				\coordinate (Q) at (1.36,3.95); \draw[fill=black] (Q) circle  (1pt);
				 \draw (B) node[left]{$B$}--(C) node[below]{$C$}--(D) node[right]{$D$}--(A) node[above]{$A$}--(C);
				\draw (B)--(A); 
				\draw[dashed]  (B)--(D);
				 \draw (M) node[left]{$M$}--(N) node[below]{$N$};
				 \draw[dashed] (N)--(P) node[right]{$P$};
				\draw[dashed]  (M)--(Q) node[right]{$Q$};
				 \draw (P)--(Q);
			\end{tikzpicture}
		} 
		$\heva{&AC\parallel MN\\&AC\subset \left( ACD\right) \\&MN\subset\left( MNP\right) \\& P \in \left( ACD\right) \cap \left( MNP\right) }\Rightarrow  \left( ACD\right) \cap\left( MNP\right) =Px\parallel AC\parallel MN$.\\
		Gọi $Q=Px \cap AD\Rightarrow  Q=AD\cap \left( MNP\right) $.\\
		Trong tam  giác $ACD$ có $PQ\parallel AC$. Nên $\dfrac{AQ}{AD}=\dfrac{CP}{CD}=\dfrac{1}{3}$.
	 }
\end{ex}

% nguồn: Ok_TraLoiNgan_Lop11_Lan1
\begin{ex}%[1H4H3-2]%[ma:a355d9b0]
	Cho tứ diện $ABCD$. Gọi $G_1$ và $G_2$ lần lượt là trọng tâm các tam giác $BCD$ và  $ACD$. Tính tỉ số $\dfrac{AB}{G_1 G_2}$ (làm tròn kết quả đến hàng phần trăm).
	 \shortans{3}
	\loigiai{ 
	\immini{Ta  có $G_1$ và $G_2$ lần lượt là trọng tâm các $\triangle BCD$ và $\triangle ACD$.\\
		Suy ra $BG_1$, $AG_2$ và $CD$ đồng qui tại $M$  (là trung điểm của $CD$).\\
		Vì $G_1 G_2 \parallel  AB$ nên $G_1 G_2 \parallel \left(ABD\right)$ và $G_1 G_2 \parallel \left(ABC\right)$.\\
		Lại có $G_1 G_2=\dfrac{1}{3} AB$ hay $AB=3G_1 G_2$.} 
		{\begin{tikzpicture}[font=\footnotesize,scale=.7]
		 \def\a{4}
		 \path (0:0) coordinate (B)
			 ++(0:\a) coordinate (D)
			 ++(-145:.68*\a) coordinate (C)
			($(B)+(70:.8*\a)$) coordinate (A) 
			($(C)!.5!(D)$) coordinate  (M)
			 ($(M)!1/3!(B)$) coordinate (G_1)
			 ($(M)!1/3!(A)$) coordinate (G_2);
		 \draw[dashed] (B)--(D)
			 (M)--(B)
			(G_1)--(G_2); 
		\draw  (B)--(A)--(D)
			 (A)--(C)
			(B)--(C)--(D) 
			(A)--(M); 
		\foreach \x/\g  in {A/90,B/180,C/-45,D/0,M/-45,G_1/-45,G_2/5}
		\fill[black] (\x) circle  (1pt)
			 ($(\g:4mm)+(\x)$) node {$\x$};
	\end{tikzpicture}
		 }
	} 
\end{ex}

% nguồn: Ok_TraLoiNgan_Lop11_Lan2
\begin{ex}%[1H4H3-2]%[ma:06295467]
	Cho  hình chóp $S.ABCD$ đáy là hình bình hành tâm $O$. Điểm $M$ thuộc cạnh $SB$. Biết $OM\parallel (SCD)$. Tính tỉ số của $\dfrac{SM}{MB}$.
	\shortans{1} 
	 \loigiai{
		\begin{center}
			\begin{tikzpicture}[scale=1,font=\footnotesize,line join=round,line cap=round,>=stealth]
					\path 
						 (0,0) coordinate (A)
						(5,0)  coordinate (D)
						 (-2,-1.5) coordinate (B)
						 ($(B)+(D)-(A)$) coordinate (C)
						($(B)!.5!(D)$)  coordinate (O)
						(1,3.5) coordinate  (S)
						($(S)!.5!(B)$)  coordinate (M);
					\draw  (S)--(B)--(C)--(D)--(S)--(C);
					\draw [dashed]  (C)--(A)--(D)--(B)--(A)--(S)
						 (O)--(M);
					\foreach \d/\g  in {A/135,B/-135,C/-45,D/45,S/90,M/135,O/-90}
					\fill (\d)  circle (1pt) + (\g:.3) node{$\d$};
				\end{tikzpicture}
		\end{center}
		Trong  mặt phẳng ($SBD$) xét hai đường thẳng $OM$ và $SD$.\\
		Vì $SD\subset(SCD)$ và $OM  \parallel (SCD)$ nên $OM\cap SD=\varnothing$ hay $OM \parallel SD$.\\
		Suy ra $OM$ là đường trung bình của tam giác $SBD$ nên $M$ là trung  điểm $SB$ hay $\dfrac{SM}{MB}=1$.
	} 
\end{ex}

% nguồn: Ok_TraLoiNgan_Lop11_Lan1
\begin{ex}%[1H4V3-2]%[ma:c18f6da8]
	Cho tứ diện $ABCD,G$ là trọng tâm của $\triangle ABD$ và $M$ là một điểm trên cạnh $BC$ sao cho $MB=2MC$. Biết độ dài $MG=2$ cm và $AD=3$ cm, tính diện tích tam giác $ACD$. 
	\shortans{4,50} 
	 \loigiai{
	\immini{Gọi $K$ là trung điểm  đoạn $AD$. \\
		 Suy ra $\dfrac{BG}{BK}=\dfrac{2}{3}$ ($G$ là trọng tâm của tam giác $ABD$).\\
		Mặt  khác $\dfrac{BM}{BC}=\dfrac{2}{3}$.\\
		Suy ra $\triangle BCK$ có $\dfrac{BM}{BC}=\dfrac{BG}{BK}$ (góc $B$  chung) nên $MG\parallel CK$.\\
		và $CK=\dfrac{3}{2}MG=\dfrac{3}{2}\cdot 2=3$  cm.\\
		Vậy  $S_{ACD}=\dfrac{1}{2}\cdot CK\cdot AD=\dfrac{1}{2}\cdot 3\cdot 3=4{,}5$ (cm$^2$).}
		{\begin{tikzpicture}[font=\footnotesize,scale=.7]
		\def\a{4} 
		 \path (0:0) coordinate (A)
			 ++(0:\a) coordinate (B)
			 ++(-145:.7*\a) coordinate (D)
			($(A)+(60:\a)$) coordinate  (C)
			($(A)!.5!(D)$) coordinate  (K)
			($(B)!2/3!(K)$) coordinate (G) 
			 ($(B)!2/3!(C)$) coordinate (M);
		\draw[dashed]  (K)--(B)--(A)
			 (M)--(G);
		\draw  (A)--(D)--(B)
			 (D)--(C)--(K)
			(A)--(C)--(B); 
		\foreach \x/\g in {A/180,B/0,C/90,D/-45,K/-135,G/-45,M/45} 
		\fill[black]  (\x) circle (1pt)
			 ($(\g:3mm)+(\x)$) node {$\x$};
	\end{tikzpicture}
		 }
	} 
\end{ex}

% nguồn: Ok_TraLoiNgan_Lop11_Lan1
\begin{ex}%[1H4V3-2]%[ma:c6d10c30]
	Cho hình chóp $S.ABCD$ có đáy $ABCD$ là hình bình hành tâm $O$, $M$ là trung điểm đoạn $SB$, $G$ là trọng tâm tam giác $SAD$. Gọi $J$ là giao  điểm của $AD$ với $(OMG)$. Tính $\dfrac{JD}{AD}$(làm tròn đến hàng phần trăm).
	 \shortans{0,33}
	\loigiai{\immini[thm]{Ta có $MO$ là đường trung bình  của tam giác $SBD$ suy ra $MO\parallel SD$.\\
		Ta có $\left\{ \begin{aligned}&MO\parallel SD\\ &G\in (OMG)\cap (SAD)
	\end{aligned}\right.\\
	\Rightarrow (OMG)\cap (SAD)=Gx\parallel SD\parallel MO\\ \Rightarrow Gx\cap  AD=J$.\\
	Ta  có $J\in Gx\subset (OMG)\Rightarrow J=AD\cap (OMG)$.\\
	 Gọi $E$ là trung điểm $SD$ với $G$ là trọng tâm ta có $\dfrac{AG}{AE}=\dfrac{2}{3}$.\\
	Do $GJ\parallel  MO\parallel SD$, áp dụng định lý Ta-lét trong tam giác $AED$ ta có $\dfrac{GE}{AE}=\dfrac{JD}{AD}=\dfrac{1}{3} \approx 0{,}33$.
	}{\begin{tikzpicture}[scale=1,font=\footnotesize,line join=round,line cap=round,>=stealth]
		\coordinate (A)  at (0,0);
		\coordinate (B) at (-1.2,-1.5); 
		 \coordinate (D) at (4,0);
		\coordinate (C)  at ($(B)+(D)-(A)$);
		\coordinate (S) at ($(A)+(1,6)$); 
		\coordinate (M) at ($(S)!1/2!(B)$); 
		\coordinate (E)  at ($(S)!1/2!(D)$);
		\coordinate  (G) at ($(A)!2/3!(E)$);
		\draw[dashed,name path=ac] (A)--(C); 
		\draw[dashed,name  path=bd] (B)--(D);
		\path[name intersections={of=ac  and bd,by={O}}];
		\coordinate (a)  at ($(G)+(O)-(M)$);
		\coordinate (a1)  at ($(a)!1.2!(G)$);
		\draw[dashed,name  path=aa1] (a1)--(a);
		 \draw[dashed,name path=ad] (A)--(D);
		\path[name intersections={of=aa1 and  ad,by={J}}];
		 \draw (S)--(B)
			(S)--(C) 
			 (S)--(D)
			 (B)--(C)
			 (C)--(D);
		\draw[dashed] (S)--(A) 
			(A)--(C) 
			 (B)--(D)
			 (A)--(B)
			 (A)--(D)
			 (M)--(O)
			(G)--(J) 
			(A)--(E) 
			 (M)--(G)
			(O)--(G); 
		\foreach \t/\g  in {A/135,S/135,B/135,M/135,D/-45,C/-45,E/45,J/45,O/-90,G/90}{
			\draw[fill=black]  (\t) circle (1pt) node[shift={(\g:7pt)},font=\scriptsize]{$ \t $};}
	\end{tikzpicture}
	 }
	 }
\end{ex}


%%%%%%%%%%%%%%%% TỰ LUẬN %%%%%%%%%%%%%%%%

% nguồn: Ok_TuanLuan_Toan11_Lan2
\begin{ex}%[1H4H3-2]%[ma:80db6b32]
	Cho hình chóp $S.ABCD$ có đáy là hình thang, đáy nhỏ $AB=a$, đáy lớn $CD=2a$. Gọi $E$ là trung điểm  của $SC$. Chứng minh rằng $BE\parallel (SAD)$.
	\loigiai{ 
		\begin{center}
			\begin{tikzpicture}[line join=round,line cap=round]
					\coordinate (S) at (0,4); 
					\coordinate (D) at  (-2,0);
					\coordinate (F) at  (0,-3);
					\coordinate  (C) at (3,0);
					\coordinate (A) at ($(D)!0.5!(F)$); 
					 \coordinate (B) at ($(C)!0.5!(F)$);
					\coordinate (E) at ($(S)!0.5!(C)$); 
					\draw (S)--(D) 
						(S)--(F) 
						(S)--(C) 
						(F)--(C) 
						(D)--(F) 
						 (E)--(B)--(S)--(A);
					 \draw[dashed](A)--(B)
						 (D)--(C);
					\foreach  \i/\g in {S/90,D/180,F/-90,C/0,A/180,B/0,E/0}{\draw[fill](\i) circle (1pt)
						($(\i)+(\g:3mm)$) node[scale=1]{$\i$};} 
				\end{tikzpicture}
		\end{center}
		Gọi $F$  là giao điểm của $BC$ và $AD$.\\
		Ta có $AB\parallel CD$  và $AB=\dfrac{1}{2}CD$, suy ra $AB$ là đường trung bình của tam giác $CDF$.\\
		Do đó $B$ là trung điểm  của $FC$.\\
		Suy ra $BE$ là đường  trung bình của tam giác $SCF$ hay $BE\parallel SF$.\\
		Mà  $SF\subset(SAD)$ nên $BE\parallel (SAD)$.
	 }
\end{ex}

% nguồn: Ok_TuanLuan_Toan11_Lan2
\begin{ex}%[1H4H3-2]%[ma:7934bda6]
	Cho hình chóp $S.ABCD$ có đáy  $ABCD$ là hình bình hành. Chứng minh rằng \break $AB \parallel ( SCD)$.
	\loigiai{ 
		\begin{center}
			\begin{tikzpicture}[scale=1,font=\footnotesize,line join=round,line cap=round]
					\foreach  \x\y\t in {0/3/S,0/0/A,-1.3/-1.6/B,2.5/-1.6/C}
					\coordinate  (\t) at (\x,\y);
					 \coordinate (D) at ($(A)+(C)-(B)$);
					\draw  (S)--(C)--(B)--(S)--(D)--(C);
					 \draw[dashed](S)--(A)
						 (B)--(A)--(D);
					\path pic[draw,angle radius=6]{right  angle=S--A--B}
						 pic[draw,angle radius=6]{right angle=S--A--D};
					\foreach \t/\g  in {S/90,A/160,B/-130,C/-60,D/0}
					 \draw[fill] (\t) circle (1pt)
						node[shift={(\g:7pt)}]{$\t$}; 
				\end{tikzpicture}
		\end{center}
		Nếu đường thẳng $AB$ và mặt phẳng $( SCD)$ có điểm chung là  $M$ thì điểm $M$ nằm trên cả hai mặt phẳng $( ABCD)$ và $( SCD)$ suy ra điểm $M$ nằm trên $CD$. \\
		Do đó  $M$ là điểm chung của hai đường thẳng $AB$ và $CD$. \\
		Điều này không xảy  ra vì $AB\parallel CD$. \\
		Vậy $AB \parallel  ( SCD)$.
	 }
\end{ex}

% nguồn: Ok_TuanLuan_Toan11_Lan2
\begin{ex}%[1H4H3-2]%[ma:476c1187]
	Cho tứ diện $ABCD$. Gọi $I$, $J$ lần lượt là trọng  tâm của hai tam giác $ABC$, $ACD$. Chứng minh rằng $IJ\parallel (BCD)$.
	 \loigiai{
		Gọi  $M$, $N$ lần lượt là trung điểm của $BC$, $CD$.
		\begin{center}
			\begin{tikzpicture}[line join=round,line cap=round]
					\coordinate (A) at (0,4); 
					\coordinate (D)  at (-2,0);
					 \coordinate (C) at (0,-3);
					\coordinate (B)  at (3,0);
					\coordinate (M)  at ($(C)!0.5!(B)$);
					\coordinate  (N) at ($(C)!0.5!(D)$);
					 \coordinate (I) at ($(A)!2/3!(M)$);
					\coordinate (J) at ($(A)!2/3!(N)$); 
					\draw  (A)--(D)
						 (A)--(C)
						(A)--(B) 
						(B)--(C) 
						 (D)--(C)
						 (N)--(A)--(M);
					 \draw[dashed](D)--(B)
						 (M)--(N)
						 (I)--(J);
					\foreach \i/\g  in {A/90,D/180,B/0,C/-90,M/0,N/180,I/0,J/180}{\draw[fill](\i) circle (1.0pt)
						 ($(\i)+(\g:3mm)$) node[scale=1]{$\i$};}
				\end{tikzpicture}
		\end{center}
		Suy ra $MN$ là đường trung bình của tam giác $BCD$. Suy ra $MN\parallel  BD$. (1)\\
		Mặt khác, $I$, $J$ lần lượt là trọng  tâm của tam giác $ABC$ và $ACD$ nên $\dfrac{AI}{AM}=\dfrac{AJ}{AN}=\dfrac{2}{3}$.\\
		Theo định lí Thalès đảo trong tam  giác $AMN$, ta có $IJ\parallel MN$. (2)\\
		Từ  (1) và (2) suy ra $IJ\parallel BD$.\\
	Mà $BD\subset (BCD)$ nên $IJ\parallel (BCD)$.} 
\end{ex}

% nguồn: Ok_TuanLuan_Toan11_Lan2
\begin{ex}%[1H4H3-2]%[ma:834a1cad]
	Cho tứ diện $ABCD$. Gọi $G_1$ và  $G_2$ lần lượt là trọng tâm của hai tam giác $ABD$ và $ACD$ . Chứng minh $G_1 G_2$ song song với các mặt phẳng $(ABC)$ và $(BCD)$.
	\loigiai{ 
		\begin{center}
			\begin{tikzpicture}[line join=round,line cap=round]
					\coordinate (A) at (2,4); 
					\coordinate (D) at (-2,0); 
					\coordinate (C) at (0,-3); 
					\coordinate  (B) at (3,0);
					\coordinate  (M) at ($(C)!0.5!(B)$);
					\coordinate (N) at ($(B)!0.5!(D)$); 
					\coordinate (G_1)  at ($(A)!2/3!(M)$);
					 \coordinate (G_2) at ($(A)!2/3!(N)$);
					 \draw (A)--(D)
						(A)--(C) 
						 (A)--(B)
						 (B)--(C)
						(D)--(C) 
						(A)--(M); 
					 \draw[dashed](D)--(B)
						 (M)--(N)
						 (N)--(A)
						 (G_1)--(G_2);
					 \foreach \i/\g in {A/90,D/180,B/0,C/-90,M/-40,N/120,G_1/0,G_2/180}{\draw[fill](\i) circle (1.0pt)
						 ($(\i)+(\g:3mm)$) node[scale=1]{$\i$};}
				\end{tikzpicture}
		\end{center}
		Gọi $M$, $N$ lần lượt là trung điểm của $DB$, $DC$.\\ 
		Ta có $MN$ là đường trung bình của tam giác $DBC$, suy ra $MN\parallel  BC$.\\
		Trong tam giác $AMN$, ta  có $\dfrac{AG_1}{AM}=\dfrac{AG_2}{AN}=\dfrac{2}{3}$.\\
		Theo định lí Thalès đảo trong tam giác $AMN$, ta có $G_1G_2\parallel MN$.\\ 
		Suy ra $G_1G_2\parallel MN\parallel BC$, suy ra $G_1G_2$  song song với các mặt phẳng $(ABC)$ và $(BCD)$.
	 }
\end{ex}

% nguồn: Ok_TuanLuan_Toan11_Lan2
\begin{ex}%[1H4H3-2]%[ma:11fc0a06]
	Cho hai hình bình hành $ABCD$ và $A B E  F$ không cùng nằm trong một mặt phẳng có tâm lần lượt là $O$ và $O'$.
	\begin{enumerate}
		\item Chứng minh $OO'$ song song với các mặt phẳng $(ADF)$ và $(BCE)$.\\ 
		\item Gọi $M$, $N$ lần lượt là hai điểm thuộc hai cạnh $AF$, $AD$  sao cho $AM=\dfrac{1}{3}AF$, $AN=\dfrac{1}{3}AD$. Chứng minh $MN\parallel (DCEF)$.
	\end{enumerate}
	\loigiai{ 
		\begin{center}
			\begin{tikzpicture}[line join=round,line cap=round]
					\coordinate  (A) at (2,3);
					\coordinate (B) at (6,3); 
					\coordinate  (D) at (0,0);
					\coordinate (F) at  (0,6);
					\coordinate (C)  at ($(B)+(D)-(A)$);
					\coordinate (E) at  ($(B)+(F)-(A)$);
					\coordinate (O') at ($(A)!0.5!(E)$); 
					\coordinate (O) at ($(A)!0.5!(C)$); 
					 \coordinate (M) at ($(A)!1/3!(F)$);
					\coordinate (N) at  ($(A)!1/3!(D)$);
					\draw (E)--(F)--(D)--(C)--(B)--(E)--(C); 
					\draw[dashed] (E)--(A)--(B) 
						 (A)--(F)--(B)
						(C)--(A)--(D)--(B) 
						 (O)--(O')
						 (M)--(N);
					\foreach \i/\g in  {A/180,B/0,C/-90,D/-90,E/90,F/90,M/0,N/0,O/-110,O'/100}{\draw[fill](\i) circle (1pt)
						($(\i)+(\g:3mm)$)  node[scale=1]{$\i$};}
				\end{tikzpicture}
		\end{center}
		\begin{enumerate}
			\item Ta có $OO' \parallel DF\parallel CE$,  suy ra $OO'$ song song với các mặt phẳng $(ADF)$ và $(BCE)$.\\
			\item Ta có $\dfrac{AM}{AF}=\dfrac{AN}{AD}$, suy ra $MN\parallel DF$ , suy ra $MN\parallel  (DCEF)$.
		\end{enumerate}
	 }
\end{ex}

% nguồn: Ok_TuanLuan_Toan11_Lan2
\begin{ex}%[1H4H3-2]%[ma:b05c8e81]
	Cho hình chóp $S.ABCD$ có đáy $ABCD$ là hình bình hành và $M$, $N$, $E$ lần lượt là trung điểm  của các đoạn thẳng $AB$, $CD$, $SA$ (hình dưới). Chứng minh rằng
	\begin{enumerate}
		\item $MN$ song  song với hai mặt phẳng $( SBC)$ và $( SAD)$;
		\item $SB$ và $SC$ song song với mặt phẳng  $( MNE)$.
	\end{enumerate}
	\loigiai{ 
		\begin{center}
			\begin{tikzpicture}[scale=1,font=\footnotesize,line join=round,line cap=round]
					 \foreach \x\y\t in {0/0/A,-1.7/-1.6/B,2.5/-1.6/C,0.4/2.8/S}
					\coordinate  (\t) at (\x,\y);
					\coordinate (D)  at ($(A)+(C)-(B)$);
					\coordinate  (M) at ($(A)!1/2!(B)$);
					\coordinate (N) at ($(C)!1/2!(D)$); 
					\coordinate (E)  at ($(S)!1/2!(A)$);
					 \draw (S)--(B)--(C)--(S)--(D)--(C);
					\draw[dashed](B)--(A)--(D) 
						 (S)--(A)
						(E)--(M)--(N)--cycle; 
					\foreach \t/\g in  {S/90,A/-60,B/-140,C/-45,D/0,M/-70,N/-50,E/40}
					\draw[fill] (\t) circle (1.0pt) 
						 node[shift={(\g:9pt)}]{$\t$};
				\end{tikzpicture}
		\end{center}
		\begin{enumerate}
			\item Do $M$, $N$ lần lượt là trung điểm của $AB,CD$ suy ra  $MN\parallel BC,MN\parallel AD$.\\
			Ta có $MN$ không nằm trong $( SBC)$ và  song song với $BC$ nằm trong $( SBC)$.\\
			Suy ra  $MN\parallel ( SBC)$ và chứng minh tương tự $MN\parallel ( SAD)$.
			\item Do $M$, $E$ lần lượt là trung điểm  của $AB$, $SA$ suy ra $ME$ là đường trung bình tam giác $SAB$ hay $MN\parallel SB$.\\
			Ta có $SB$ không nằm trong $( MNE)$ và song song với $ME$ nằm trong  $( MNE) \Rightarrow SB\parallel ( MNE)$. \\
			Gọi $O$ là giao điểm của $AC$ và $BD$ do $ABCD$ là hình  bình hành nên $O$ là trung điểm của $AC$ mà $E$ là trung điểm của $SA$ nên $OE$ là đường trung bình của tam giác $SAC \Rightarrow OE \parallel SC$.\\
			Ta  có $SC$ không nằm trong $( MNE)$ và song song với $OE$ nằm trong $( MNE) \Rightarrow SC\parallel ( MNE)$.
		\end{enumerate}
	 }
\end{ex}

% nguồn: Ok_TuanLuan_Toan11_Lan2
\begin{ex}%[1H4H3-2]%[ma:6617dcb4]
	Cho hình chóp $S.ABCD$ có đáy $ABCD$ là hình bình hành. Gọi $M$, $N$ lần lượt là trung điểm của  hai cạnh $AB$ và $CD$, $P$ là trung điểm của $SA$. Chứng minh:
	\begin{enumerate}
		\item $MN$ song song với  các mặt phẳng $(SBC)$ và $(SAD)$;\\
		\item  $SB$ song song với $(MNP)$;\\
		\item $SC$ song song  với $(MNP)$.\\
		 \item Gọi $G_1$ và $G_2$ theo thứ tự là trọng tâm của hai tam giác $A B C$ và $S B C$ . Chứng minh $G_1 G_2$ song song với $(SAD)$.
	\end{enumerate}
	 \loigiai{
		\begin{center}
			\begin{tikzpicture}[line join=round,line cap=round,scale=0.7]
					\coordinate  (A) at (0,0);
					\coordinate  (B) at (2,-2);
					\coordinate (D) at (5,0); 
					\coordinate (C)  at ($(B)+(D)-(A)$);
					\coordinate  (N) at ($(C)!0.5!(D)$);
					\coordinate  (S) at ($(A)+(8,6)$);
					\coordinate (M)  at ($(A)!0.5!(B)$);
					 \coordinate (P) at ($(S)!0.5!(A)$);
					\coordinate (d1) at ($(S)+(B)-(A)$); 
					 \coordinate (d2) at ($(S)+(A)-(B)$);
					\coordinate  (E) at (intersection of M--P and d1--d2);
					\coordinate  (G_1) at ($(C)!2/3!(M)$);
					\coordinate (g2) at  ($(B)!0.5!(C)$);
					\coordinate (G_2) at ($(S)!2/3!(g2)$); 
					\coordinate (I) at (intersection of S--G_2 and  B--C);
					\coordinate  (e) at (intersection of E--N and S--A);
					 \draw (S)--(A)
						 (S)--(B)
						 (S)--(C)
						 (A)--(B)
						(B)--(C) 
						(M)--(E) 
						 (S)--(I)
						(E)--(e); 
					\draw[dashed] (C)--(A)--(D) 
						 (C)--(D)
						(S)--(D) 
						(M)--(N) 
						(A)--(I) 
						 (G_1)--(G_2)
						(e)--(N); 
					 \draw (d2)--(d1) node[above]{$d$};
					\foreach \i/\g in  {S/90,A/180,B/-90,C/-90,D/120,M/180,N/0,P/180,E/50,G_1/-90,G_2/-50,I/-90}{\draw[fill](\i) circle (1pt)
						 ($(\i)+(\g:3mm)$) node[scale=0.7]{$\i$};}
				\end{tikzpicture}
		\end{center}
		\begin{enumerate}
			\item Ta có $BC\subset(SBC)$ và $MN\parallel BC$, suy  ra $MN\parallel (SBC)$; $AD\subset(SAD)$ và $MN\parallel AD$, suy ra $MN\parallel (SAD)$.
			\item Trong  tam giác $SAB$, có $PM$ là đường trung bình, suy ra $SB\parallel MP$, suy ra $S B\parallel (M N P)$.
			\item Trong mặt phẳng $(SAB)$ vẽ đường thẳng $d$ đi qua $S$ và $d\parallel  AB$.\\
			Gọi $E$ là giao điểm của $MP$ và $d$, ta có $MBSE$ là hình bình hành, suy ra $S E\parallel M B$ và $S E=M B$, suy ra $SE\parallel CN$ và $SE=CN$, suy ra $SC\parallel NE$.\\ 
			 Ta lại có $NE\subset(MNP)$, suy ra $SC\parallel (MNP)$.
			\item Trong mặt phẳng $(SBC)$, gọi $I$ là giao  điểm của $SG_2$ và $BC$.
			Trong tam giác $SIA$, ta có $\dfrac{IG_1}{I A}=\dfrac{IG_2}{IS}=\dfrac{1}{3}$ , theo định lí Thalès đảo suy  ra $G_1G_2\parallel SA$, suy ra $G_1G_2\parallel (SAD)$.
		\end{enumerate}
	 }
\end{ex}

% nguồn: Ok_TuanLuan_Toan11_Lan2
\begin{ex}%[1H4H3-2]%[ma:7572af23]
	Cho hình chóp $S.ABCD$ có đáy $ABCD$ là một hình bình hành. Gọi $G$ là trọng tâm của tam giác $SAB$, $I$ là trung điểm của $AB$ và $M$ là điểm thuộc cạnh $AD$ sao cho $AM=\dfrac{1}{3}AD$. Đường thẳng đi qua $M$ và song song với $AB$  cắt $CI$ tại $N$. Chứng minh:
	\begin{enumerate}
		\item $NG\parallel  (SCD)$;
		\item  $MG\parallel (SCD)$.
	\end{enumerate}
	 \loigiai{
		\begin{center}
			\begin{center}
				\begin{tikzpicture}[line join=round,line cap=round]
						\coordinate (A)  at (0,0);
						\coordinate  (B) at (2,-2);
						\coordinate (D) at  (5,0);
						\coordinate (C) at  ($(B)+(D)-(A)$);
						\coordinate (M) at  ($(A)!1/3!(D)$);
						\coordinate (S) at  ($(A)+(8,6)$);
						 \coordinate (I) at ($(A)!0.5!(B)$);
						\coordinate (m) at  ($(M)+(B)-(A)$);
						\coordinate (G) at  ($(S)!2/3!(I)$);
						\coordinate (N)  at (intersection of M--m and I--C);
						\coordinate (F) at (intersection of M--N  and B--C);
						\coordinate (O) at (intersection of M--I and  D--C);
						\draw (S)--(A) 
						 (S)--(B)
						 (S)--(C)
						(A)--(B) 
						(B)--(C) 
						 (S)--(I);
						\draw[dashed]  (A)--(D)
						 (C)--(D)
						 (S)--(D)
						 (D)--(O)--(M)--(F)
						(N)--(G)--(M)--(I)--(C); 
						\foreach \i/\g in {S/90,A/180,B/-90,C/-90,D/0,M/120,I/180,F/-90,N/-90,G/-60,O/0}{\draw[fill](\i) circle (1pt) 
						 ($(\i)+(\g:3mm)$) node[scale=1]{$\i$};}
					\end{tikzpicture}
			\end{center}
		\end{center}
		\begin{enumerate}
			\item Gọi  $F$ là giao điểm của $MN$ và $BC$.\\
			Ta có $MN\parallel AB$, suy ra $NF\parallel BI$ (vì  $F\in MN$, $I\in AB)$.\\
			Trong tam giác $CIB$  có $NF\parallel BI$, nên theo định lí Thalès ta có $\dfrac{NI}{CI}=\dfrac{BF}{BC}=\dfrac{1}{3}$.\\
			Trong tam giác $SAB$, ta có $G$ là trọng tâm nên  $\dfrac{GI}{SI}=\dfrac{1}{3}$.\\
			Trong tam giác  $SIC$, ta có $\dfrac{GI}{SI}=\dfrac{NI}{CI}=\dfrac{1}{3}$, suy ra $NG\parallel SC$ (định lí Thalès đảo).\\
			Do đó  $NG\parallel (SDC)$.\\
			 \item Trong mặt phẳng $(ABCD)$, gọi $O$ là giao điểm của $MI$ và $DC$ .\\
			Trong tam giác $OCI$ có $MN\parallel OC$ , suy ra $\dfrac{MI}{OI}=\dfrac{IN}{IC}=\dfrac{1}{3}$ (theo  định lí Thalès).\\
			Mà $\dfrac{IG}{SI}=\dfrac{1}{3}$($G$ là trọng tâm của tam  giác $SAB$).\\
			Do đó,  trong tam giác $SOI$ có $\dfrac{MI}{OI}=\dfrac{IG}{SI}=\dfrac{1}{3}$, suy ra $MG\parallel SO$(định lí Thalès đảo).\\
			Do đó $MG\parallel  (SDC)$.
		\end{enumerate}
	} 
\end{ex}

% nguồn: Ok_TuanLuan_Toan11_Lan2
\begin{ex}%[1H4H3-2]%[ma:9498a543]
	Cho tứ diện $ABCD$. Gọi  $M$, $N$ lần lượt là trung điểm của $AB$ và $AD$. Chứng minh rằng $MN\parallel (BCD)$.
	 \loigiai{
		\begin{center}
			\begin{tikzpicture}[line join=round,line cap=round]
					\coordinate (A)  at (0,4);
					 \coordinate (D) at (-2,0);
					\coordinate (B) at (0,-3); 
					\coordinate (C) at  (3,0);
					\coordinate (M) at  ($(A)!0.5!(B)$);
					\coordinate (N)  at ($(A)!0.5!(D)$);
					\draw (A)--(D) 
						 (A)--(B)
						(A)--(C) 
						 (B)--(C)
						 (D)--(B)
						 (M)--(N);
					\draw[dashed](D)--(C); 
					\foreach \i/\g  in {A/90,D/180,B/-90,C/0,M/0,N/180}{\draw[fill](\i) circle (1.0pt)
						 ($(\i)+(\g:3mm)$) node[scale=1]{$\i$};}
				\end{tikzpicture}
		\end{center}
		Vì $M$,  $N$ lần lượt là trung điểm của $AB$, $AD$ nên $MN$ là đường trung bình của tam giác $ABD$.\\
		Suy ra $MN\parallel  BD$.\\
		Mà $BD\subset (BCD)$ nên $MN\parallel  (BCD)$.
	} 
\end{ex}

% nguồn: Ok_TuanLuan_Toan11_Lan2
\begin{ex}%[1H4H3-2]%[ma:6896a474]
	Cho hình chóp $S.ABCD$ có đáy $ABCD$ là hình bình hành. Gọi $G_1$, $G_2$ lần lượt là trọng tâm  của tam giác $ABC$ và $SBC$. Chứng minh $G_1G_2\parallel ( SAB)$.
	\loigiai{ 
		Gọi $I$  là trung điểm $BC$.\\
		Xét tam giác $SIA$ Ta có $\dfrac{IG_1}{IA}=\dfrac{IG_2}{IS}=\dfrac{1}{3} \Rightarrow G_1G_2\parallel  SA$.
		\begin{center}
			\begin{tikzpicture}[scale=1,font=\footnotesize,line join=round,line cap=round]
					 \foreach \x\y\t in {0/0/A,-2.7/-2.0/B,2.5/-2.0/C,-1.8/3.8/S}
					\coordinate  (\t) at (\x,\y);
					\coordinate (D) at  ($(A)+(C)-(B)$);
					\coordinate (I) at ($(B)!1/2!(C)$); 
					\coordinate (G1)  at ($(I)!1/3!(A)$);
					\coordinate (G2) at ($(I)!1/3!(S)$); 
					\draw (I)--(S)--(B)--(C)--(S)--(D)--(C); 
					 \draw[dashed](B)--(A)--(D)
						(S)--(A)--(I) 
						 (G1)--(G2);
					\foreach \t/\g in {S/90,A/160,B/-140,C/-45,D/0,I/-90,G1/0,G2/180} 
					 \draw[fill] (\t) circle (1.0pt)
						node[shift={(\g:9pt)}]{$\t$}; 
				\end{tikzpicture}
		\end{center}
		Do đó $\heva{& G_1G_2\not \subset ( SAB) \\ & G_1G_2\parallel SA \\ & SA \subset ( SAB) } \Rightarrow G_1G_2\parallel ( SAB)$. 
	} 
\end{ex}

% nguồn: Ok_TuanLuan_Toan11_Lan2
\begin{ex}%[1H4H3-2]%[ma:3b2fb1f7]
	Cho hình chóp $S.ABC$ có $M$ là trung điểm của $AB$. Gọi $(P)$ là mặt phẳng chứa $CB$  và song song với $SA$, Gọi $(Q)$ là mặt phẳng chứa $CM$ và song song với $SA$.
	\begin{enumerate}
		\item Tìm giao tuyến của hai mặt phẳng $(P)$ và  $(Q)$.
		\item Vẽ  đường thẳng $b$ qua $B$ và $b\parallel SA$. Chứng minh $b\subset(P)$.
	\end{enumerate}
	\loigiai{ 
		\begin{center}
			\begin{tikzpicture}[scale=1,font=\footnotesize,line join=round,line cap=round,>=stealth]
					 \path
						(0,0) coordinate (A) 
						(5,0) coordinate  (C)
						(3,-1.5) coordinate (B) 
						(1,3.5)  coordinate (S)
						($(B)!.5!(A)$)  coordinate (M)
						($(B)+(S)-(A)$)  coordinate (b)
						($(b)+(C)-(B)$) coordinate (a); 
					\draw (S)--(A)--(B)--(C)--(S)--(B) 
						 (B)--(b)
						 (C)--(a);
					\draw [dashed] (M)--(C)--(A); 
					\foreach \d/\g  in {S/90,A/135,B/-90,C/45,M/-135}
					\fill  (\d) circle (1pt) + (\g:.3) node{$\d$};
					\foreach \d/\g in {a/-45,b/-45} 
					\fill  (\d) + (\g:.3) node{$\d$};
					\draw pic["$P$",draw=black,angle eccentricity=.5,angle radius=.5cm]{angle=a--C--B}; 
					\draw  pic["$Q$",draw=black,angle eccentricity=.75,angle radius=1cm]{angle=a--C--M};
				\end{tikzpicture}
		\end{center}
		\begin{enumerate}
			\item Ta có hai mặt phẳng ($P$) và ($Q$) cùng có điểm chung $C$ và  cùng song song với $SA$, suy ra giao tuyến của $(P)$ và $(Q)$ là đường thằng $a$ đi qua $C$ và $a\parallel SA$.
			\item Ta có $SA\parallel (P)$ và $B$ thuộc $(P)$, $b$ là đường thẳng đi  qua $B$ và $b\parallel SA$ suy ra $b\subset(P)$.
		\end{enumerate}
	 }
\end{ex}

% nguồn: Ok_TuanLuan_Toan11_Lan2
\begin{ex}%[1H4H3-2]%[ma:c2cc9e85]
	Cho hình chóp $S.ABCD$ có đáy $ABCD$ là hình bình hành. Trên  các cạnh $SB$ và $AD$ lần lượt lấy các điểm $N,P$ sao cho $\dfrac{SN}{SB}=\dfrac{PD}{AD}$. Chứng minh $NP\parallel ( SCD)$.
	 \loigiai{
		\begin{center}
			\begin{tikzpicture}[scale=1,font=\footnotesize,line join=round,line cap=round]
					\foreach \x\y\t in  {0/0/A,-1.7/-1.6/B,2.5/-1.6/C,0.4/2.8/S}
					\coordinate (\t) at  (\x,\y);
					\coordinate (D) at ($(A)+(C)-(B)$); 
					 \coordinate (Q) at ($(B)!1/3!(C)$);
					\coordinate  (P) at ($(D)!1/3!(A)$);
					\coordinate (N) at ($(S)!1/3!(B)$); 
					\coordinate  (R) at ($(S)!1/3!(C)$);
					\draw (S)--(B)--(C)--(S)--(D)--(C) 
						(N)--(R)--(D) 
						(R)--(Q); 
					 \draw[dashed](N)--(P)--(B)--(A)--(D)--(Q)
						 (S)--(A);
					\foreach  \t/\g in {S/90,A/160,B/-140,C/-45,D/0,Q/-90,P/50,N/160,R/60}
					\draw[fill] (\t)  circle (1.0pt)
						node[shift={(\g:9pt)}]{$\t$}; 
				\end{tikzpicture}
		\end{center}
		Trong mặt  phẳng $( SBC)$ kẻ $NR\parallel BC$ với $R \in SC$ và kẻ $RQ\parallel SB$ với $Q \in BC$ thì Ta có \\
		$\dfrac{SN}{SB}=\dfrac{SR}{SC}$ và  $\dfrac{SR}{SC}=\dfrac{BQ}{BC}$.\\
		Kết hợp với giả thiết $\dfrac{SN}{SB}=\dfrac{PD}{AD}$ ta được  $\dfrac{BQ}{BC}=\dfrac{PD}{AD}$ suy ra $BQ=PD$ do $BC=AD$.\\
		Lại có $NR=BQ$ suy ra $NR=PD$. Hơn nữa ta có $\heva{& NR\parallel  BQ \\ & PD\parallel BQ }$ suy ra $NR\parallel PD$.\\
		Do đó $PDRN$ là hình bình hành nên $NP\parallel DR$. Ta có $\heva{& NP\not \subset ( SCD) \\ & NP\parallel DR \\ & DR \subset ( SCD) } \Rightarrow NP\parallel (  SCD)$.
	} 
\end{ex}

% nguồn: Ok_TuanLuan_Toan11_Lan2
\begin{ex}%[1H4H3-2]%[ma:4e52a184]
	Cho hai hình bình hành $ABCD$ và $ABEF$ không cùng nằm trong $1$ mặt phẳng. Gọi $O$,  $O'$ lần lượt là tâm của $ABCD$ và $ABEF$. Chứng minh $OO'$ song song với các mặt phẳng $(ADF)$ và $(BCE)$.
	\loigiai{ 
		\begin{center}
			\begin{tikzpicture}[line join=round,line cap=round]
					 \coordinate (A) at (2,3);
					 \coordinate (B) at (6,3);
					 \coordinate (D) at (0,0);
					\coordinate (F) at (0,6); 
					\coordinate (C)  at ($(B)+(D)-(A)$);
					\coordinate (E)  at ($(B)+(F)-(A)$);
					 \coordinate (O') at ($(A)!0.5!(E)$);
					\coordinate (O)  at ($(A)!0.5!(C)$);
					\draw (E)--(F)--(D)--(C)--(B)--(E)--(C); 
					\draw[dashed] (E)--(A)--(B) 
						 (A)--(F)--(B)
						 (C)--(A)--(D)--(B)
						(O)--(O'); 
					 \foreach \i/\g in {A/180,B/0,C/-90,D/-90,E/90,F/90,O/-110,O'/100}{\draw[fill](\i) circle (1pt)
						($(\i)+(\g:3mm)$) node[scale=1]{$\i$};} 
				\end{tikzpicture}
		\end{center}
		Ta có  $BO=\dfrac{1}{2}BD$; $BO'=\dfrac{1}{2}BF$.\\
		Suy ra $OO'\parallel DF$.\\ 
		Mà $DF\subset  (ADF)\Rightarrow OO'\parallel (ADF)$.
	 }
\end{ex}

% nguồn: Ok_TuanLuan_Toan11_Lan2
\begin{ex}%[1H4H3-2]%[ma:c14d7600]
	Cho hình lăng trụ tam giác  $ABC.A'B'C'$. Gọi $M$ và $M'$ lần lượt là trung điểm của các cạnh $BC$ và $B'C'$. Chứng minh rằng
	\begin{enumEX}{2}
		 \item $AA'\parallel\left(BCC'B'\right)$.
		 \item $AM\parallel A'M$.
	\end{enumEX}
	 \loigiai{
		\begin{center}
			\begin{tikzpicture}[line join=round,line cap=round,line width=.6pt,font=\footnotesize,scale=1]
					 \coordinate[label=left:$A$] (A) at (0,0);
					 \coordinate[label=below left:$B$] (B) at (1,-1);
					\coordinate[label=right:$C$] (C)  at (4,0);
					\coordinate (H)  at ($(B)!.5!(C)$);
					\coordinate[label=left:$A'$]  (A1) at ($(H)+(90:4)$);
					 \coordinate[label=left:$B'$] (B1) at ($(B)-(A)+(A1)$);
					\coordinate[label=right:$C'$] (C1) at  ($(C)-(A)+(A1)$);
					\coordinate[label=below:$M$]  (M) at ($(B)!1/2!(C)$);
					\coordinate[label=above:$M'$] (M1) at  ($(B1)!1/2!(C1)$);
					 \draw (A1)--(A)--(B)--(C)--(C1)--(A1)--(B1)--(C1)
						(B)--(B1) 
						 (A1)--(M1)--(M);
					\draw[dashed] (A)--(C) 
						(A)--(M); 
					\fill  (A) circle (1pt)
						(B) circle  (1pt)
						(C) circle  (1pt)
						 (A1) circle (1pt)
						 (B1) circle (1pt)
						(C1) circle (1pt) 
						(M)  circle (1pt)
						(M1) circle (1pt); 
				\end{tikzpicture}
		\end{center}
		\begin{enumerate}
			\item Trong hình lăng trụ $ABC.A'B'C'$ ta có $AA'\parallel BB'$ và $BB'\subset\left(BCC'B'\right)\Rightarrow  AA'\parallel\left(BCC'B'\right)$.
			\item Vì $MM'\parallel BB'$, $MM'=BB'$ và $BB'\parallel AA'$,  $BB'=AA'$ nên $MM'\parallel AA'$, $MM'=AA'$.\\
			Suy ra $AMM'A'$ là  hình bình hành nên $AM\parallel A'M'$.
	\end{enumerate} }
\end{ex}

% nguồn: Ok_TuanLuan_Toan11_Lan2
\begin{ex}%[1H4H3-2]%[ma:e731d878]
	Cho hình thang $ABCD$ có $AB\parallel CD$ và $S \notin (ABCD)$. Trên $SA$, $BD$ lấy hai điểm $M$, $N$ sao cho $\dfrac{SM}{SA}=\dfrac{DN}{DB}=\dfrac{2}{3}$. Kẻ $NI\parallel AB$, ($I \in AD)$.  Chứng minh $MN\parallel (SCD)$.
	 \loigiai{
		\begin{center}
			\begin{tikzpicture}[scale=1,font=\footnotesize,line join=round,line cap=round,>=stealth]
					\path 
						(0:0) coordinate (A)--++(130:2) coordinate (D)--++(0:5)  coordinate (C)
						(A)--++(0:2) coordinate (B)--++(100:6) coordinate (S); 
					\draw ($(S)!.67!(A)$) coordinate  (M);
					\draw ($(D)!.67!(B)$) coordinate  (N);
					\coordinate[label =:$E$]  (E) at ($(A)+(N)-(B)$);
					 \path[name path=ne] (N)--(E);
					 \path[name path=ad] (A)--(D);
					\path[name intersections={of=ad 
						and  ne}]
						 (intersection-1) coordinate (I);
					\foreach \x/\g in {A/-90,B/-90,C/0,D/180,S/90,M/0,N/30,I/180} 
					 \fill(\x) circle (1pt)
						($(\x)+(\g:3mm)$) node{$\x$}; 
					\draw[dashed]  (N)--(I)
						 (D)--(C)
						 (M)--(N)
						(B)--(D); 
					\draw (I)--(M) 
						 (S)--(D)--(A)--(B)--(C)--(S)--(B)
						 (S)--(A);
				\end{tikzpicture}
		\end{center}
		Ta có  $\dfrac{AM}{AS}=\dfrac{1}{3}$. Do $NI\parallel AB$ nên $\dfrac{AI}{AD}=\dfrac{BN}{BD}=\dfrac{1}{3}$.\\
		Suy ra $\dfrac{AM}{AS}=\dfrac{AI}{AD}  \Rightarrow MI \parallel SD \Rightarrow M I\parallel (SCD)$.\\
		 Do $NI\parallel AB$ ta suy ra $NI\parallel CD\Rightarrow NI \parallel (SDC)$.\\
		Vậy $(MNI)\parallel (SCD) \Rightarrow  MN\parallel (SCD)$.
	 }
\end{ex}

% nguồn: Ok_TuanLuan_Toan11_Lan2
\begin{ex}%[1H4H3-2]%[ma:c60cd14c]
	Cho hình chóp $S.ABCD$ có đáy là hình bình hành. Gọi $G$ là trọng tâm của tam giác $SAD$, $M$ là điểm trên đoạn $DC$ sao  cho $DC=3DM$.\\
	Chứng minh rằng $MG\parallel  (SBC)$.
	\loigiai{ 
		\begin{center}
			\begin{tikzpicture}[line join=round,line cap=round,scale=0.7]
					 \coordinate (A) at (0,0);
					\coordinate  (B) at (2,-2);
					\coordinate  (D) at (5,0);
					\coordinate  (C) at ($(B)+(D)-(A)$);
					\coordinate (N) at ($(A)!0.5!(D)$); 
					\coordinate  (S) at ($(A)+(8,7)$);
					\coordinate (M) at  ($(D)!1/3!(C)$);
					\coordinate  (G) at ($(S)!2/3!(N)$);
					\coordinate (E) at (intersection  of M--N and B--C);
					\draw (S)--(A) 
						 (S)--(B)
						(S)--(C) 
						(A)--(B) 
						(B)--(C)--(E)--(S); 
					\draw[dashed](A)--(C) 
						 (A)--(D)
						(C)--(D) 
						(S)--(D) 
						(B)--(D) 
						(S)--(N)--(E) 
						(G)--(M); 
					\foreach \i/\g in  {S/90,A/180,B/-90,C/-90,D/110,M/50,N/90,E/0,G/180}{\draw[fill](\i) circle (1pt)
						($(\i)+(\g:3mm)$) node[scale=0.7]{$\i$};} 
				\end{tikzpicture}
		\end{center}
		Gọi $N$ là trung điểm của $AD$.\\ 
		Ta có  $MG\subset(SMN)$.\\
		Trong mặt phẳng  $(ABCD)$, gọi $E=MN\cap BC$.\\
		Ta có $S\in(SNM)\cap(SBC)$; $E\in MN$ và $MN\subset(SMN)$; $E\in BC$ và $BC\subset(SBC)$.\\ 
		 Suy ra $(SMN)\cap(SBC)=SE$.\\
		Dễ thấy  $\triangle MND$ và $\triangle MEC$ đồng dạng, suy ra $\dfrac{MN}{ME}=\dfrac{MD}{MC}=\dfrac{1}{2}$ , suy ra $\dfrac{MN}{NE}=\dfrac{1}{3}$.\quad (1)\\
		 Mặt khác, $\dfrac{GN}{SN}=\dfrac{1}{3}$ ($G$ là trọng tâm của tam giác $SAD$).\quad (2)\\
		Từ (1) và (2)  suy ra $\dfrac{GN}{SN}=\dfrac{MN}{NE}$.\\
		Theo định lí Thalès đảo trong tam giác $SNE$, ta có $MG\parallel SE$ .\\ 
		Mà $SE\subset(SBC)$ nên $MG\parallel (SBC)$. 
	} 
\end{ex}

% nguồn: Ok_TuanLuan_Toan11_Lan2
\begin{ex}%[1H4H3-2]%[ma:b494f44f]
	Cho hai hình bình hành $ABCD$ và $ABEF$ không cùng nằm trong một mặt phẳng. Gọi $M$, $N$ lần lượt là hai điểm trên các cạnh $AE,BD$ sao cho $AM=\dfrac{1}{3}AE,BN=\dfrac{1}{3}BD$. Chứng minh $MN$  song song với $(CDEF)$.
	\loigiai{ 
		\begin{center}
			\begin{tikzpicture}[line join=round,line cap=round]
					\coordinate (A) at (2,2.5); 
					 \coordinate (B) at (7.5,2.5);
					\coordinate  (D) at (0,0);
					\coordinate  (F) at (0,5);
					\coordinate (C) at  ($(B)+(D)-(A)$);
					\coordinate (E) at  ($(B)+(F)-(A)$);
					 \coordinate (O') at ($(A)!0.5!(E)$);
					\coordinate (O) at ($(A)!0.5!(C)$); 
					\coordinate  (M) at ($(A)!1/3!(E)$);
					\coordinate  (N) at ($(B)!1/3!(D)$);
					\coordinate (I) at (intersection of A--N and  C--D);
					\coordinate (i) at (intersection of  A--I and C--B);
					\draw  (E)--(F)--(D)--(C)--(B)--(E)--(C)
						 (i)--(I)--(C);
					\draw[dashed]  (E)--(A)--(B)
						(A)--(F)--(B) 
						 (C)--(A)--(D)--(B)
						(O)--(O') 
						 (M)--(N)
						 (A)--(i);
					\foreach \i/\g  in {A/180,B/0,C/-90,D/-90,E/90,F/90,M/160,N/-80,O/-110,O'/100,I/0}{\draw[fill](\i) circle (1pt)
						 ($(\i)+(\g:3mm)$) node[scale=1]{$\i$};}
				\end{tikzpicture}
		\end{center}
		Trong $(ABCD)$, gọi $I=AN\cap CD$.\\ 
		 Do $AB\parallel CD$ nên $\dfrac{AN}{AI}=\dfrac{BN}{BD}\Rightarrow\dfrac{AN}{AI}=\dfrac{1}{3}$.\\
		Lại có $\dfrac{AM}{AE}=\dfrac{1}{3}\Rightarrow\dfrac{AN}{AI}=\dfrac{AM}{AE}\Rightarrow MN\parallel  IE$.\\
		Mà $I\in CD\Rightarrow IE\subset(CDEF)\Rightarrow  MN\parallel (CDEF)$.
	 }
\end{ex}

% nguồn: Ok_TuanLuan_Toan11_Lan2
\begin{ex}%[1H4H3-2]%[ma:cbfc8074]
	Cho tứ diện $ABCD$. $G$ là trọng tâm của $\triangle ABD$. $M$ là điểm trên cạnh  $BC$ sao cho $MB=2MC$. Chứng minh $MG\parallel (ACD)$.
	 \loigiai{
		\begin{center}
			\begin{tikzpicture}[line join=round,line cap=round]
					 \coordinate (A) at (2,4);
					\coordinate (D) at  (-2,0);
					\coordinate  (C) at (0,-3);
					 \coordinate (B) at (3,0);
					\coordinate (M) at  ($(C)!1/3!(B)$);
					\coordinate (E) at ($(B)!0.5!(C)$); 
					\coordinate (G)  at ($(D)!2/3!(E)$);
					 \draw (A)--(D)
						 (A)--(C)
						 (A)--(B)
						 (B)--(C)
						 (D)--(C);
					\draw[dashed](D)--(B) 
						 (E)--(D)
						 (M)--(G);
					\foreach \i/\g  in {A/90,D/180,B/0,C/-90,M/-40,G/90,E/0}{\draw[fill](\i) circle (1.0pt)
						($(\i)+(\g:3mm)$) node[scale=1]{$\i$};} 
				\end{tikzpicture}
		\end{center}
		Gọi $E$ là  trung điểm cạnh $BC$.\\
		Do $G$ là trọng tâm tam giác $BCD$, nên ta có  $GD=\dfrac{2}{3}ED$.\\
		 Mặt khác $3MC=BC\Rightarrow 3MC=2EC\Rightarrow\dfrac{MC}{EC}=\dfrac{2}{3}$.\\
	 Suy ra $MG\parallel CD$, mà $CD\subset (ACD)$ nên $MG\parallel (ACD)$.}
\end{ex}

% nguồn: Ok_TuanLuan_Toan11_Lan2
\begin{ex}%[1H4H3-2]%[ma:48485f35]
	Cho hai hình bình hành $ABCD$ và $ABMN$ không cùng nằm trong một phẳng. Gọi $M$, $N$ lần lượt là trọng tâm của hai tam giác $ABF$ và $ABC$. Chứng minh rằng đường thẳng $MN$ song song  với mặt phẳng $( ACF)$.
	\loigiai{ 
		\begin{center}
			\begin{tikzpicture}[scale=1,line cap=round,line join=round,font=\footnotesize,>=stealth]
					 \def\dai{4} \def\rong{1.6} \def\goc{-110}
					\path (0,0)  coordinate (A)++(0:\dai) coordinate (B)
						++(\goc:\rong) coordinate  (C)++(-180:\dai) coordinate (D)
						($(A)+(110:\rong)$) coordinate (F)++(0:\dai)  coordinate (E);
					\coordinate (H) at  ($(A)!1/2!(F)$);
					\coordinate  (K) at ($(A)!1/2!(C)$);
					\coordinate (M) at  ($(B)!2/3!(H)$);
					\coordinate (N)  at ($(B)!2/3!(K)$);
					\draw  (A)--(D)--(C)--(B)--(E)--(F)--(A)--(C)--(F)--(B)
						 (H)--(K);
					\draw[dashed]  (H)--(B)--(K)
						(M)--(N) 
						(A)--(B); 
					\foreach \t/\g in {A/180,B/0,C/0,D/180,E/0,F/180,K/-110,H/-170,M/90,N/-30} 
					\draw[fill] (\t)  circle (1.0pt)
						 node[shift={(\g:9pt)}]{$\t$};
				\end{tikzpicture}
		\end{center}
		Gọi $H$, $K$ lần lượt là trung điểm của $AF$,  $AC$.\\
		Ta có  $\dfrac{BM}{BH}=\dfrac{BN}{BK}=\dfrac{2}{3} \Rightarrow MN\parallel HK$.\\
		Mà $HK \subset (  ACF)$ và $MN\not \subset ( ACF)$ nên suy ra $MN\parallel ( ACF)$.
	 }
\end{ex}
